\makeatletter
\ifdefined\XeTeXversion
  \providecommand{\headerps@out}[1]{}
\else\ifdefined\directlua
  \providecommand{\headerps@out}[1]{}
\fi\fi
\makeatother
\documentclass[review,onefignum,onetabnum,nohypdvips]{siamart220329}
% 共享文档信息 (可以从主文档和补充文档中调用)

% Encoding and typography (load early)
\usepackage[utf8]{inputenc} % UTF-8 input encoding
\usepackage[T1]{fontenc}    % 8-bit T1 fonts
\usepackage{microtype}      % Microtypography improvements

% General utilities
\usepackage{lipsum}         % Placeholder text
\usepackage{graphicx}       % Image insertion
% \usepackage{epstopdf}     % Uncomment only if EPS conversion is explicitly needed
\usepackage{xcolor}         % Colors
\usepackage{colortbl}       % Table-row shading
\usepackage{url}            % URL typesetting
\usepackage{nicefrac}       % Compact fraction symbols
\usepackage{xspace}         % Smart spacing for macros
\usepackage{geometry}       % Page layout

% Mathematical packages
\usepackage{amsmath}        % Core math support
% \theoremstyle{remark}
\usepackage{amssymb}        % Math symbols (includes amsfonts)
\usepackage{bbm}            % Blackboard bold 1 for \mathbbm
\usepackage{mathrsfs}       % Script fonts
\usepackage{bm}             % Bold math symbols
\usepackage{booktabs}       % Professional tables
\usepackage{threeparttable} % Table notes
\usepackage{makecell}       % Multi-line table cells
\usepackage{subcaption}     % Subfigures
\usepackage{multirow}       % Multirow tables
\usepackage{array}          % Enhanced table formatting
\usepackage{tabularx}
\usepackage{setspace}       % Line spacing control

% Algorithms
\makeatletter
\@ifpackageloaded{algorithm}{}{\usepackage{algorithm}}
\makeatother
\usepackage{algorithmic}    % Consider algorithmicx or algorithm2e for modern alternatives
\usepackage{tikz}
\usetikzlibrary{arrows.meta,positioning,fit,shapes.misc}
\makeatletter
\@ifpackageloaded{hyperref}{}{\usepackage{hyperref}}
\makeatother
\usepackage{nameref}        % Reference names
\makeatletter
\@ifpackageloaded{cleveref}{}{\usepackage{cleveref}}
\makeatother

\geometry{left=3.0cm,right=3.0cm}

% Ensure hyperref settings are robust
\hypersetup{
    colorlinks=true,
    linkcolor=blue,
    citecolor=blue,
    urlcolor=blue
}
%-----------------------------数学符号新命令------------------------------
\newcommand{\eucalB}{\EuScript{B}}
\newcommand{\eucalP}{\EuScript{P}}
\newcommand{\Var}{\mathrm{Var}}
\newcommand{\Cov}{\mathrm{Cov}}
\newcommand{\Ex}{\mathsf{E}}
\newcommand{\Prob}{\mathsf{P}}
\newcommand{\one}{\mathbbm{1}}
\newcommand{\One}{\bm{1}}
\newcommand{\zero}{\bm{0}}
\newcommand{\RR}{\mathbb{R}}
\newcommand{\R}{\mathbb{R}}
\newcommand{\ZZ}{\mathbb{Z}}
\newcommand{\EE}{\mathbb{E}}
\newcommand{\PP}{\mathbb{P}}
\newcommand{\Rp}{\RR_+}
\newcommand{\TT}{\mathbb{T}}
\newcommand{\NN}{\mathbb{N}}
\newcommand{\VV}{\mathbb{V}}
\newcommand{\diff}{\mathrm{d}}
\newcommand{\prox}{\mathrm{prox}}
\newcommand{\rprox}{\mathrm{rprox}}
\newcommand{\ran}{\mathrm{ran}}
\newcommand{\Id}{\mathrm{Id}}
\newcommand{\e}{\mathrm{e}}
\newcommand{\calA}{\mathcal{A}}
\newcommand{\A}{\mathcal{A}}
\newcommand{\dist}{\operatorname{dist}}
\newcommand{\rKKT}{r_{\mathrm{KKT}}}
\newcommand{\bB}{\bm{B}}
\newcommand{\bH}{\bm{H}}
\newcommand{\bF}{\bm{F}}
\newcommand{\bS}{\bm{S}}
\newcommand{\bX}{\bm{X}}
\newcommand{\bx}{\bm{x}}
\newcommand{\bb}{\bm{b}}
\newcommand{\br}{\bm{r}}
\newcommand{\bbeta}{\bm{\beta}}
\newcommand{\bt}{\bm{t}}
% \newcommand{\eproof}{\hfill $\square$}
\newcommand{\bh}{\bm{h}}
\newcommand{\bu}{\bm{u}}
\newcommand{\bz}{\bm{z}}
\newcommand{\ba}{\bm{a}}
\newcommand{\by}{\bm{y}}
\newcommand{\be}{\bm{e}}
\newcommand{\bc}{\bm{c}}
\newcommand{\bU}{\bm{U}}
\newcommand{\bZ}{\bm{Z}}
\newcommand{\bY}{\bm{Y}}
\newcommand{\blmbd}{\bm{\lambda}}
\newcommand{\bLmbd}{\bm{\Lambda}}
\newcommand{\bd}{\bm{d}}
\newcommand{\bq}{\bm{q}}
\newcommand{\bmu}{\bm{\mu}}
\DeclareMathOperator*{\argmin}{\mathrm{argmin}}
\DeclareMathOperator*{\Argmin}{\mathrm{Argmin}}
\DeclareMathOperator*{\minimize}{\mathrm{minimize}\,}
\DeclareMathOperator*{\maximize}{\mathrm{maximize}}
\DeclareMathOperator*{\argmax}{\mathrm{argmax}}
\DeclareMathOperator*{\Diag}{\mathrm{Diag}}
\newcommand{\calX}{\mathcal{X}}
\newcommand{\calY}{\mathcal{Y}}
\newcommand{\calB}{\mathcal{B}}
\newcommand{\calL}{\mathcal{L}}
\newcommand{\calG}{\mathcal{G}}
\newcommand{\calS}{\mathcal{S}}
\newcommand{\calC}{\mathcal{C}}
\newcommand{\calJ}{\mathcal{J}}
\newcommand{\calD}{\mathcal{D}}
\newcommand{\calQ}{\mathcal{Q}}
\newcommand{\calO}{\mathcal{O}}
\newcommand{\ent}{\mathrm{ent}}
% \newcommand{\diag}{\mathrm{diag}}
\newcommand{\dom}{\mathrm{dom}}
\newcommand{\cl}{\mathrm{cl}}
\DeclareMathOperator{\conv}{conv}
\newcommand{\interior}{\mathrm{int}}
\newcommand{\zer}{\mathrm{zer}}
\newcommand{\epi}{\mathrm{epi}}
\newcommand{\amax}{a_{\max}}
\newcommand{\amin}{a_{\min}}
\newcommand{\bmax}{b_{\max}}
\newcommand{\bmin}{b_{\min}}
\newcommand{\bp}{\bm{p}}
\newcommand{\bD}{\bm{D}}
\newcommand{\mW}{\mathrm{W}}
\newcommand{\bv}{\bm{v}}
\newcommand{\bw}{\bm{w}}
\newcommand{\ta}{\tilde{a}}
\newcommand{\tb}{\tilde{b}}
\newcommand{\tp}{\widetilde{\bp}}
\newcommand{\sumn}{\sum_{i=1}^n}
\newcommand{\sumQ}{\sum_{q=1}^Q}
\newcommand{\sumq}{\sum_{k=1}^q}
\newcommand{\summ}{\sum_{i=1}^m}
\newcommand{\sumd}{\sum_{i=1}^d}
\newcommand{\sumT}{\sum\limits_{ k =1}^{T}}
\newcommand{\sumTa}{\frac{1}{T} \sum\limits_{ k =1}^{T}}
\newcommand{\bC}{\bm{C}}
\newcommand{\bI}{\bm{I}}
\newcommand{\bM}{\bm{M}}
\newcommand{\bR}{\bm{R}}
\newcommand{\bV}{\bm{V}}
\newcommand{\bW}{\bm{W}}
\newcommand{\bxi}{\bm{\xi}}
\newcommand{\tmu}{\widetilde{\bmu}}
\newcommand{\tU}{\widetilde{\bU}}
\newcommand{\tV}{\widetilde{\bV}}
\newcommand{\tM}{\widetilde{\bM}}
\newcommand{\bDelta}{\bm{\Delta}}
\newcommand{\bs}{\bm{s}}
\newcommand{\ox}{\overline{x}}
\newcommand{\ob}{\overline{b}}
\newcommand{\op}{\overline{p}}
\newcommand{\obp}{\overline{\bp}}
\newcommand{\obmu}{\overline{\bmu}}
\newcommand{\obC}{\overline{\bC}}
\newcommand{\obD}{\overline{\bD}}
\newcommand{\bsigma}{\bm{\sigma}}
\newcommand{\bomega}{\bm{\omega}}
\newcommand{\tomega}{\widetilde{\omega}}
\newcommand{\tr}{\mathrm{tr}}
\newcommand{\an}{\mathrm{an}}
\newcommand{\bj}{\bm{j}}
\newcommand{\bG}{\bm{G}}
\newcommand{\bcalG}{\bm{\calG}}
\newcommand{\bcalK}{\bm{\calK}}
\newcommand{\bcalX}{\bm{\calX}}
\newcommand{\fX}{\mathfrak{X}}
\newcommand{\bfX}{\bm{\fX}}
\newcommand{\bell}{\bm{\ell}}
\newcommand{\bL}{\bm{L}}
\newcommand{\bzeta}{\bm{\zeta}}
\renewcommand{\vec}{\mathrm{vec}}
\newcommand{\hy}{\widehat{y}}
\newcommand{\ul}[1]{\underline{#1}}
\makeatletter
% Each environment has its own counter, reset at the start of each section.
\@ifundefined{theorem}{\newtheorem{theorem}{Theorem}[section]}{}
\@ifundefined{newsiamremark}{
  \@ifundefined{lemma}{\newtheorem{lemma}{Lemma}[section]}{}
  \@ifundefined{corollary}{\newtheorem{corollary}{Corollary}[section]}{}
  \@ifundefined{proposition}{\newtheorem{proposition}{Proposition}[section]}{}
  \@ifundefined{definition}{\newtheorem{definition}{Definition}[section]}{}
  \@ifundefined{assumption}{\newtheorem{assumption}{Assumption}[section]}{}
  \@ifundefined{remark}{\newtheorem{remark}{Remark}[section]}{}
  \@ifundefined{example}{\newtheorem{example}{Example}[section]}{}
}{
  % Override the shared counters supplied by the SIAM class.
  \theoremstyle{plain}
  \theoremheaderfont{\normalfont\scshape}
  \theorembodyfont{\normalfont\itshape}
  \theoremseparator{.}
  \theoremsymbol{}
  \renewtheorem{lemma}{Lemma}[section]
  \renewtheorem{corollary}{Corollary}[section]
  \renewtheorem{proposition}{Proposition}[section]
  \renewtheorem{definition}{Definition}[section]
  \theoremheaderfont{\normalfont\itshape}
  \theorembodyfont{\normalfont}
  \@ifundefined{assumption}{
    \newtheorem{assumption}{Assumption}[section]
  }{}
  \@ifundefined{remark}{\newtheorem{remark}{Remark}[section]}{}
  \@ifundefined{example}{\newtheorem{example}{Example}[section]}{}
}
\@ifundefined{proof}{\newenvironment{proof}[1][Proof]{\par\noindent\textbf{#1.}\space}{\hfill$\square$\par}}{}
\@ifundefined{noproof}{\newenvironment{noproof}{\par\noindent\textbf{Proof.}\space}{\hfill\par}}{}
\makeatother

\newcommand{\bcL}{\overline{\calL}}
\newcommand{\KL}{K\L\xspace}
\newcommand{\res}{\mathrm{res}}
\newcommand{\eproof}{\hfill $\square$}
\newcommand{\calF}{\mathcal{F}}

\newcommand{\commyy}[1]{{\color{red}(yy: #1)}} % Yuan Yao's comments

\numberwithin{equation}{section}
\newcommand{\blue}[1]{{\color{blue}#1}}
\newcommand{\red}[1]{{\color{red}#1}}
\newcommand{\todo}[1]{{\color{red}\textbf{TODO:} #1}}

%% Packages and macros go here
\usepackage{lipsum}
\usepackage{amsfonts}
\ifpdf
  \DeclareGraphicsExtensions{.eps,.pdf,.png,.jpg}
\else
  \DeclareGraphicsExtensions{.eps}
\fi

% Prevent itemized lists from running into the left margin inside theorems and proofs
\usepackage{enumitem}
\setlist[enumerate]{leftmargin=.5in}
\setlist[itemize]{leftmargin=.5in}

% Add a serial/Oxford comma by default.
\newcommand{\creflastconjunction}{, and~}

% \headers{Shrink the Dual Regularization, Keep the Penalty}{L. Li, K. Deng, C. Cui, and J. Hu}
\headers{Shrink the Dual Regularization, Keep the Penalty}{}
\title{{Shrink the Dual Regularization, Keep the Penalty: Tight Complexity of Linearized Augmented Lagrangian Methods}}
% \author{
% Luxuan Li\thanks{School of Mathematical Sciences, Beihang University,
% Beijing 100191, China}
% \and Kangkang Deng\thanks{Department of Mathematics, National University of Defense Technology,
% Changsha 410073, China.
% Corresponding author. Email: \texttt{freedeng1208@gmail.com}.
% }
% \and Chunfeng Cui\thanks{School of Mathematical Sciences and State Key Laboratory of Virtual Reality Technology and Systems,
% Beihang University, Beijing 100191, China.}
% \and Jiang Hu\thanks{Yau Mathematical Sciences Center, Tsinghua University,
% Beijing 100084, China}
% }

\begin{document}

\maketitle

\begin{abstract}
A classical advantage of augmented Lagrangian methods over penalty methods is
that the penalty parameter can stay bounded. For single-loop linearized
augmented Lagrangian methods applied to the nonconvex composite problem
$\min_{\bx} f(\bx)+h(\calA(\bx))$, where $f$ and $\calA$ are smooth and $h$ is
convex and Lipschitz, this advantage has not been reflected in complexity
guarantees: in existing single-loop methods that find an $\epsilon$-KKT point
in $\mathcal O(\epsilon^{-3})$ iterations, the penalty or the dual
regularization grows with $1/\epsilon$. We ask whether such growth is
necessary, and whether it must be carried by the penalty. We study a family of
methods that perform one gradient descent--ascent step per iteration on the
reduced augmented Lagrangian with a fixed penalty $\rho$, a dual
regularization with parameter $s$, and moving primal and dual centers. Every
member of this family has the combined penalty $\rho+s$. By constructing hard
instances, we show that if $\epsilon(\rho+s)$ is at most a constant, then
every choice of constant center stepsizes requires $\Omega(\epsilon^{-4})$
iterations, and for fixed $\rho$ this order is attained within the family. In contrast, with
$\rho$ fixed and $s$ of order $1/\epsilon$, a member of the family finds an
$\epsilon$-KKT point in $\mathcal O(\epsilon^{-3})$ iterations. Hence the
combined penalty must reach the order $1/\epsilon$ to improve on
$\epsilon^{-4}$, and this growth can be carried by the dual regularization
alone while the penalty stays fixed.
\end{abstract}

\begin{keywords}
augmented Lagrangian method, nonconvex composite optimization, iteration
complexity, lower bound, dual regularization, gradient descent--ascent
\end{keywords}

\begin{AMS}
65K05, 90C26, 90C30, 90C46, 90C60
\end{AMS}

\section{Introduction}
\label{sec:introduction}
In this paper, we study the structured nonsmooth nonconvex composite problem
\begin{align}
    \min_{\bx\in\RR^n}\quad & F(\bx) := f(\bx) + h(\calA(\bx)),
    \tag{\textbf{P}}
    \label{prob:main}
\end{align}
where $f\colon\RR^n\to\RR$ is a possibly nonconvex function with a Lipschitz
continuous gradient, $\calA\colon\RR^n\to\RR^m$ is a smooth and possibly nonlinear
mapping, and $h\colon\RR^m\to\RR\cup\{+\infty\}$ is a proper closed convex function
whose proximal mapping is easy to evaluate; the complete assumptions are stated in
Section~\ref{sec:assumptions}. The two components play complementary roles: $f$ is
usually a smooth loss or data-fidelity term, whereas $h$ is a nonsmooth convex function
that imposes robustness, sparsity, or other structure on the transformed variable
$\calA(\bx)$
\cite{charisopoulos2021low,markovsky2008structured,parekh2015convex}. Several models
of practical interest take this form. 

\subsection{Motivating Examples}

\textbf{Robust phase retrieval}~\cite{duchi2019solving}.
Given measurement vectors $\mathbf a_i\in\RR^n$ and observations
$b_i\approx\langle\mathbf a_i,\bx^\star\rangle^2$, a fraction of which may be grossly
corrupted, one solves
\begin{equation*}
  \min_{\bx\in\RR^n}\;\frac1m\sum_{i=1}^m
  \bigl|\langle\mathbf a_i,\bx\rangle^2-b_i\bigr|.
\end{equation*}
This is~Problem \eqref{prob:main} with $f\equiv0$, $h=\frac1m\|\cdot\|_1$, and
$\calA(\bx)_i=\langle\mathbf a_i,\bx\rangle^2-b_i$; the $\ell_1$ loss makes the
model robust to outliers.

\textbf{Robust low-rank matrix recovery}~\cite{charisopoulos2021low}.
Given a linear measurement map $\mathcal M\colon\RR^{d_1\times d_2}\to\RR^m$ and
observations $\mathbf b$ containing outliers, a rank-$r$ matrix is recovered in
factorized form via
\begin{equation*}
  \min_{\mathbf U\in\RR^{d_1\times r},\,\mathbf V\in\RR^{d_2\times r}}
  \frac1m\bigl\|\mathcal M(\mathbf U\mathbf V^\top)-\mathbf b\bigr\|_1.
\end{equation*}
Here $\bx=(\mathbf U,\mathbf V)$, $h=\frac1m\|\cdot\|_1$, and
$\calA(\bx)=\mathcal M(\mathbf U\mathbf V^\top)-\mathbf b$ is bilinear in
$(\mathbf U,\mathbf V)$, so the model again has the form~Problem \eqref{prob:main}.

\textbf{Minimizing a pointwise maximum}~\cite{zhang2020single}.
For smooth losses $f_1,\dots,f_m$, such as the losses on $m$ data subgroups, one solves
\begin{equation*}
  \min_{\bx\in\RR^n}\;\max_{1\le i\le m} f_i(\bx).
\end{equation*}
This is~Problem \eqref{prob:main} with $h(u)=\max_{1\le i\le m}u_i$ and
$\calA(\bx)=(f_1(\bx),\dots,f_m(\bx))$; a smooth regularizer, if present, enters
through $f$, and the proximal mapping of $h$ reduces to a projection onto the unit
simplex.

\subsection{Related Work}

Introducing the auxiliary variable $\by=\calA(\bx)$ rewrites~Problem \eqref{prob:main} as
the equality-constrained problem
\begin{equation}\label{prob:split}
  \min_{\bx,\by}\; f(\bx)+h(\by)\quad\text{s.t.}\quad \calA(\bx)=\by,
\end{equation}
which brings it within the scope of penalty methods and augmented Lagrangian methods
(ALMs). For general nonlinearly constrained nonconvex problems, both families control
feasibility through the penalty parameter. Quadratic penalty methods increase the
penalty without bound~\cite{kong2019complexity,lin2022complexity}. ALMs were
originally introduced to avoid such unbounded penalties
\cite{hestenes1969multiplier,powell1969method}, and in convex programming any fixed
penalty suffices~\cite{rockafellar1973multiplier}; in the nonconvex setting, however,
inexact ALMs still increase the penalty across outer iterations to obtain complexity
guarantees~\cite{sahin2019inexact,li2021rate}, and proximal ALM, which keeps the
penalty fixed within a run, chooses its value according to the target accuracy
\cite{xie2021proximal}. Single-loop methods have attracted growing attention because
they avoid solving subproblems to high accuracy, yet they rely on the same mechanism:
single-loop quadratic penalty schemes increase the penalty~\cite{alacaoglu2024complexity},
and linearized ALMs use small dual steps oraccuracy $eps$-dependent parameters
\cite{shi2025momentum}. A large penalty inflates the smoothness constant of the
primal update, which forces small step sizes and makes the primal subproblems
increasingly ill-conditioned.

The split formulation~\eqref{prob:split} also carries favorable structure. Its
constraint Jacobian $[J_{\calA}(\bx),\,-I]$ has full row rank at every $\bx$, so
constraint qualifications hold automatically, and its objective separates into a
smooth part in $\bx$ and a prox-friendly part in $\by$. Nonconvex ADMM and related
Lagrangian-based splitting methods exploit exactly this structure
\cite{bot2020proximal,cohen2022dynamic,hallak2023adaptive,elbourkhissi2025perturbed}.
The structure, however, does not remove the need for a large penalty. For nonlinear
$\calA$, adaptive Lagrangian schemes require penalty levels proportional to
$1/\epsilon$~\cite{hallak2023adaptive}, and linearized perturbed ALMs use
accuracy-dependent penalty and perturbation parameters~\cite{elbourkhissi2025perturbed}.
The nonconvex--linear minimax route attains $\mathcal O(\epsilon^{-3})$
\cite{pan2021efficient,he2024approximation,xu2026rada}, but its dual regularization
acts as a penalty parameter that grows as the target accuracy decreases. Fixed-penalty
guarantees are available for smooth nonlinear equality constraints under a uniform LICQ
\cite{elbourkhissi2025linearized,liu2026fixed}, but they do not cover the nonsmooth
term $h(\by)$ in~\eqref{prob:split}. It therefore remains open whether a single-loop
ALM can keep every penalty-type parameter independent of both the accuracy
and the target accuracy while retaining an explicit first-order update for
$h\circ\calA$.

From another perspective, Problem~\eqref{prob:main} also admits a minimax reformulation. In particular, eliminating $\by$ from the Lagrangian of~\eqref{prob:split} gives
\begin{equation}\label{eq:ncl}
  \min_{\bx\in\RR^n}\max_{\blmbd\in\RR^m}\;
  \calL(\bx,\blmbd) =  \Phi(\bx,\blmbd)-h^*(\blmbd), \qquad \Phi(\bx,\blmbd):=f(\bx)+\langle\blmbd,\calA(\bx)\rangle,
\end{equation}
where $\operatorname{dom}h^*\subseteq\mathbb B(\mathbf 0,L_h)$ by
Assumption~\ref{ass:h}. This is a nonconvex--linear instance of the smooth
nonconvex--concave (NC-C) minimax problem. Since $\blmbd$ is exactly the multiplier of $\calA(\bx)=\by$, KKT points
of~\eqref{prob:split} and game-stationary (GS) points of~\eqref{eq:ncl} agree up to
constants, and ALMs can be read as single-loop minimax methods. 
For large-scale problems, we focus on single-loop first-order methods, which
update each variable by one explicit step per iteration without an inner
solver. The basic method of this kind for~\eqref{eq:ncl} is gradient
descent--ascent (GDA). It converges slowly because regularity fails on both
sides: $\Phi$ is only linear in $\blmbd$, so the best response
$\blmbd^\star(\bx)$ is in general set-valued, and $\Phi$ is nonconvex in $\bx$,
so $\bx^\star(\blmbd)$ need not be unique or stable. Even with two timescales,
GDA needs $\mathcal O(\epsilon^{-6})$ iterations to reach an $\epsilon$-GS
point~\cite{lin2020gradient}.

Two techniques improve on this rate, each restoring regularity on one side.
The first is primal smoothing, which adds $\frac{1}{2\gamma}\|\bx-\bz\|^2$ with
a fixed $\gamma$ and moves the center slowly by
$\bz^{t+1}=\bz^t+\beta(\bx^{t+1}-\bz^t)$. The target accuracy enters only
through $\beta$: with $\beta=\Theta(\epsilon^2)$, the resulting Smoothed GDA
finds an $\epsilon$-GS point in $\mathcal O(\epsilon^{-4})$ iterations. For the
pointwise maximum of finitely many smooth functions, i.e.,~\eqref{eq:ncl} with
$h(\bu)=\max_i u_i$, strict complementarity permits a constant $\beta$ and
improves the bound to $\mathcal O(\epsilon^{-2})$~\cite{zhang2020single}.

The second technique is dual regularization, which subtracts
$\frac{1}{2s}\|\blmbd-\bmu\|^2$ from the dual objective and makes it strongly
concave. In our reduced augmented-Lagrangian formulation, this regularization
has the following combined penalty interpretation:
\begin{equation}\label{eq:intro-effective-penalty}
  \max_{\blmbd}\mathcal P_s(\bx,\blmbd;\bz,\bmu)
  =\calD_{\rho+s}(\bx,\bmu)
    +\frac{1}{2\gamma}\|\bx-\bz\|^2,
\end{equation}
where $\mathcal P_s$ is defined in~\eqref{eq:potential} and
$\calD_{\rho+s}$ is the reduced augmented Lagrangian~\eqref{eq:reduced-al}
with the combined penalty parameter $\rho+s$; see Lemma~\ref{lem:effective-penalty}.
A zero center $\bmu=\mathbf 0$ recovers a quadratic penalty with
coefficient $\rho+s$, which can leave a nonzero feasibility residual
when $s$ is fixed. Updating $\bmu$ instead provides a multiplier-update
mechanism that may recover feasibility while keeping $\rho+s$ fixed;
its convergence and complexity depend on the coupled updates.

With the dual center fixed at $\bmu=\mathbf 0$, Perturbed GDA
attains an $\mathcal O(\epsilon^{-4})$ complexity for finding an
$\epsilon$-GS point~\cite{lu2020hybrid,xu2023unified}, and Li et al. establish
a matching algorithm-specific lower bound~\cite{li2026smoothing}.
Methods exploiting the NC-L structure attain
$\mathcal O(\epsilon^{-3})$~\cite{pan2021efficient,he2024approximation,xu2026rada}.
These guarantees rely on accuracy-dependent regularization parameters,
corresponding to a penalty parameter that increases as $\epsilon$ decreases. To improve the iteration complexity, Li et al. combine primal smoothing
with zero-centered dual perturbation to develop Perturbed Smoothed GDA,
which yields a genuine improvement in complexity order, reducing
the GS complexity of Perturbed GDA from $\mathcal O(\epsilon^{-4})$
to $\mathcal O(\epsilon^{-3})$ with a matching lower
bound~\cite{li2026smoothing}. However, the dual regularization coefficient still satisfies
$1/s=\Theta(\epsilon)$.

With a moving dual center $\bmu$ and fixed primal and dual regularization
coefficients, DS-GDA achieves $\mathcal O(\epsilon^{-4})$ complexity
for nonconvex--concave problems. Faster rates are available under a
uniform one-sided K{\L} condition~\cite{zheng2023universal}.
Without primal smoothing, PPALA reports
$\widetilde{\mathcal O}(\epsilon^{-2})$ complexity with fixed parameters
for functionally constrained optimization~\cite{kim2024fast}.
However, its feasibility argument relies on asymptotic reasoning,
and a corresponding finite-iteration feasibility bound remains
to be verified.

In all guarantees above that reach $\mathcal O(\epsilon^{-3})$, some
parameter grows with $1/\epsilon$: either the penalty, as in penalty methods
and GDPA, or the dual regularization, as in Perturbed Smoothed GDA. Practical
ALM implementations, in contrast, usually keep the penalty fixed. Growing the
penalty is costly for a single-loop method, because the Lipschitz constant of
$\nabla_{\bx}\calD_\rho$ increases with $\rho$ and forces the primal stepsize
to shrink. This leads to the following question.
\begin{quote}
  \textbf{(Q)}
  \emph{To reach $\mathcal O(\epsilon^{-3})$, must a single-loop linearized ALM
  let some parameter grow with $1/\epsilon$? If so, must this growth be
  carried by the penalty parameter?}
\end{quote}
In this paper, we provide a precise characterization addressing (Q). We study
SPLALM, a family of single-loop linearized ALMs that perform one gradient
descent--ascent step per iteration on $\calD_\rho$ with a fixed penalty
$\rho$, a dual regularization with parameter $s$, and moving primal and dual
centers (Definition~\ref{def:splalm}). Every member of the family has the
combined penalty $\rho+s$ (Lemma~\ref{lem:effective-penalty}). We answer the
first part of (Q) affirmatively and the second part negatively. If
$\rho+s\le c_1/\epsilon$, no choice of constant center stepsizes improves on
$\Omega(\epsilon^{-4})$ (Theorem~\ref{thm:lower-bound}). If $\rho$ is fixed
and $s$ is of order $1/\epsilon$, suitable center stepsizes give
$\mathcal O(\epsilon^{-3})$ (Theorem~\ref{thm:upper-beta0}). Hence some
parameter must grow, and the dual regularization alone can carry this growth.
The lower bound is new. The upper bounds adapt the analyses of Smoothed GDA
and Perturbed Smoothed GDA, both members of the family, thus demonstrating the
tightness of the lower bound and of the critical order of the combined
penalty. Table~\ref{tab:known-results} compares these results with known
guarantees. Our contribution can be summarized as follows: 
\begin{table*}[tp]
\centering
\begin{threeparttable}
\footnotesize
\setlength{\tabcolsep}{6pt}
\renewcommand{\arraystretch}{1.15}
\begin{tabular}{@{}lccc@{}}
\toprule
Method & Parameter scale & Upper bound & Lower bound \\
\midrule
\makecell[l]{Smoothed GDA\\\cite{zhang2020smoothed,li2026smoothing}}
 & Fixed primal smoothing
 & $\mathcal O(\epsilon^{-4})$ & $\Omega(\epsilon^{-4})$ \\[3pt]
\makecell[l]{Perturbed Smoothed GDA\\\cite{li2026smoothing}}
 & $s=\Theta(\epsilon^{-1})$
 & $\mathcal O(\epsilon^{-3})$ & $\Omega(\epsilon^{-3})$ \\[3pt]
% \midrule 
\makecell[l]{GDPA\\\cite{luSingleLoopGradientDescent2022}}
 & $\rho_k=\Theta(\epsilon^{-1})$
 & $\mathcal O(\epsilon^{-3})$ & --- \\[3pt]
\midrule
\rowcolor{gray!15}
\makecell[l]{\textbf{SPLALM family}\\Theorems~\ref{thm:upper-beta1} and~\ref{thm:lower-bound}}
 & \makecell{$\rho=\Theta(1)$,\\$\rho+s\le c_1\epsilon^{-1}$}
 & $\mathcal O(\epsilon^{-4})$ & $\Omega(\epsilon^{-4})$ \\[3pt]
\makecell[l]{\textbf{SPLALM family}\\Theorem~\ref{thm:upper-beta0}}
 & \makecell{$\rho=\Theta(1)$,\\$s=c_s\epsilon^{-1}$}
 & $\mathcal O(\epsilon^{-3})$ & --- \\[3pt]
\bottomrule
\end{tabular}
\caption{Complexity guarantees for representative single-loop methods.
For SPLALM, the bounds concern the family complexity in
Definition~\ref{def:complexity-measure} over
$\mathcal I(\mathbf C_0,\Delta)$, under the conditions of the cited theorems.
In the last row, $c_s$ is chosen as in Theorem~\ref{thm:upper-beta0}.
The rates use each method's stated stationarity criterion and assumptions;
lower bounds are algorithm-specific. ``---'' denotes an unreported bound.}
\label{tab:known-results}
\end{threeparttable}
\end{table*}
\begin{enumerate}
    \item \textbf{A lower bound independent of the center stepsizes.}
    We construct hard instances on which
    SPLALM, for every choice of constant center stepsizes
    $(\alpha,\beta)\in[0,1]^2$, requires $\Omega(\epsilon^{-4})$ iterations
    whenever $\epsilon(\rho+s)\le c_1$ (Theorem~\ref{thm:lower-bound}).
    Moving the centers therefore cannot compensate for a combined penalty
    $\rho+s=o(\epsilon^{-1})$.
    % TODO: GS version. A strictly convex perturbation of the outer function
    % would convert this bound into a lower bound of the same order for game
    % stationarity of the output pairs (x^{k+1},q^k). Restore this sentence
    % only after the perturbation lemma is proved.

    \item \textbf{Tightness for a fixed penalty.}
    For $\beta=1$, SPLALM reduces to Smoothed GDA applied to $\calD_\rho$ and
    attains $\mathcal O(\epsilon^{-4})$ uniformly over
    $\mathcal I(\mathbf C,\Delta)$ (Theorem~\ref{thm:upper-beta1}). Hence the
    lower bound is tight when $\rho$ is fixed and $\epsilon(\rho+s)\le c_1$,
    for sufficiently small $\gamma$, $\eta$, and $\tau$.

    \item \textbf{The critical scale of the combined penalty.}
    For $\beta=0$, $\rho$ fixed, and $s=c_s/\epsilon$, SPLALM reduces to
    Perturbed Smoothed GDA applied to $\calD_\rho$ and attains
    $\mathcal O(\epsilon^{-3})$ (Theorem~\ref{thm:upper-beta0}). Together with
    the lower bound, $\epsilon^{-1}$ is the critical order of $\rho+s$, and
    this growth can be carried by $s$ alone
    (Corollary~\ref{cor:effective-penalty}).
\end{enumerate}

% Open questions moved from the former (Q2) and (Q3); to be placed in the
% concluding section: (i) the complexity gain of primal smoothing under a
% fixed s and a moving dual center; (ii) structural conditions on
% (f,h,A) that yield a center--primal error bound and faster rates under a
% fixed combined penalty.

\paragraph{Organization}
Section~\ref{sec:preliminaries} introduces the problem class and the KKT and
game-stationarity criteria. Section~\ref{sec:algorithm-framework} introduces
the SPLALM family and the complexity measure, and shows that every member
of the family has the combined penalty
$\rho+s$. Section~\ref{sec:lower-bound} states the main results, and
Section~\ref{sec:lb-proof} proves the lower bound.
Section~\ref{sec:conclusion} discusses open questions.

\section{Preliminaries}
\label{sec:preliminaries}
% \subsection{Notation}
% \label{sec:notation}
\textbf{Notation.} We use $\langle\cdot,\cdot\rangle$ and $\|\cdot\|$ for the
Euclidean inner product and norm; for a matrix, $\|\cdot\|$ is the spectral
norm. For $\calA=(A_1,\ldots,A_m)^\top$, the Jacobian
$J_{\calA}(\bx)\in\RR^{m\times n}$ has $i$th row $\nabla A_i(\bx)^\top$. For a
nonempty closed convex set $C$, $\dist(\bu,C):=\min_{\bv\in C}\|\bu-\bv\|$, and
$\Pi_C(\bu)$ is the Euclidean projection of $\bu$ onto $C$. When convenient, we
identify $\RR^2$ with $\mathbb C$, so that $\|p\|=|p|$ and
$\langle p,q\rangle=\operatorname{Re}(\bar pq)$.
% Throughout the paper, $\partial h$ denotes the limiting
% subdifferential of $h$.

For a proper closed convex function \(h:\RR^m\to\RR\cup\{+\infty\}\), its
subdifferential, convex conjugate, and proximal mapping are
\[
    \partial h(\by):=\{\bv:\ h(\bw)\ge h(\by)+\langle\bv,\bw-\by\rangle
    \ \text{for all }\bw\},
    \qquad
    h^*(\bq):=\sup_{\by\in\RR^m}\{\langle\bq,\by\rangle-h(\by)\},
\]
\begin{equation}
    \prox_{t h}(\bu)
    :=\argmin_{\by\in\RR^m}
       \Bigl\{h(\by)+\frac{1}{2t}\|\by-\bu\|^2\Bigr\},
    \qquad t>0;
    \label{eq:prox-definition}
\end{equation}
see~\cite{rockafellar1970convex}. We write
\(\dom\partial h^*:=\{\bq:\partial h^*(\bq)\neq\varnothing\}\).


\subsection{Assumptions}
\label{sec:assumptions}
We impose the following assumptions.
% \begin{assumption}[Uniform Jacobian regularity]
% \label{ass:strong-licq}
% There exists a constant $\nu>0$ such that
% \[
%     \sigma_{\min}\bigl(J_{\calA}(\bx)\bigr)\ge \nu,
%     \qquad
%     \forall \bx\in\calS.
% \]
% Equivalently,
% \[
%     \|J_{\calA}(\bx)^\top \by\|
%     \ge
%     \nu \|\by\|,
%     \qquad \forall \bx\in\calS,\ \forall \by\in\RR^m.
% \]
% \end{assumption}
\begin{assumption}[Smooth components]
\label{ass:smoothness}
There exist constants $L_f,G_f,L_{\calA},G_{\calA}\ge0$ such that, for
all $\bx,\bx'\in\RR^n$,
\[
\begin{aligned}
    \|\nabla f(\bx)-\nabla f(\bx')\|
    &\le L_f\|\bx-\bx'\|,
    &\|\nabla f(\bx)\|&\le G_f,\\
    \|J_{\calA}(\bx)-J_{\calA}(\bx')\|
    &\le L_{\calA}\|\bx-\bx'\|,
    &\|J_{\calA}(\bx)\|&\le G_{\calA}.
\end{aligned}
\]
\end{assumption}

\begin{assumption}[Lipschitz outer function]
\label{ass:h}
The function $h:\RR^m\to\RR$ is convex and $L_h$-Lipschitz:
\[
    |h(\bu)-h(\bv)|\le L_h\|\bu-\bv\|,
    \qquad \bu,\bv\in\RR^m.
\]
\end{assumption}

\subsection{Problem class, stationarity, and oracle}
\label{subsec:setup}
To study~(Q), we introduce the following problem class.
\begin{definition}[Problem class]
\label{def:problem-class}
Let \(\Delta>0\) and
\(\mathbf C=(L_f,G_f,L_{\calA},G_{\calA},L_h)\in\RR^5_{\ge0}\). 
We denote by \(\mathcal I(\mathbf C,\Delta)\) the set of all triples
\((f,\calA,h)\), over all dimensions \(n,m\in\NN\), such that
Assumptions~\ref{ass:smoothness} and~\ref{ass:h} hold with the constants
in \(\mathbf C\) and 
\(
    F(\zero)-\inf_{\bx\in\RR^n}F(\bx)\le\Delta .
\)
\end{definition}
The constants \(\mathbf C\) and \(\Delta\) do not depend on the target
accuracy or on \(n,m\), and the gap condition implies \(\inf F>-\infty\).
% The gap condition implies 
% \(\inf F>-\infty\). Assumption~\ref{ass:level-boundedness} is not part of
% \(\mathcal I(\mathbf C,\Delta)\); it is imposed separately in the complexity
% results that require it. 

{We use the following stationarity measure for Problem~\eqref{prob:main}.} 
\begin{definition}[\(\epsilon\)-KKT point]
\label{def:eps-kkt}
Given \(\epsilon>0\). A triple \((\bx,\by,\bq)\) with
\(\bq\in\partial h(\by)\) is an \(\epsilon\)-KKT point of
\eqref{prob:main} if
\begin{equation}
    \mathcal{R}_{\mathrm{KKT}}(\bx,\by,\bq)
    :=
    \max\Bigl\{
        \bigl\|\nabla f(\bx)+J_{\calA}(\bx)^\top\bq\bigr\|,\;
        \|\calA(\bx)-\by\|
    \Bigr\} \leq \epsilon
    \label{eq:kkt-residual}
\end{equation}
\end{definition}

We next relate Definition~\ref{def:eps-kkt} to game stationarity (GS) for the
minimax reformulation~Problem~\eqref{eq:ncl} of Problem~\eqref{prob:main}. We use the
following GS definition from \cite[Definition~3.1(b)]{li2025nonsmooth} and
\cite[(2.10)--(2.11)]{zhang2026matching}.
\begin{definition}[\(\epsilon\)-GS point]
\label{def:eps-gs}
Let \(\epsilon>0\). A pair \((\bx,\bq)\in\RR^n\times\dom\partial h^*\) is an
\(\epsilon\)-game stationary point (\(\epsilon\)-GS) of~Problem~\eqref{eq:ncl} if
\[
    \calG(\bx,\bq)
    :=
    \max\Bigl\{
        \bigl\|\nabla f(\bx)+J_{\calA}(\bx)^\top\bq\bigr\|,\;
        \dist\bigl(\calA(\bx),\partial h^*(\bq)\bigr)
    \Bigr\}
    \le\epsilon .
\]
\end{definition}
We now relate $\epsilon$-KKT points of Problem~\eqref{prob:main}
to $\epsilon$-GS points of its Fenchel saddle reformulation~Problem~\eqref{eq:ncl}. 
\begin{proposition}[KKT points of Problem~\eqref{prob:main} and GS points of~Problem~\eqref{eq:ncl}]
\label{prop:kkt-gs}
Suppose that Assumption~\ref{ass:h} holds, and let
\((\bx,\bq)\in\RR^n\times\dom\partial h^*\). Then the following measurement relationship holds: 
\[
    \calG(\bx,\bq)
    =
    \min_{\by\in\partial h^*(\bq)}
    \mathcal{R}_{\mathrm{KKT}}(\bx,\by,\bq),
\]
with minimum attained at
\(\hat\by=\Pi_{\partial h^*(\bq)}\bigl(\calA(\bx)\bigr)\). Consequently, for
every \(\epsilon>0\):
\begin{enumerate}[label=(\alph*)]
\item if \((\bx,\by,\bq)\) is an \(\epsilon\)-KKT point
  of~Problem~\eqref{prob:main}, then \((\bx,\bq)\) is an \(\epsilon\)-GS point
  of~Problem~\eqref{eq:ncl};
\item if \((\bx,\bq)\) is an \(\epsilon\)-GS point of~Problem~\eqref{eq:ncl}, then
  \((\bx,\hat\by,\bq)\) is an \(\epsilon\)-KKT point of~Problem~\eqref{prob:main}.
\end{enumerate}
\end{proposition}

\begin{proof}
Since \(h\) is proper, closed, and convex,
\(\bq\in\partial h(\by)\) if and only if \(\by\in\partial h^*(\bq)\)
\cite[Theorem~23.5]{rockafellar1970convex}. The first term of
\(\mathcal{R}_{\mathrm{KKT}}\) does not depend on \(\by\). The set
\(\partial h^*(\bq)\) is nonempty, closed, and convex, so
\(\min_{\by\in\partial h^*(\bq)}\|\calA(\bx)-\by\|
=\dist(\calA(\bx),\partial h^*(\bq))\), attained at \(\hat\by\).
Statements~(a) and~(b) follow.
\end{proof}

We next introduce the following deterministic first-order oracle model.
\begin{definition}[Oracle model]
\label{def:oracle-model}
We assume access to a deterministic first-order oracle
\[
    \mathcal{O}_{f,\mathcal{A}}(\bx)
    :=
    \bigl(f(\bx),\nabla f(\bx),\mathcal{A}(\bx),J_{\mathcal{A}}(\bx)\bigr),
\]
and an exact proximal oracle
\[
    \mathcal{O}_{h}(\bv,t)
    :=
    \operatorname{prox}_{t h}(\bv),
    \qquad t>0.
\]
Each call to either oracle counts as one oracle evaluation. In particular,
the cost of solving the proximal subproblem is represented by one exact
proximal-oracle call.
\end{definition}


\section{Algorithm Family and Effective Penalty}
\label{sec:algorithm-framework}

\subsection{Algorithm family and worst-case iteration complexity}
\label{subsec:algorithm-family}

For \(\rho>0\), the reduced augmented Lagrangian of~\eqref{prob:split} is
\begin{equation}
    \calD_\rho(\bx,\blmbd)
    :=\min_{\by\in\RR^m}\Bigl\{f(\bx)+h(\by)+\langle\blmbd,\calA(\bx)-\by\rangle
    +\frac{\rho}{2}\|\calA(\bx)-\by\|^2\Bigr\}.
    \label{eq:reduced-al}
\end{equation}
Its minimizer and the associated multiplier are
\begin{equation}
    \by(\bx,\blmbd):=\prox_{h/\rho}\Bigl(\calA(\bx)+\frac{\blmbd}{\rho}\Bigr),
    \qquad
    \bq(\bx,\blmbd):=\blmbd+\rho\bigl[\calA(\bx)-\by(\bx,\blmbd)\bigr],
    \label{eq:prox-multiplier}
\end{equation}
and \(\bq(\bx,\blmbd)\in\partial h(\by(\bx,\blmbd))\).
We study the smoothed perturbed linearized augmented Lagrangian method
(SPLALM), which performs alternating gradient descent--ascent steps on the
smoothed perturbed potential function
\begin{equation}
    \mathcal P_s(\bx,\blmbd;\bz,\bmu)
    :=\calD_\rho(\bx,\blmbd)+\frac{1}{2\gamma}\|\bx-\bz\|^2
    -\frac{1}{2s}\|\blmbd-\bmu\|^2 .
    \label{eq:potential}
\end{equation}
The proximal term smooths the primal problem, the term
\(-\frac{1}{2s}\|\blmbd-\bmu\|^2\) perturbs the dual problem, and
\(\bz\in\RR^n\) and \(\bmu\in\RR^m\) are the primal and dual centers.

\begin{definition}[SPLALM]
\label{def:splalm}
Let \(\theta=(\rho,s,\gamma,\eta,\tau)\) satisfy \(\rho>0\),
\(0<\eta\le\gamma<\infty\), and \(0<\tau\le s<\infty\), and let
\((\alpha,\beta)\in[0,1]^2\). The method \(\mathrm{SPLALM}_{\theta,\alpha,\beta}\)
starts from \((\bx^0,\blmbd^0,\bz^0,\bmu^0)=\zero\) and, for \(k\ge0\), sets
\begin{equation}
\begin{aligned}
    \bx^{k+1}&=\bx^k-\eta\,\nabla_{\bx}\mathcal P_s(\bx^k,\blmbd^k;\bz^k,\bmu^k),\\
    \blmbd^{k+1}&=\blmbd^k+\tau\,\nabla_{\blmbd}\mathcal P_s(\bx^{k+1},\blmbd^k;\bz^k,\bmu^k),\\
    \bz^{k+1}&=\bz^k+\alpha(\bx^{k+1}-\bz^k),\\
    \bmu^{k+1}&=\bmu^k+\beta(\blmbd^{k+1}-\bmu^k).
\end{aligned}
    \label{eq:splalm-updates}
\end{equation}
At iteration \(k\), it outputs
\((\bx^{k+1},\by^{k+1},\bq^k)
:=(\bx^{k+1},\by(\bx^{k+1},\blmbd^k),\bq(\bx^{k+1},\blmbd^k))\).
\end{definition}

With first-order information cached between iterations, \(K\) iterations and
their outputs use \(K+1\) calls to \(\mathcal O_{f,\calA}\) and \(2K\) calls to
\(\mathcal O_h\). Thus the total oracle count is \(3K+1\), and iteration
and oracle complexity agree up to a factor of three, apart from one initial
first-order oracle call.

\begin{definition}[SPLALM family]
\label{def:algorithm-family}
For \(\theta\) as in Definition~\ref{def:splalm},
\(\calC_\theta:=\{\mathrm{SPLALM}_{\theta,\alpha,\beta}:(\alpha,\beta)\in[0,1]^2\}\).
\end{definition}

\begin{remark}[Initialization]
\label{rem:lb-initialization}
Every method in \(\calC_\theta\) starts from zero on every instance.
Theorem~\ref{thm:lower-bound} constructs a hard instance for each
\((\theta,\alpha,\beta,\epsilon)\) while keeping this initialization fixed.
Thus, the lower bound holds for the prescribed starting point.
The constructions in~\cite[Section~5.1.1]{li2026smoothing}
and~\cite[Appendix~C.1]{mahdavinia2022tight} choose both the hard function
and the initialization.
\end{remark}

\begin{definition}[Worst-case iteration complexity]
\label{def:complexity-measure}
Let \(\epsilon>0\) and \(I\in\mathcal I(\mathbf C,\Delta)\), and let
\((\bx^{k+1},\by^{k+1},\bq^k)\) be the outputs of
\(\mathrm{SPLALM}_{\theta,\alpha,\beta}\) on \(I\). Define
\[
\begin{aligned}
    T^{\mathrm{KKT}}_\epsilon(\theta,\alpha,\beta;I)
    &:=\inf\Bigl\{K\ge1:\ \min_{0\le k<K}
        \mathcal R_{\mathrm{KKT}}(\bx^{k+1},\by^{k+1},\bq^k)\le\epsilon\Bigr\},\\
    T^{\mathrm{GS}}_\epsilon(\theta,\alpha,\beta;I)
    &:=\inf\Bigl\{K\ge1:\ \min_{0\le k<K}
        \calG(\bx^{k+1},\bq^k)\le\epsilon\Bigr\},
\end{aligned}
\]
with \(\inf\varnothing=+\infty\). For \(\bullet\in\{\mathrm{KKT},\mathrm{GS}\}\),
the worst-case iteration complexity of \(\calC_\theta\) over
\(\mathcal I(\mathbf C,\Delta)\) is
\[
    T^{\bullet}_\epsilon\bigl(\calC_\theta,\mathcal I(\mathbf C,\Delta)\bigr)
    :=\inf_{(\alpha,\beta)\in[0,1]^2}\
    \sup_{I\in\mathcal I(\mathbf C,\Delta)}\
    T^{\bullet}_\epsilon(\theta,\alpha,\beta;I).
\]
\end{definition}

The infimum precedes the supremum, so \((\alpha,\beta)\) may depend on
\(\epsilon\), \(\theta\), and \((\mathbf C,\Delta)\), but not on the instance.
At least one of the first \(K\) outputs meets the tolerance if and only if
\(K\ge T^{\bullet}_\epsilon\); hence the definition covers every rule that
returns one of these outputs, including the best-so-far rule. By
Proposition~\ref{prop:kkt-gs},
\(\calG(\bx^{k+1},\bq^k)\le\mathcal R_{\mathrm{KKT}}(\bx^{k+1},\by^{k+1},\bq^k)\)
for every \(k\), so
\(T^{\mathrm{GS}}_\epsilon(\calC_\theta,\mathcal I)
\le T^{\mathrm{KKT}}_\epsilon(\calC_\theta,\mathcal I)\). An upper bound for
the KKT measure therefore yields the same upper bound for the GS measure, and
a lower bound for the GS measure yields the same lower bound for the KKT
measure.

\begin{remark}[Family complexity and oracle complexity]
\label{rem:family-not-oracle}
Let \(\mathcal A_{\det}\) be the class of deterministic methods whose iterates
are arbitrary maps of all previous responses of \(\mathcal O_{f,\calA}\) and
\(\mathcal O_h\), defined as in~\cite[Section~2.2]{carmon2020lower}
and~\cite[Definition~2.5]{pan2026optimal} with the oracles of
Definition~\ref{def:oracle-model}. Then \(\calC_\theta\subset\mathcal A_{\det}\), and
\[
    T_\epsilon(\mathcal A_{\det},\mathcal I)
    =\inf_{A\in\mathcal A_{\det}}\sup_{I\in\mathcal I}T_\epsilon(A;I)
    \le\inf_{(\alpha,\beta)\in[0,1]^2}\sup_{I\in\mathcal I}
    T_\epsilon(\theta,\alpha,\beta;I)
    =T_\epsilon(\calC_\theta,\mathcal I).
\]
\begin{enumerate}[label=(\roman*)]
\item (Scope) A lower bound on \(T_\epsilon(\calC_\theta,\mathcal I)\) does not
  bound \(T_\epsilon(\mathcal A_{\det},\mathcal I)\); methods outside
  \(\calC_\theta\) may be faster. When \(h\) is the support function of a
  compact convex set, \eqref{eq:ncl} is a smooth nonconvex--concave problem,
  and the \(O(\epsilon^{-5/2})\) GS bound of~\cite{zhang2026matching} and the
  \(O(\epsilon^{-3})\) OS bound of~\cite{pan2026optimal} are attained by
  methods outside \(\calC_\theta\).
\item (Uniformity) The infimum ranges over all \((\alpha,\beta)\in[0,1]^2\),
  including the endpoint cases in Table~\ref{tab:special-cases}, and
  Theorem~\ref{thm:lower-bound} holds for all \((\rho,s)\) with
  \(\epsilon(\rho+s)\le c_1\). A lower bound on
  \(T_\epsilon(\calC_\theta,\mathcal I)\) therefore cannot be removed by
  tuning the center stepsizes; it is governed by the effective penalty
  \(\rho+s\) of Lemma~\ref{lem:effective-penalty}.
\item (Transfer) Conversely, every lower bound on
  \(T_\epsilon(\mathcal A_{\det},\mathcal I)\) for the same problem class and
  measure also holds for \(\calC_\theta\).
\end{enumerate}
\end{remark}

\subsection{Effective penalty and special cases}
\label{subsec:effective-penalty}

The reduced augmented Lagrangian is a dual proximal smoothing of the saddle
function \(\calL(\bx,\bq)=f(\bx)+\langle\calA(\bx),\bq\rangle-h^*(\bq)\)
in~\eqref{eq:ncl}~\cite{rockafellar1976augmented}. The center of this
smoothing is the multiplier, and its strength is the penalty. Adding the dual
regularizer of~\eqref{eq:potential} keeps this form; it only raises the
penalty and moves the center.

\begin{lemma}[Effective penalty]
\label{lem:effective-penalty}
Suppose that Assumption~\ref{ass:h} holds. For every \(\sigma>0\),
\begin{equation}
    \calD_\sigma(\bx,\blmbd)
    =\max_{\bq\in\RR^m}\Bigl\{\calL(\bx,\bq)-\frac{1}{2\sigma}\|\bq-\blmbd\|^2\Bigr\},
    \label{eq:dual-smoothing}
\end{equation}
and consequently
\begin{equation}
    \max_{\blmbd\in\RR^m}\mathcal P_s(\bx,\blmbd;\bz,\bmu)
    =\calD_{\rho+s}(\bx,\bmu)+\frac{1}{2\gamma}\|\bx-\bz\|^2 .
    \label{eq:effective-penalty}
\end{equation}
\end{lemma}

\begin{proof}
Write \(\bu=\calA(\bx)\) and \(h(\by)=\max_{\bq}\{\langle\bq,\by\rangle-h^*(\bq)\}\).
Since \(\dom h^*\) lies in a ball by Assumption~\ref{ass:h}, Sion's minimax
theorem allows the minimum over \(\by\) and the maximum over \(\bq\) in
\(\min_{\by}\max_{\bq}\{\langle\bq-\blmbd,\by\rangle-h^*(\bq)
+\frac{\sigma}{2}\|\bu-\by\|^2\}\) to be interchanged. The inner minimizer is
\(\by=\bu-(\bq-\blmbd)/\sigma\), which gives~\eqref{eq:dual-smoothing}.
Applying~\eqref{eq:dual-smoothing} with \(\sigma=\rho\) and exchanging the two
maximizations,
\[
    \max_{\blmbd}\mathcal P_s
    =\max_{\bq}\Bigl\{\calL(\bx,\bq)-\min_{\blmbd}\Bigl[\frac{\|\bq-\blmbd\|^2}{2\rho}
    +\frac{\|\blmbd-\bmu\|^2}{2s}\Bigr]\Bigr\}+\frac{\|\bx-\bz\|^2}{2\gamma},
\]
and the inner minimum equals \(\|\bq-\bmu\|^2/(2(\rho+s))\).
Then~\eqref{eq:dual-smoothing} with \(\sigma=\rho+s\)
gives~\eqref{eq:effective-penalty}.
\end{proof}

Lemma~\ref{lem:effective-penalty} shows that every member of
\(\calC_\theta\) has the effective penalty \(\rho+s\), whatever the center
stepsizes are. The center stepsize \(\beta\) only places the dual center
between two classical choices, and \(\alpha\) decides whether the primal
variable is smoothed.

\emph{Zero center} (\(\beta=0\)). Then \(\bmu^k=\zero\), and
\(\calD_{\rho+s}(\bx,\zero)=f(\bx)+\min_{\by}\{h(\by)
+\frac{\rho+s}{2}\|\calA(\bx)-\by\|^2\}\) is the quadratic penalty function
with parameter \(\rho+s\). The dual step carries the perturbation
\(-\blmbd^k/s\), as in perturbed GDA~\cite{lu2020hybrid,xu2023unified} and
perturbed smoothed GDA~\cite{li2026smoothing}.

\emph{Current multiplier as center} (\(\beta=1\)). Then
\(\bmu^k=\blmbd^k\), the gradient of the dual regularizer vanishes at every
dual step, and
\[
    \blmbd^{k+1}=\Bigl(1-\frac{\tau}{\rho}\Bigr)\blmbd^k+\frac{\tau}{\rho}\,\bq^k,
    \qquad
    \bq^k=\prox_{\rho h^*}\bigl(\blmbd^k+\rho\calA(\bx^{k+1})\bigr).
\]
For \(\tau=\rho\), this is the multiplier update of the augmented Lagrangian
method~\cite{hestenes1969multiplier,powell1969method,rockafellar1973multiplier},
that is, a proximal ascent step on \(\calL(\bx^{k+1},\cdot)\)
by~\eqref{eq:dual-smoothing}. The iterates do not depend on \(s\), so this
choice also reproduces the case without dual regularizer, \(s=\infty\).

\emph{Moving center} (\(0<\beta<1\)). The dual center is an exponentially
weighted average of past multipliers, and the method has the structure of doubly smoothed
GDA~\cite{zheng2023universal}.

Table~\ref{tab:special-cases} collects these cases. The identities
\(\bmu^k=\zero\), \(\bmu^k=\blmbd^k\), and \(\bz^k=\bx^k\) for
\(\beta=0\), \(\beta=1\), and \(\alpha=1\), respectively, follow by induction
from the zero initialization. Every entry refers to the cited update structure
applied to \(\calD_\rho\) with \(\rho>0\); methods acting on \(\calL\) itself
are compared in Section~\ref{subsec:gda-embedding}.

\begin{table}[t]
\centering
\small
\caption{Special cases of \(\calC_\theta\). All methods act on
\(\calD_\rho\) with \(\rho>0\) and have the effective penalty \(\rho+s\).}
\label{tab:special-cases}
\begin{tabular}{@{}llll@{}}
\toprule
\(\beta\) & Dual center & \(\alpha=1\) & \(0<\alpha<1\) \\
\midrule
\(0\) & \(\zero\)
  & perturbed GDA~\cite{lu2020hybrid,xu2023unified}
  & perturbed smoothed GDA~\cite{li2026smoothing} \\
\((0,1)\) & \(\sum_{j=1}^{k}\beta(1-\beta)^{k-j}\blmbd^j\)
  & ---
  & DS-GDA~\cite{zheng2023universal}, \(r_1=1/\gamma\), \(r_2=1/s\) \\
\(1\) & \(\blmbd^k\)
  & linearized ALM (\(\tau=\rho\))
  & smoothed GDA~\cite{zhang2020smoothed} \\
\bottomrule
\end{tabular}
\end{table}


\section{Main Results}
\label{sec:lower-bound}

\subsection{Main result}
\label{subsec:lb-main}

We use the problem class of Definition~\ref{def:problem-class}, the family
\(\calC_\theta\) of Definition~\ref{def:algorithm-family}, and the complexity
of Definition~\ref{def:complexity-measure}, with the normalized constants
\(\mathbf C_0:=(1,1,0,1,1)\). Since \(L_{\calA}=0\), every instance in
\(\mathcal I(\mathbf C_0,\Delta)\) has an affine inner mapping.

We now present the iteration complexity of \(\calC_\theta\) under the choice
\(\beta=1\). We choose this setting for two reasons. First, when
\(\beta=1\), SPLALM reduces to Smoothed GDA applied to \(\calD_\rho\)
(Table~\ref{tab:special-cases}), whose iteration complexity is known. Second,
Theorem~\ref{thm:lower-bound} below holds for any arbitrary choice of
\((\alpha,\beta)\in[0,1]^2\), including \(\beta=1\). Stating the upper bound
under the same setting yields a direct tightness guarantee.

\begin{theorem}[Iteration complexity of SPLALM with \(\beta=1\)]
\label{thm:upper-beta1}
There exist \(\bar\gamma>0\), depending only on \(\mathbf C\) and \(\rho\),
and \(\bar\tau>0\), depending only on \(\mathbf C\), \(\rho\), and
\(\eta/\gamma\), such that the following holds. Let \(\gamma\le\bar\gamma\),
\(\eta\le\gamma/2\), \(\tau\le\bar\tau\), and \(s\ge\tau\). Then there exist
\(c>0\), depending only on \(\theta\) and \(\mathbf C\), and \(C>0\),
depending only on \(\theta\), \(\mathbf C\), and \(\Delta\), such that, given
any \(\epsilon\in(0,1]\) and \(\alpha:=c\,\epsilon^2\),
\[
T^{\mathrm{KKT}}_\epsilon(\theta,\alpha,1;I)\le C\,\epsilon^{-4}
\qquad\text{for all } I\in\mathcal I(\mathbf C,\Delta).
\]
\end{theorem}
The proof is given in Appendix~\ref{app:upper-beta1}.

\begin{remark}
\label{rem:upper-beta1}
(i) When \(\beta=1\), SPLALM reduces to Smoothed GDA applied to
\(\calD_\rho\). A result analogous to Theorem~\ref{thm:upper-beta1} has been
obtained in~\cite[Theorem~3.4]{zhang2020smoothed}
and~\cite[Theorem~B.1]{li2026smoothing}. However, the initial gap there
depends on the initialization, the dual domain is a given compact set, and
the stationarity measure is evaluated at \((\bx^{k+1},\blmbd^{k+1})\).
Theorem~\ref{thm:upper-beta1} bounds the initial gap by the constants of
\(\mathcal I(\mathbf C,\Delta)\), keeps the multipliers in the interior of a
ball without any projection, and evaluates the KKT residual at
\((\bx^{k+1},\blmbd^k)\), the point that SPLALM outputs.
(ii) The iterates do not depend on \(s\) when \(\beta=1\). Hence the constants
\(c\) and \(C\) in Theorem~\ref{thm:upper-beta1} can be chosen independently
of \(s\): the choices for \(s=\tau\) apply to every \(s\ge\tau\).
The upper bound
requires only one member of \(\calC_\theta\); for \(\beta=0\), the instance
in Section~\ref{subsec:lb-huber} shows that no upper bound holds when
\(\epsilon(\rho+s)\le c_1\).
\end{remark}

\begin{theorem}[Lower bound uniform in \(\rho\) and \(s\)]
\label{thm:lower-bound}
Fix \(0<\eta\le\gamma\), \(\tau>0\), and \(\Delta>0\). There exist constants
\(c_1=c_1(\gamma,\eta,\tau)>0\) and \(c_2=c_2(\gamma,\tau,\Delta)>0\),
independent of \(\rho\), \(s\), \(\alpha\), and \(\beta\), such that the
following holds. For every \(\epsilon\in(0,c_2)\), every \(\rho>0\) and
\(s\ge\tau\) with \(\epsilon(\rho+s)\le c_1\), and every
\((\alpha,\beta)\in[0,1]^2\), there exists an instance
\(I\in\mathcal I(\mathbf C_0,\Delta)\) on which
\(\mathrm{SPLALM}_{\theta,\alpha,\beta}\) requires at least
\[
    \frac{\Delta\,\epsilon^{-4}}{16384\,\gamma\max\{1,4\tau,4\gamma\tau\}}
\]
iterations before one of its outputs is an \(\epsilon\)-KKT point.
Equivalently, 
\[T^{\mathrm{KKT}}_\epsilon(\calC_\theta,\mathcal I(\mathbf C_0,\Delta))
\ge\Delta\epsilon^{-4}/(16384\,\gamma\max\{1,4\tau,4\gamma\tau\})\].
\end{theorem}
Explicit choices of \(c_1\) and \(c_2\) are given in~\eqref{eq:rot-thresholds}.
% \todo{State the GS version $T^{\mathrm{GS}}_\epsilon$ after the strict-convex perturbation lemma is proved.}
% perturbation lemma is proved.

\begin{theorem}[Iteration complexity of SPLALM with \(\beta=0\)]
\label{thm:upper-beta0}
Fix \(\rho,\gamma,\eta,\tau>0\) satisfying the step-size
conditions~\eqref{eq:beta0-stepsizes} for all sufficiently large \(s\).
There exist constants \(c_s>0\) and \(c_\alpha>0\), depending only on
\(\rho,\gamma,\eta,\tau\) and \(\mathbf C\), and \(C'>0\), depending in
addition on \(\Delta\), such that the following holds. Given any
sufficiently small \(\epsilon>0\), let
\[
s:=\frac{c_s}{\epsilon},\qquad
\alpha:=c_\alpha\epsilon,\qquad \beta:=0,\qquad
\theta(\epsilon):=(\rho,s,\gamma,\eta,\tau).
\]
Then
\[
T^{\mathrm{KKT}}_\epsilon(\theta(\epsilon),\alpha,0;I)
\le C'\epsilon^{-3}
\qquad\text{for all }I\in\mathcal I(\mathbf C,\Delta).
\]
All parameters remain constant within each run. The coefficient \(c_s\)
is chosen sufficiently large, and \(c_\alpha\) sufficiently small, as
specified in Appendix~\ref{app:upper-beta0}.
\end{theorem}
The proof is given in Appendix~\ref{app:upper-beta0}.

\begin{remark}
\label{rem:upper-beta0}
(i) When \(\beta=0\), SPLALM reduces to Perturbed Smoothed GDA applied to
\(\calD_\rho\). A result analogous to Theorem~\ref{thm:upper-beta0} has
been obtained in~\cite[Theorem~3.3]{li2026smoothing}.
Theorem~\ref{thm:upper-beta0} bounds the initial gap by the constants of
\(\mathcal I(\mathbf C,\Delta)\), keeps the multipliers in the interior of
a ball without any projection, and evaluates the KKT residual at
\((\bx^{k+1},\blmbd^k)\). The change in the evaluation point contributes
an additional term \(L_h/(10s)\) in~\eqref{eq:beta0-residual-transfer}.
(ii) The coefficient \(c_s\) must be sufficiently large. For fixed \(\rho\),
the choice \(s=c/\epsilon\) with \(0<c<c_1\) satisfies
\(\epsilon(\rho+s)\le c_1\) for all sufficiently small \(\epsilon\), and
Theorem~\ref{thm:lower-bound} applies.
(iii) The primal center stepsize is \(\alpha=\Theta(\epsilon)\), while the
dual center remains zero. Whether a primal center stepsize independent of
\(\epsilon\) attains the same order is outside the scope of this result.
\end{remark}

\begin{corollary}[Scale of the combined penalty]
\label{cor:effective-penalty}
Fix \(0<\eta\le\gamma\), \(\tau>0\), and \(\Delta>0\), and let \(c_1\)
be the constant in Theorem~\ref{thm:lower-bound}.
\begin{enumerate}[label=(\alph*)]
\item Let \(\rho=\rho(\epsilon)>0\) and \(s=s(\epsilon)\ge\tau\) satisfy
\(\rho+s\le c_1/\epsilon\) for all sufficiently small \(\epsilon\).
Then
\[
T^{\mathrm{KKT}}_\epsilon
\bigl(\calC_{\theta(\epsilon)},\mathcal I(\mathbf C_0,\Delta)\bigr)
=\Omega(\epsilon^{-4}).
\]
\item Fix \(\rho>0\) and let \(\rho,\gamma,\eta,\tau\) satisfy the
conditions of Theorem~\ref{thm:upper-beta0} for the class constants
\(\mathbf C_0\). For \(s=c_s/\epsilon\) with the sufficiently large
coefficient specified there, the choice \(\beta=0\) and
\(\alpha=c_\alpha\epsilon\) gives
\[
T^{\mathrm{KKT}}_\epsilon
\bigl(\calC_{\theta(\epsilon)},\mathcal I(\mathbf C_0,\Delta)\bigr)
=O(\epsilon^{-3}).
\]
\end{enumerate}
Thus an \(O(\epsilon^{-3})\) guarantee over
\(\mathcal I(\mathbf C_0,\Delta)\) requires
\(\rho(\epsilon)+s(\epsilon)>c_1/\epsilon\) for all sufficiently small
\(\epsilon\), and is attained with fixed \(\rho\) and
\(s=c_s/\epsilon\) under the conditions in~(b).
The constants \(c_1\) and \(c_s\) differ; the conclusion identifies the
order \(\epsilon^{-1}\) of the required combined penalty.
\end{corollary}
\begin{proof}
Part~(a) follows from Theorem~\ref{thm:lower-bound}. Part~(b) follows from
Theorem~\ref{thm:upper-beta0}, since the infimum in
Definition~\ref{def:complexity-measure} includes
\((\alpha,\beta)=(c_\alpha\epsilon,0)\). If
\(\rho(\epsilon)+s(\epsilon)\le c_1/\epsilon\) along a sequence of
accuracies tending to zero, the lower bound in
Theorem~\ref{thm:lower-bound} precludes an \(O(\epsilon^{-3})\) guarantee
along that sequence. This proves the necessary condition.
\end{proof}

\section{Hard Instances and Proof of Theorem~\ref{thm:lower-bound}}
\label{sec:lb-proof}

In Section~\ref{sec:lower-bound}, we established a lower bound on the iteration
complexity of the SPLALM family that holds for every choice of constant center
stepsizes. In this section, we prove this result by constructing hard instances.

Fix \(\epsilon\) and \(\theta\) satisfying the conditions of
Theorem~\ref{thm:lower-bound}. The lower bound has the quantified form
\[
\forall\,(\alpha,\beta)\in[0,1]^2,\quad
\exists\,I\in\mathcal I(\mathbf C_0,\Delta):\quad
T^{\mathrm{KKT}}_\epsilon(\theta,\alpha,\beta;I)
\geq \Omega(\Delta\epsilon^{-4}).
\]
Thus, different choices of \((\alpha,\beta)\) may require different hard
instances. For each pair, we construct an instance in
\(\mathcal I(\mathbf C_0,\Delta)\) on which
\(\mathrm{SPLALM}_{\theta,\alpha,\beta}\), started from the zero
initialization prescribed in Definition~\ref{def:splalm}, requires at least
the stated number of iterations to produce an \(\epsilon\)-KKT output.
For the outputs of \(\mathrm{SPLALM}_{\theta,\alpha,\beta}\), we write
\(R_k:=\mathcal R_{\mathrm{KKT}}(\bx^{k+1},\by^{k+1},\bq^k)\) for \(k\ge0\).

We split \([0,1]^2\) into three parts by the threshold
\[
\alpha_\epsilon:=1024\max\{1,4\tau,4\gamma\tau\}\,\epsilon^2,
\]
namely
\[
\begin{aligned}
[0,1]^2
={}&\bigl\{(\alpha,\beta)\in[0,1]^2:\alpha\le\alpha_\epsilon\bigr\}\\
&\cup\bigl\{(\alpha,\beta)\in[0,1]^2:\alpha>\alpha_\epsilon,\ \beta>0\bigr\}\\
&\cup\bigl\{(\alpha,\beta)\in[0,1]^2:\alpha>\alpha_\epsilon,\ \beta=0\bigr\},
\end{aligned}
\]
and use a separate hard instance for each part. For
\(\alpha\le\alpha_\epsilon\), the bound obtained in
Section~\ref{subsec:lb-plateau} is of order \(\Delta\epsilon^{-4}\). This
threshold ensures the parameter estimates required by the rotating instance
in Section~\ref{subsec:lb-rotation}.

\subsection{Reduction to a gradient descent--ascent recursion}
\label{subsec:gda-embedding}

Where \(\by(\bx,\blmbd)=\zero\), the updates of
\(\mathrm{SPLALM}_{\theta,\alpha,\beta}\) reduce to gradient
descent--ascent steps on a smooth saddle function. To state this, define
\[
\mathcal R_\rho:=\bigl\{(\bx,\blmbd):\ \blmbd+\rho\calA(\bx)\in\partial h(\zero)\bigr\},
\qquad
\Psi_\rho(\bx,\blmbd):=f(\bx)+\frac{\rho}{2}\|\calA(\bx)\|^2+\langle\blmbd,\calA(\bx)\rangle .
\]

\begin{proposition}[Reduction on \(\mathcal R_\rho\)]
\label{prop:reduction}
Let \(\rho>0\) and \((\bx,\blmbd)\in\mathcal R_\rho\). Then
\(\by(\bx,\blmbd)=\zero\), \(\bq(\bx,\blmbd)=\blmbd+\rho\calA(\bx)\), and
\[
\nabla f(\bx)+J_{\calA}(\bx)^\top\bq(\bx,\blmbd)=\nabla_{\bx}\Psi_\rho(\bx,\blmbd),
\qquad
\calA(\bx)-\by(\bx,\blmbd)=\nabla_{\blmbd}\Psi_\rho(\bx,\blmbd).
\]
Moreover,
\(\calD_\rho(\bx,\blmbd)=\Psi_\rho(\bx,\blmbd)+h(\zero)\).
\end{proposition}

\begin{proof}
Set \(\bv:=\calA(\bx)+\blmbd/\rho\). By~\eqref{eq:prox-multiplier}, \(\by(\bx,\blmbd)\)
is the unique minimizer of \(h(\by)+\frac{\rho}{2}\|\by-\bv\|^2\), and
\(\by\) is this minimizer if and only if \(\rho(\bv-\by)\in\partial h(\by)\).
For \(\by=\zero\) the inclusion reads
\(\blmbd+\rho\calA(\bx)\in\partial h(\zero)\), which holds because
\((\bx,\blmbd)\in\mathcal R_\rho\). Hence \(\by(\bx,\blmbd)=\zero\) and
\(\bq(\bx,\blmbd)=\blmbd+\rho\calA(\bx)\). Since
\(\nabla_{\bx}\bigl(\frac{\rho}{2}\|\calA(\bx)\|^2\bigr)=\rho J_{\calA}(\bx)^\top\calA(\bx)\),
\[
\nabla f(\bx)+J_{\calA}(\bx)^\top\bigl(\blmbd+\rho\calA(\bx)\bigr)
=\nabla_{\bx}\Psi_\rho(\bx,\blmbd),
\]
and \(\calA(\bx)-\zero=\nabla_{\blmbd}\Psi_\rho(\bx,\blmbd)\).
Finally, completing the square in~\eqref{eq:reduced-al} gives
\(\calD_\rho(\bx,\blmbd)=f(\bx)+\min_{\by}\{h(\by)+\frac{\rho}{2}\|\by-\bv\|^2\}
-\frac{1}{2\rho}\|\blmbd\|^2\), and the minimum equals
\(h(\zero)+\frac{\rho}{2}\|\bv\|^2\) because its minimizer is \(\zero\). Hence
\[
\calD_\rho(\bx,\blmbd)
=f(\bx)+h(\zero)+\frac{\rho}{2}\|\calA(\bx)\|^2+\langle\blmbd,\calA(\bx)\rangle
+\frac{1}{2\rho}\|\blmbd\|^2-\frac{1}{2\rho}\|\blmbd\|^2
=\Psi_\rho(\bx,\blmbd)+h(\zero).
\]
\end{proof}

For \(\theta=(\rho,s,\gamma,\eta,\tau)\) and \((\alpha,\beta)\in[0,1]^2\), consider
the recursion
\begin{equation}\label{eq:gda-recursion}
\begin{aligned}
\bx^{k+1}&=\bx^k-\eta\Bigl[\nabla_{\bx}\Psi_\rho(\bx^k,\blmbd^k)
  +\frac{1}{\gamma}(\bx^k-\bz^k)\Bigr],\\
\blmbd^{k+1}&=\blmbd^k+\tau\Bigl[\nabla_{\blmbd}\Psi_\rho(\bx^{k+1},\blmbd^k)
  -\frac{1}{s}(\blmbd^k-\bmu^k)\Bigr],\\
\bz^{k+1}&=\bz^k+\alpha(\bx^{k+1}-\bz^k),\qquad
\bmu^{k+1}=\bmu^k+\beta(\blmbd^{k+1}-\bmu^k).
\end{aligned}
\end{equation}

\begin{proposition}[Embedding]
\label{prop:embedding}
Let \(\rho>0\), and let the sequence
\(\bw^k:=(\bx^k,\blmbd^k,\bz^k,\bmu^k)\) be generated by
\(\mathrm{SPLALM}_{\theta,\alpha,\beta}\). Suppose that for some \(K\ge1\),
\((\bx^k,\blmbd^k)\in\mathcal R_\rho\) and
\((\bx^{k+1},\blmbd^k)\in\mathcal R_\rho\) for all \(0\le k<K\).
Then \eqref{eq:gda-recursion} holds for \(0\le k<K\), and
\[
R_k=\max\bigl\{\|\nabla_{\bx}\Psi_\rho(\bx^{k+1},\blmbd^k)\|,\;
\|\nabla_{\blmbd}\Psi_\rho(\bx^{k+1},\blmbd^k)\|\bigr\},
\qquad 0\le k<K.
\]
\end{proposition}

\begin{proof}
Apply Proposition~\ref{prop:reduction} at \((\bx^k,\blmbd^k)\) in the
\(\bx\)-update of~\eqref{eq:splalm-updates}, and at \((\bx^{k+1},\blmbd^k)\) in
the \(\blmbd\)-update. The same point gives \(\by^{k+1}=\zero\) and
\(\bq^k=\blmbd^k+\rho\calA(\bx^{k+1})\), so the formula for \(R_k\) follows
from~\eqref{eq:kkt-residual}.
\end{proof}

Recursion~\eqref{eq:gda-recursion} is the update of
DS-GDA~\cite{zheng2023universal} for the saddle function \(\Psi_\rho\), with
the same order of updates, whenever the projections in DS-GDA are inactive.
Proposition~\ref{prop:embedding} identifies
\(\mathrm{SPLALM}_{\theta,\alpha,\beta}\) with this recursion only while both
evaluation points \((\bx^k,\blmbd^k)\) and \((\bx^{k+1},\blmbd^k)\) lie in
\(\mathcal R_\rho\); outside \(\mathcal R_\rho\), the two recursions differ.

\subsection{Small primal center stepsize $\alpha$}
\label{subsec:lb-plateau}

In this subsection, we treat the pairs \((\alpha,\beta)\) with
\(\alpha\le\alpha_\epsilon\). We construct a hard instance on which, for
every \(\beta\in[0,1]\), \(\mathrm{SPLALM}_{\theta,\alpha,\beta}\) requires
more than \(\Delta/(16\gamma\alpha\epsilon^2)\) iterations to produce an
\(\epsilon\)-KKT output. The hard instance is given below.

\begin{example}
\label{ex:plateau}
Let \(0<\epsilon\le\min\{1/2,\sqrt{\Delta}/2\}\) and
\(x_\star:=\Delta/(2\epsilon)+\epsilon\). Define \(\calA\equiv0\),
\(h(y)=|y|\), and \(f:\RR\to\RR\) by
\[
f(x)=
\begin{cases}
2\epsilon(x_\star-x)-2\epsilon^2, & x<x_\star-2\epsilon,\\[4pt]
\dfrac12(x-x_\star)^2, & |x-x_\star|\le2\epsilon,\\[4pt]
2\epsilon(x-x_\star)-2\epsilon^2, & x>x_\star+2\epsilon.
\end{cases}
\]
Then \(f\) is \(1\)-smooth with \(\|\nabla f\|\le2\epsilon\le1\), \(\calA\) satisfies
\(L_{\calA}=G_{\calA}=0\), and \(h\) is convex and \(1\)-Lipschitz. Since
\(\epsilon\le\sqrt{\Delta}/2\), we have \(x_\star\ge2\epsilon\), thus \(f(0)=2\epsilon x_\star-2\epsilon^2=\Delta\). Moreover, \(\inf f=f(x_\star)=0\). Hence
\((f,\calA,h)\in\mathcal I(\mathbf C_0,\Delta)\).
\end{example}

\begin{remark}
\label{rem:plateau}
The construction ensures that \(\inf f=0\) and \(f(0)=\Delta\), so the
initial objective gap is exactly \(\Delta\). This allows us to make the
dependence of the iteration lower bound on \(\Delta\) explicit.
\end{remark}

This instance shows that a small primal center stepsize can force
\(\Omega\bigl(\Delta/(\gamma\alpha\epsilon^2)\bigr)\) iterations even on a
smooth one-dimensional problem, regardless of the dual center stepsize.
The iterates must traverse a long constant-gradient interval before an
output can satisfy the KKT tolerance. A small \(\alpha\) limits how quickly
the center moves, and the primal variable stays within a bounded distance
of that center. Since the multipliers remain zero, every output on this
interval has KKT residual \(2\epsilon>\epsilon\). The next lemma bounds the
distance traveled and converts it into a lower bound on the number of
iterations.

\begin{lemma}
\label{lem:plateau}
Let \((\alpha,\beta)\in[0,1]^2\), and suppose that \(\epsilon\) satisfies the
condition in Example~\ref{ex:plateau} and \(\epsilon^2\le\Delta/(16\gamma)\).
Run \(\mathrm{SPLALM}_{\theta,\alpha,\beta}\) on Example~\ref{ex:plateau}.
Then the following hold.
\begin{enumerate}[label=(\roman*)]
\item For all \(k\ge0\), \(\lambda^k=\mu^k=0\), \(y^{k+1}=0\), and
  \(R_k=\|\nabla f(x^{k+1})\|\).
\item If \(\alpha>0\), then for every integer \(k\) with \(0\le k\le\Delta/(16\gamma\alpha\epsilon^2)\),
  \[
  \|z^k-z^0\|\le2\alpha\gamma\epsilon\,k,\qquad
  \|x^k-z^k\|\le2\gamma\epsilon,
  \]
  and
  \[
  \|x^k-x^0\|\le2\gamma\epsilon(1+\alpha k)
  \le x_\star-2\epsilon .
  \]
  When \(\alpha=0\), these bounds hold for all \(k\ge0\).
\item Let \(I\) denote the instance of Example~\ref{ex:plateau}. Then
  \(T^{\mathrm{KKT}}_\epsilon(\theta,\alpha,\beta;I)>\Delta/(16\gamma\alpha\epsilon^2)\)
  if \(\alpha>0\), and \(T^{\mathrm{KKT}}_\epsilon(\theta,\alpha,\beta;I)=+\infty\)
  if \(\alpha=0\).
\end{enumerate}
\end{lemma}

\begin{proof}
(i) Since \(\calA\equiv0\) and \(\partial h(0)=[-1,1]\), every pair
\((x,0)\) lies in \(\mathcal R_\rho\), and \(\Psi_\rho(x,\lambda)=f(x)\). The
\(\lambda\)-update in~\eqref{eq:gda-recursion} reads
\(\lambda^{k+1}=\lambda^k-\frac{\tau}{s}(\lambda^k-\mu^k)\). Starting from
\(\lambda^0=\mu^0=0\), induction gives \(\lambda^k=\mu^k=0\) for all \(k\).
Proposition~\ref{prop:embedding} then gives \(y^{k+1}=0\) and
\(R_k=\max\{\|\nabla f(x^{k+1})\|,0\}\).

(ii) We use the Euclidean norm, which coincides with the absolute value
in one dimension. By (i), the primal update gives
\[
x^{k+1}-z^k
=\left(1-\frac{\eta}{\gamma}\right)(x^k-z^k)-\eta\nabla f(x^k).
\]
Since \(\|\nabla f(x)\|\le2\epsilon\) for every \(x\) and
\(0<\eta\le\gamma\),
\[
\|x^{k+1}-z^k\|
\le\left(1-\frac{\eta}{\gamma}\right)\|x^k-z^k\|+2\eta\epsilon.
\]
The center update implies
\[
\|x^{k+1}-z^{k+1}\|=(1-\alpha)\|x^{k+1}-z^k\|,
\qquad
\|z^{k+1}-z^k\|=\alpha\|x^{k+1}-z^k\|.
\]
Starting from \(x^0=z^0\), induction gives
\(\|x^k-z^k\|\le2\gamma\epsilon\) and
\(\|x^{k+1}-z^k\|\le2\gamma\epsilon\) for all \(k\ge0\).
Consequently,
\[
\|z^k-z^0\|
\le\sum_{j=0}^{k-1}\|z^{j+1}-z^j\|
\le2\alpha\gamma\epsilon k.
\]
Using \(x^0=z^0=0\), we obtain
\[
\|x^k-x^0\|
\le\|x^k-z^k\|+\|z^k-z^0\|
\le2\gamma\epsilon(1+\alpha k).
\]
If \(\alpha>0\), then for every integer \(k\le\Delta/(16\gamma\alpha\epsilon^2)\), the
assumptions \(\epsilon^2\le\Delta/(16\gamma)\) and
\(\epsilon^2\le\Delta/4\) imply
\[
\|x^k-x^0\|
\le2\gamma\epsilon+\frac{\Delta}{8\epsilon}
\le\frac{\Delta}{4\epsilon}
\le\frac{\Delta}{2\epsilon}-\epsilon
=x_\star-2\epsilon.
\]
In particular, \(x^k\le\|x^k-x^0\|\le x_\star-2\epsilon\), so the
iterate lies in the constant-gradient region. When \(\alpha=0\), the
center remains at \(z^k=z^0=0\), and the same conclusion holds for every
\(k\ge0\).

(iii) Suppose \(\alpha>0\). By (ii), every output with
\(k+1\le\Delta/(16\gamma\alpha\epsilon^2)\) satisfies
\(x^{k+1}\le x_\star-2\epsilon\). By (i), its residual is
\(R_k=\|\nabla f(x^{k+1})\|=2\epsilon>\epsilon\).
Thus, none of the first
\(\lfloor\Delta/(16\gamma\alpha\epsilon^2)\rfloor\) outputs is an
\(\epsilon\)-KKT point. Definition~\ref{def:complexity-measure} yields
\[
T^{\mathrm{KKT}}_\epsilon(\theta,\alpha,\beta;I)
\ge\left\lfloor\frac{\Delta}{16\gamma\alpha\epsilon^2}\right\rfloor+1
>\frac{\Delta}{16\gamma\alpha\epsilon^2}.
\]
When \(\alpha=0\), (ii) keeps every iterate in the constant-gradient region.
Hence \(R_k=2\epsilon>\epsilon\) for all \(k\ge0\), and
\(T^{\mathrm{KKT}}_\epsilon(\theta,\alpha,\beta;I)=+\infty\).
\end{proof}

For \(0<\alpha\le\alpha_\epsilon\), Lemma~\ref{lem:plateau}~(iii) gives
\[
T^{\mathrm{KKT}}_\epsilon(\theta,\alpha,\beta;I)
>\frac{\Delta}{16\gamma\alpha\epsilon^2}
\ge\frac{\Delta\,\epsilon^{-4}}{16384\,\gamma\max\{1,4\tau,4\gamma\tau\}}.
\]

\subsection{Large primal center stepsize \texorpdfstring{$\alpha$}{alpha}, and \texorpdfstring{$\beta>0$}{beta>0}}
\label{subsec:lb-rotation}

In this subsection, we prove the lower bound for
\(\alpha>\alpha_\epsilon\) and \(\beta>0\). For every such pair, we construct
a hard instance on which SPLALM, initialized at zero, never produces an
\(\epsilon\)-KKT point.

% On the region where the objective is quadratic, and as long as the prox
% outputs vanish, every update of SPLALM is linear. One iteration is therefore
% a fixed matrix acting on the current state, and the eigenvalues of this
% matrix determine the trajectory. If all eigenvalues had modulus less than
% one, the iterates would converge to a KKT point and the residual would
% vanish; if some eigenvalue had modulus greater than one, the iterates would
% leave the quadratic region. We therefore choose the parameters so that the
% matrix has a non-real eigenvalue of modulus one. The corresponding component
% of the trajectory rotates with constant amplitude and never settles at a
% fixed point. The initialization is fixed at zero and is not on this
% component. We show that the trajectory from zero differs from the rotation
% by a decaying term that stays small at every iteration, so the residual
% stays above \(\epsilon\). The hard instance is given below.

Motivated by the quadratic hard instances in~\cite[Example~5.1]{li2026smoothing}
and~\cite[Appendix~C.1]{mahdavinia2022tight}, we construct the following instance.

\begin{example}
\label{ex:rotation}
Let \(b,\kappa,X>0\) and \(x_c,\lambda_c\in\mathbb C\) be given in
Appendix~\ref{app:lb-rotation-params}, and set \(R_*:=\frac54X\). Define
\(\calA(x)=b(x-x_c)\), \(h(y)=\|y\|_2\), and
\[
f(x)=\psi\bigl(\|x-x_c\|\bigr)-b\langle\lambda_c,x-x_c\rangle,
\qquad
\psi(r)=
\begin{cases}
-\dfrac{\kappa}{2}r^2, & 0\le r\le R_*,\\[4pt]
\dfrac{\kappa}{2}(r-2R_*)^2-\kappa R_*^2, & R_*\le r\le3R_*,\\[4pt]
\kappa R_*(r-3R_*)-\dfrac{\kappa}{2}R_*^2, & r\ge3R_*.
\end{cases}
\]
\end{example}
To establish this lower bound, we analyze the SPLALM updates on
Example~\ref{ex:rotation}. Specifically, since \(h=\|\cdot\|_2\), Proposition~\ref{prop:reduction} ensures that both proximal outputs vanish
whenever
\[
\max_{j\in\{k,k+1\}}\|\lambda^k+\rho\calA(x^j)\|\le1.
\] 
Under this condition, SPLALM follows~\eqref{eq:gda-recursion}. In the regime where \(\|x-x_c\|\le R_*\) and \(\|\lambda+\rho\calA(x)\|\le1\),
the reduced function takes the quadratic form
\[
\Psi_\rho(x,\lambda)
=-\frac{a}{2}\|x-x_c\|^2+b\langle\lambda-\lambda_c,x-x_c\rangle .
\]
Here \(a:=\kappa-\rho b^2\). The gradients of this quadratic function are
linear in \(x-x_c\) and \(\lambda-\lambda_c\). 
Therefore, let define the iteration sequence \(\bw^k:=(x^k,\lambda^k,z^k,\mu^k)\) and
\(\bw_c:=(x_c,\lambda_c,x_c,\lambda_c)\). Whenever the proximal condition
above and \(\|x^k-x_c\|\le R_*\) hold, the four updates take the linear form
\[
\bw^{k+1}-\bw_c=\mathbb T(\bw^k-\bw_c),
\] with $\mathbb T$ defined as 
\begin{equation}
\label{eq:rot-matrix}
\mathbb T=\begin{pmatrix}
1+\eta a-\tfrac{\eta}{\gamma} & -\eta b & \tfrac{\eta}{\gamma} & 0\\[2pt]
\tau b\bigl(1+\eta a-\tfrac{\eta}{\gamma}\bigr) & 1-\tfrac{\tau}{s}-\eta\tau b^2
  & \tfrac{\eta\tau b}{\gamma} & \tfrac{\tau}{s}\\[2pt]
\alpha\bigl(1+\eta a-\tfrac{\eta}{\gamma}\bigr) & -\alpha\eta b
  & 1-\alpha+\tfrac{\alpha\eta}{\gamma} & 0\\[2pt]
\beta\tau b\bigl(1+\eta a-\tfrac{\eta}{\gamma}\bigr)
  & \beta\bigl(1-\tfrac{\tau}{s}-\eta\tau b^2\bigr)
  & \tfrac{\beta\eta\tau b}{\gamma} & 1-\beta+\tfrac{\beta\tau}{s}
\end{pmatrix}.
\end{equation}
Under the conditions above, analyzing the SPLALM iterates reduces to
studying the spectral properties of \(\mathbb T\).
The core idea is to construct a bounded trajectory whose KKT residual remains
above \(\epsilon\). We choose \(a\) and \(b\) so that \(\mathbb T\) has a
non-real eigenvalue of modulus one, which gives a component that does not
decay.

The construction in~\cite[Section~5.1.1]{li2026smoothing} chooses both the
hard instance and the initial point, with the latter lying along an
eigenvector of the iteration matrix. In our setting, the initial point is
fixed at zero, and we choose the instance alone. We therefore need to
obtain the required trajectory from this prescribed initial point. To this
end, we choose the centers \(x_c\) and \(\lambda_c\) in the instance so that
the shifted initial state is the sum of an eigenvector of \(\mathbb T\)
and a decaying component. The following lemma gives the spectral
properties and this decomposition for the parameters defined
in Appendix~\ref{app:lb-rotation-params}.

\begin{lemma}
\label{lem:rot-spectral}
Under the conditions of Lemma~\ref{lem:rot-estimates}, let
\(\mathbb T\) be the matrix in~\eqref{eq:rot-matrix}. Then:
\begin{enumerate}[label=(\roman*)]
\item \(\mathbb T\bv=\zeta\bv\), where \(|\zeta|=1\) and \(\zeta\notin\RR\);
\item \(\br\) lies in a \(\mathbb T\)-invariant subspace on which the spectral
  radius of \(\mathbb T\) is less than one;
\item \(\bv+\br=-\bw_c\).
\end{enumerate}
\end{lemma}
 
By Lemma~\ref{lem:rot-spectral}, the linear recursion initialized at zero
has the explicit solution
\[
\bw^k-\bw_c=\zeta^k\bv+\mathbb T^k\br.
\]
We next bound \(\mathbb T^k\br\) uniformly in \(k\) to verify that this
sequence stays in the quadratic region, both proximal outputs vanish,
and the KKT residual remains above \(\epsilon\).

\begin{lemma}
\label{lem:rot-estimates}
Let \(\alpha>\alpha_\epsilon\), \(0<\beta\le1\), \(0<\epsilon<c_2\), and
\(\epsilon(\rho+s)\le c_1\), with \(c_1,c_2\) as specified in
Theorem~\ref{thm:lower-bound}. Then:
\begin{enumerate}[label=(\roman*)]
\item the instance of Example~\ref{ex:rotation} belongs to
  \(\mathcal I(\mathbf C_0,\Delta)\);
\item for all \(k\ge0\),
\[
\|(\mathbb T^k\br)_x\|\le\frac X4,
\qquad
\|(\mathbb T^k\br)_\lambda\|\le
s\Bigl(\frac{\gamma}{\eta}+\frac54\Bigr)bX.
\]
\end{enumerate}
\end{lemma}
The proofs of Lemmas~\ref{lem:rot-spectral} and~\ref{lem:rot-estimates} are
given in Appendix~\ref{app:lb-rotation}.

Using these estimates, we now show that the sequence generated by the linear
recursion coincides with the SPLALM iterates and that its KKT residual
remains above \(\epsilon\) at every iteration.

\begin{lemma}
\label{lem:rot-trajectory}
Under the conditions of Lemma~\ref{lem:rot-estimates}, the iterates of
\(\mathrm{SPLALM}_{\theta,\alpha,\beta}\) on Example~\ref{ex:rotation}
satisfy, for all \(k\ge0\),
\[
x^k-x_c=\zeta^kX+(\mathbb T^k\br)_x,
\qquad
\lambda^k-\lambda_c=\zeta^k\Lambda+(\mathbb T^k\br)_\lambda,
\]
and \(R_k\ge\frac76\epsilon\). Hence
\(T^{\mathrm{KKT}}_\epsilon(\theta,\alpha,\beta;I)=+\infty\).
\end{lemma}

\begin{proof}
We prove by induction that, for all \(k\ge0\),
\[
\bw^k-\bw_c=\zeta^k\bv+\mathbb T^k\br.
\]
\emph{Base case \(k=0\).} By Lemma~\ref{lem:rot-spectral}~(iii),
\(\bv+\br=-\bw_c\). Since \(\bw^0=\zero\), the identity holds for \(k=0\).

\emph{Inductive step.} Assume that the identity holds for some \(k\ge0\).
Since \(|\zeta|=1\), Lemma~\ref{lem:rot-estimates}~(ii) gives
\[
\frac34X\le\|x^k-x_c\|\le\frac54X=R_*.
\]
Thus, \(x^k\) lies in the quadratic region. Moreover, the induction
hypothesis and Lemma~\ref{lem:rot-spectral}~(iii) imply
\(\|\lambda^k-\lambda_c\|\le\|\Lambda\|+\|(\mathbb T^k\br)_\lambda\|\)
and \(\|\lambda_c\|\le\|\Lambda\|+\|\br_\lambda\|\).
Lemma~\ref{lem:rot-estimates}~(ii) and the parameter bounds in
Lemma~\ref{lem:rot-parameter-bounds} therefore yield
\[
\|\lambda^k+\rho\calA(x^k)\|
\le2\|\Lambda\|+2s\Bigl(\frac{\gamma}{\eta}+\frac54\Bigr)bX+\frac54\rho bX
\le\frac13\Bigl(2+\frac14\Bigr)=\frac34.
\]
Hence \((x^k,\lambda^k)\in\mathcal R_\rho\). By
Proposition~\ref{prop:reduction}, the primal update is given by the first
row of~\eqref{eq:rot-matrix}. Using the induction hypothesis and
\(\mathbb T\bv=\zeta\bv\), we obtain
\[
x^{k+1}-x_c=\zeta^{k+1}X+(\mathbb T^{k+1}\br)_x.
\]
Lemma~\ref{lem:rot-estimates}~(ii) then gives
\(\frac34X\le\|x^{k+1}-x_c\|\le\frac54X\).
The same estimate as above, with \(x^{k+1}\) in place of \(x^k\), yields
\[
\|\lambda^k+\rho\calA(x^{k+1})\|\le\frac34.
\]
Thus, \((x^{k+1},\lambda^k)\in\mathcal R_\rho\) as well, so
Proposition~\ref{prop:reduction} also applies to the dual update.
Together with the center updates, this gives
\[
\bw^{k+1}-\bw_c
=\mathbb T(\bw^k-\bw_c)
=\zeta^{k+1}\bv+\mathbb T^{k+1}\br.
\]
This completes the induction. In particular, both proximal outputs vanish
at every iteration, and all primal iterates lie in the prescribed annulus.

Finally, we bound the residual. By Proposition~\ref{prop:embedding},
\[
R_k=\max\bigl\{\|{-a}(x^{k+1}-x_c)+b(\lambda^k-\lambda_c)\|,\
b\|x^{k+1}-x_c\|\bigr\}.
\]
The second term is at least \(\frac34bX\). For the first term,
Lemma~\ref{lem:rot-estimates}~(ii) and Lemma~\ref{lem:rot-parameter-bounds} give
\[
\|{-a}(x^{k+1}-x_c)+b(\lambda^k-\lambda_c)\|
\ge b\|\Lambda\|-s\Bigl(\frac{\gamma}{\eta}+\frac54\Bigr)b^2X-\frac54aX
\ge\frac{27}{32}b\|\Lambda\|-s\Bigl(\frac{\gamma}{\eta}+\frac52\Bigr)b^2X .
\]
If \(\|\Lambda\|\ge X\), Lemma~\ref{lem:rot-parameter-bounds} bounds the
right-hand side below by \(\frac{23}{32}b\|\Lambda\|\). If
\(X\ge\|\Lambda\|\), we use the second term instead. In both cases,
\(R_k\ge\frac b2\max\{X,\|\Lambda\|\}=\frac b6\ge\frac76\epsilon\).
\end{proof}
 


\subsection{Large primal center stepsize \texorpdfstring{$\alpha$}{alpha}, and \texorpdfstring{$\beta=0$}{beta=0}}
\label{subsec:lb-huber}

In this subsection, we prove the lower bound for \(\alpha>\alpha_\epsilon\)
and \(\beta=0\). The instance below keeps the objective of
Example~\ref{ex:plateau} and changes the inner mapping and the outer
function. On this instance, \(\mathrm{SPLALM}_{\theta,\alpha,0}\), initialized
at zero, never produces an \(\epsilon\)-KKT point, for every
\(\alpha\in[0,1]\).

\begin{example}
\label{ex:huber}
Let \(\epsilon\) satisfy the condition in Example~\ref{ex:plateau}, and let
\(f\) be the function defined there. Set
\[
q_h:=2s+1,
\qquad
a_0:=\frac{4\epsilon(\rho+3s+1)}{q_h}.
\]
Define \(\calA\equiv a_0\) and \(h:\RR\to\RR\) by
\[
h(y)=
\begin{cases}
\dfrac{q_h}{2}y^2, & |y|\le q_h^{-1},\\[4pt]
|y|-\dfrac{1}{2q_h}, & |y|>q_h^{-1}.
\end{cases}
\]
Then \(h\) is convex and \(1\)-Lipschitz, and \(\calA\) satisfies
\(L_{\calA}=G_{\calA}=0\). Since \(F=f+h(a_0)\) differs from \(f\) by a
constant, Remark~\ref{rem:plateau} gives
\(F(0)-\inf F=f(0)-\inf f=\Delta\). Together with the properties of \(f\)
established in Example~\ref{ex:plateau}, this gives
\((f,\calA,h)\in\mathcal I(\mathbf C_0,\Delta)\).
\end{example}

When \(\beta=0\), the dual center stays at \(\mu^k=\mu^0=0\). The dual
perturbation term \(-\frac{1}{2s}\|\lambda\|^2\) in~\eqref{eq:potential}
therefore pulls \(\lambda^k\) toward zero at every iteration. On this
instance, the output \(y^{k+1}\) equals \(a_0\) only when
\(\lambda^k=\nabla h(a_0)=q_ha_0\), whereas the dual iterates converge to a
smaller value. The difference keeps the feasibility residual
\(\|\calA(x^{k+1})-y^{k+1}\|\) above \(\epsilon\) at every iteration. Since
\(J_{\calA}=0\), the multiplier does not enter the primal update, and the
dual iterates have the closed form given in the next lemma.

\begin{lemma}
\label{lem:huber}
Let \(\alpha\in[0,1]\) and \(\beta=0\). Suppose that \(\epsilon\) satisfies
the condition in Example~\ref{ex:plateau} and \(\epsilon(\rho+s)\le c_1\),
with \(c_1\) as specified in Theorem~\ref{thm:lower-bound}. Run
\(\mathrm{SPLALM}_{\theta,\alpha,0}\) on Example~\ref{ex:huber}, and set
\[
\varphi:=1-\frac{\tau}{s}-\frac{\tau}{\rho+q_h},
\qquad
\bar\lambda:=\frac{s\,q_ha_0}{\rho+q_h+s}.
\]
Then the following hold.
\begin{enumerate}[label=(\roman*)]
\item For all \(k\ge0\),
  \[
  \mu^k=0,\qquad
  \lambda^k=\bar\lambda\bigl(1-\varphi^k\bigr),\qquad
  y^{k+1}=\frac{\rho a_0+\lambda^k}{\rho+q_h}.
  \]
\item For all \(k\ge0\),
  \(R_k\ge\|\calA(x^{k+1})-y^{k+1}\|\ge3\epsilon\). Hence, with \(I\) the
  instance of Example~\ref{ex:huber},
  \(T^{\mathrm{KKT}}_\epsilon(\theta,\alpha,0;I)=+\infty\).
\end{enumerate}
\end{lemma}

\begin{proof}
We first collect three bounds on the parameters. Since \(0<\tau\le s\) and
\(\rho+q_h>2s\),
\begin{equation}
\label{eq:huber-ratio}
-\frac12<-\frac{\tau}{\rho+q_h}\le\varphi<1,
\qquad
\frac{s}{\rho+q_h}<\frac12 .
\end{equation}
Hence \(\varphi^k\ge-\frac12\) and \(0\le1-\varphi^k\le2\) for all
\(k\ge0\). Since \(2s\le\rho+q_h+s\), it follows that
\(0\le\bar\lambda(1-\varphi^k)\le2\bar\lambda\le q_ha_0\). Finally,
\(s\ge\tau\) gives \(1\le(\rho+s)/\tau\), so the third term in the
definition of \(c_1\) yields
\[
q_ha_0=4\epsilon(\rho+3s+1)
\le4\epsilon\bigl(3+\tau^{-1}\bigr)(\rho+s)
\le4c_1\bigl(3+\tau^{-1}\bigr)
\le\frac14 .
\]

(i) Since \(\beta=0\) and \(\mu^0=0\), the dual center update gives
\(\mu^k=0\) for all \(k\). Next, fix \(x\in\RR\) and
\(\lambda\in[0,q_ha_0]\), and set
\(y_\lambda:=(\rho a_0+\lambda)/(\rho+q_h)\). Then
\(0\le y_\lambda\le a_0\le q_h^{-1}\), where the last inequality uses
\(q_ha_0\le\frac14\). On \([-q_h^{-1},q_h^{-1}]\), the function \(h\) is
differentiable with \(\nabla h(y)=q_hy\), and \(y_\lambda\) satisfies
\[
q_hy_\lambda-\lambda-\rho\bigl(a_0-y_\lambda\bigr)=0,
\]
which is the optimality condition of the strongly convex problem
in~\eqref{eq:reduced-al}. Hence \(y(x,\lambda)=y_\lambda\), and
\[
\calA(x)-y(x,\lambda)=\frac{q_ha_0-\lambda}{\rho+q_h}.
\]
By Danskin's theorem,
\(\nabla_\lambda\calD_\rho(x,\lambda)=\calA(x)-y(x,\lambda)\).

We now prove the formula for \(\lambda^k\) by induction. It holds for
\(k=0\) because \(\lambda^0=0\). Assume that it holds for some \(k\ge0\).
Then \(\lambda^k\in[0,q_ha_0]\), so
\(y^{k+1}=y(x^{k+1},\lambda^k)=y_{\lambda^k}\). Together with
\(\mu^k=0\), the dual update in~\eqref{eq:splalm-updates} reads
\[
\lambda^{k+1}
=\lambda^k+\tau\Bigl[\frac{q_ha_0-\lambda^k}{\rho+q_h}
-\frac{\lambda^k}{s}\Bigr]
=\varphi\lambda^k+\frac{\tau q_ha_0}{\rho+q_h}.
\]
A direct computation shows that
\(\bar\lambda=\varphi\bar\lambda+\tau q_ha_0/(\rho+q_h)\). Subtracting the
two identities gives
\(\lambda^{k+1}-\bar\lambda=\varphi(\lambda^k-\bar\lambda)
=-\bar\lambda\varphi^{k+1}\), which completes the induction.

(ii) Since \(J_{\calA}=0\), \eqref{eq:kkt-residual} gives
\(R_k=\max\{\|\nabla f(x^{k+1})\|,\|\calA(x^{k+1})-y^{k+1}\|\}\). By (i)
and \(q_ha_0-\bar\lambda=q_ha_0(\rho+q_h)/(\rho+q_h+s)\),
\[
\calA(x^{k+1})-y^{k+1}
=\frac{q_ha_0-\lambda^k}{\rho+q_h}
=\frac{q_ha_0}{\rho+q_h+s}\Bigl(1+\frac{s}{\rho+q_h}\varphi^k\Bigr)
=4\epsilon\Bigl(1+\frac{s}{\rho+q_h}\varphi^k\Bigr),
\]
where the last equality uses \(\rho+q_h+s=\rho+3s+1\).
By~\eqref{eq:huber-ratio}, \(\frac{s}{\rho+q_h}\varphi^k>-\frac14\), so
\(\|\calA(x^{k+1})-y^{k+1}\|>3\epsilon\). Thus no output is an
\(\epsilon\)-KKT point, and
\(T^{\mathrm{KKT}}_\epsilon(\theta,\alpha,0;I)=+\infty\).
\end{proof}

\subsection{Proof of Theorem~\ref{thm:lower-bound}}
\label{subsec:lb-combine}

We now combine the three instances. Let \(c_1\) and \(c_2\) be given
by~\eqref{eq:rot-thresholds}. The constant \(c_1\) depends only on
\((\gamma,\eta,\tau)\), and \(c_2\) depends only on \((\gamma,\tau,\Delta)\),
as required in Theorem~\ref{thm:lower-bound}.

\begin{proof}[Proof of Theorem~\ref{thm:lower-bound}]
Fix \(\epsilon\in(0,c_2)\), \(\rho>0\), and \(s\ge\tau\) with
\(\epsilon(\rho+s)\le c_1\), and fix \((\alpha,\beta)\in[0,1]^2\). The
definition of \(c_2\) gives \(\epsilon\le\frac1{64}\le\frac12\),
\(\epsilon\le\sqrt\Delta/2\), and
\(\epsilon^2\le\Delta/(64\gamma)\le\Delta/(16\gamma)\). Hence the conditions
of Example~\ref{ex:plateau}, Lemma~\ref{lem:plateau}, and
Example~\ref{ex:huber} on \(\epsilon\) hold. We treat the three parts of the
partition of \([0,1]^2\) separately.

\emph{Case 1: \(\alpha\le\alpha_\epsilon\).} Let \(I\) be the instance of
Example~\ref{ex:plateau}, which belongs to \(\mathcal I(\mathbf C_0,\Delta)\).
If \(\alpha=0\), Lemma~\ref{lem:plateau}~(iii) gives
\(T^{\mathrm{KKT}}_\epsilon(\theta,\alpha,\beta;I)=+\infty\). If
\(0<\alpha\le\alpha_\epsilon\), Lemma~\ref{lem:plateau}~(iii) and the
definition of \(\alpha_\epsilon\) give
\[
T^{\mathrm{KKT}}_\epsilon(\theta,\alpha,\beta;I)
>\frac{\Delta}{16\gamma\alpha\epsilon^2}
\ge\frac{\Delta}{16\gamma\alpha_\epsilon\epsilon^2}
=\frac{\Delta\,\epsilon^{-4}}{16384\,\gamma\max\{1,4\tau,4\gamma\tau\}}.
\]
Both bounds hold for every \(\beta\in[0,1]\).

\emph{Case 2: \(\alpha>\alpha_\epsilon\) and \(\beta>0\).} The conditions of
Lemma~\ref{lem:rot-estimates} hold. Let \(I\) be the instance of
Example~\ref{ex:rotation}. By Lemma~\ref{lem:rot-estimates}~(i),
\(I\in\mathcal I(\mathbf C_0,\Delta)\), and Lemma~\ref{lem:rot-trajectory}
gives \(T^{\mathrm{KKT}}_\epsilon(\theta,\alpha,\beta;I)=+\infty\).

\emph{Case 3: \(\alpha>\alpha_\epsilon\) and \(\beta=0\).} The conditions of
Lemma~\ref{lem:huber} hold. Let \(I\) be the instance of
Example~\ref{ex:huber}, which belongs to \(\mathcal I(\mathbf C_0,\Delta)\).
Lemma~\ref{lem:huber}~(ii) gives
\(T^{\mathrm{KKT}}_\epsilon(\theta,\alpha,0;I)=+\infty\).

The three cases cover \([0,1]^2\). Hence, for every
\((\alpha,\beta)\in[0,1]^2\), some \(I\in\mathcal I(\mathbf C_0,\Delta)\)
satisfies
\[
T^{\mathrm{KKT}}_\epsilon(\theta,\alpha,\beta;I)
\ge\frac{\Delta\,\epsilon^{-4}}{16384\,\gamma\max\{1,4\tau,4\gamma\tau\}}.
\]
Taking the supremum over \(I\) and then the infimum over \((\alpha,\beta)\)
in Definition~\ref{def:complexity-measure} gives the bound on
\(T^{\mathrm{KKT}}_\epsilon(\calC_\theta,\mathcal I(\mathbf C_0,\Delta))\).
\end{proof}

In Cases~2 and~3, the method never produces an \(\epsilon\)-KKT output.
The finite bound comes only from Case~1, and it determines the order
\(\epsilon^{-4}\) in Theorem~\ref{thm:lower-bound}.

% !TEX root = ../main.tex
\section{Conclusion}
\label{sec:conclusion}

In this paper, we studied the iteration complexity of single-loop linearized
augmented Lagrangian methods with a fixed penalty. For the family
\(\calC_\theta\), we established a lower bound of \(\Omega(\epsilon^{-4})\)
that holds for any choice of constant center stepsizes whenever
\(\epsilon(\rho+s)\le c_1\), and we showed that the choice \(\beta=1\)
attains this order. With \(\rho\) fixed and \(s=c_s/\epsilon\) for a
sufficiently large constant \(c_s\), the choice \(\beta=0\) and a suitable
primal center stepsize \(\alpha=\Theta(\epsilon)\) attains
\(\mathcal O(\epsilon^{-3})\).

These results delineate the roles of the penalty, dual regularization, and
center updates within \(\calC_\theta\). The combined penalty \(\rho+s\)
connects penalty selection in ALMs with dual regularization in the reduced
minimax formulation. Adjusting the constant center stepsizes cannot
substitute for the growth of \(\rho+s\) required to attain
\(\mathcal O(\epsilon^{-3})\). This growth can be realized by reducing the
dual regularization coefficient \(1/s\) as the target accuracy decreases,
while keeping the penalty \(\rho\) fixed.

An interesting direction for future research is to extend this analysis to
smoothed gradient descent--ascent methods for general nonconvex--concave
minimax problems.


\bibliographystyle{tools/siamplain}
\bibliography{tools/refer}

\appendix
% !TEX root = ../main.tex
\section{Details for the rotating instance}
\label{app:lb-rotation}

\subsection{Parameters}
\label{app:lb-rotation-params}
Let \(V:=\beta+\frac{\tau}{s}(1-\beta)\) and
\begin{equation}
\nu:=\frac{64\eta\tau\alpha\beta\epsilon^2}
  {\bigl(\alpha+\frac{\eta}{\gamma}(1-\alpha)\bigr)V},
\qquad
\zeta:=1-\frac{\nu}{2}+i\sqrt{\nu-\frac{\nu^2}{4}} .
\label{eq:rot-basic}
\end{equation}
Under the conditions of Lemma~\ref{lem:rot-estimates}, \(0<\nu<4\); hence
\(\zeta\notin\RR\), \(|\zeta|=1\), and \(|\zeta-1|^2=\nu\). Let
\(W:=1+\frac{\eta}{\gamma}\alpha(1-\alpha)/|\zeta-1+\alpha|^2\), and define
\begin{equation}
b^2:=\frac{\nu W|\zeta-1+V|^2}{\eta\tau\bigl(\beta V+(1-V)\nu\bigr)},
\qquad
a:=\frac{\nu(W-1)}{\eta\alpha}+\frac{\tau b^2(V-\beta)}{|\zeta-1+V|^2},
\qquad
\kappa:=\rho b^2+a.
\label{eq:rot-b-a}
\end{equation}
Let \(X>0\) be such that \(\max\{X,\|\Lambda\|\}=\frac13\), where
\(\Lambda:=\frac{\tau b\zeta(\zeta-1+\beta)}{(\zeta-1)(\zeta-1+V)}X\), and set
\begin{equation}
\bv:=\Bigl(X,\ \Lambda,\ \frac{\alpha\zeta}{\zeta-1+\alpha}X,\
  \frac{\beta\zeta}{\zeta-1+\beta}\Lambda\Bigr).
\label{eq:rot-eigvec}
\end{equation}
The shifts \(x_c\) and \(\lambda_c\) are given in~\eqref{eq:rot-shift}. They
are chosen so that \(-\bw_c-\bv\) lies in the \(\mathbb T\)-invariant subspace
on which the spectral radius of \(\mathbb T\) is less than one.


We use the following choices of the constants in Theorem~\ref{thm:lower-bound}:
\begin{equation}
\label{eq:rot-thresholds}
\begin{aligned}
c_1&:=\min\left\{
\frac1{32}\sqrt{\frac{\tau}{\eta(1+4\gamma/\eta)}},\,
\frac1{64(\gamma/\eta+5/2)},\,
\frac1{16(3+\tau^{-1})}\right\},\\
c_2&:=\min\left\{\frac1{64},\,
\frac1{32\sqrt{\max\{1,4\tau,4\gamma\tau\}}},\,
\frac\Delta{12},\,\frac{\sqrt\Delta}{2},\,
\sqrt{\frac\Delta{64\gamma}}\right\}.
\end{aligned}
\end{equation}
Throughout this appendix, assume the parameter conditions of
Lemma~\ref{lem:rot-estimates}. All complex scalars represent vectors in
\(\RR^2\), with the Euclidean norm equal to the complex modulus.
For the calculations below, write
\[
\chi:=\frac\eta\gamma,\qquad d_\tau:=\frac\tau s,\qquad
U:=\alpha+\chi(1-\alpha),\qquad D:=\alpha W,
\]
\[
E:=|\zeta-1+\alpha|^2=\alpha^2+(1-\alpha)\nu,\qquad
H:=|\zeta-1+V|^2=V^2+(1-V)\nu,
\]
\[
\Upsilon:=\beta V+(1-V)\nu,\qquad
\xi:=\frac{\eta\tau b^2(V-\beta)}H,\qquad
C_{\rm tr}:=\frac1{d_\tau}+\frac{2(1+d_\tau^{-1})}{\chi}.
\]
These quantities are used only to prove the estimates and spectral identities.

\subsection{Parameter bounds and proof of Lemma~\ref{lem:rot-estimates}~(i)}
\label{app:lb-rotation-params-proof}
\begin{lemma}
\label{lem:rot-parameter-bounds}
Under the conditions of Lemma~\ref{lem:rot-estimates}, the parameters satisfy
\[
7\epsilon\le b\le8\epsilon,
\qquad \max\{X,\|\Lambda\|\}=\frac13,
\]
\[
aX\le\frac18b\|\Lambda\|+sb^2X,
\]
\[
bs\Bigl(\frac{\gamma}{\eta}+\frac52\Bigr)\le\frac18,
\qquad
b\Bigl(2s\Bigl(\frac{\gamma}{\eta}+\frac54\Bigr)+\frac54\rho\Bigr)\le\frac14.
\]
\end{lemma}

\begin{proof}[Proof of Lemma~\ref{lem:rot-parameter-bounds}]
We first establish the bounds needed in all three proofs. Since
\(U\ge\chi\), \(V\ge\beta\), and
\(\alpha>1024\max\{1,4\tau,4\gamma\tau\}\epsilon^2\),
\[
0<\frac\nu{\alpha^2}
\le\frac{64\gamma\tau\epsilon^2}{\alpha}<\frac1{64}.
\]
In particular, \(0<\nu<4\), so \(\zeta\) is well defined and
\(|\zeta|=1\). The definition of \(W\) gives
\[
\frac{64}{65}U\le
D=\alpha+\frac{\chi(1-\alpha)}{1+(1-\alpha)\nu/\alpha^2}
\le U\le1,
\qquad D\ge\frac\chi2.
\]
Since \(V\Upsilon-\beta H=(1-V)(V-\beta)\nu\ge0\),
\[
\frac{b^2}{64\epsilon^2}=\frac D U\frac{\beta H}{V\Upsilon}\le1.
\]
Also, \(H\ge V^2\), \(V\ge d_\tau\), and \(V-\beta\le V\) imply
\(0\le\xi\le\eta s b^2\). Using
\(C_{\rm tr}\le(1+4/\chi)/d_\tau\) and~\eqref{eq:rot-thresholds}, we obtain
\begin{equation}
\label{eq:rot-small-feedback}
\xi C_{\rm tr}
\le\frac{64\eta(1+4/\chi)}\tau(\epsilon s)^2
\le\frac1{16}.
\end{equation}
Thus \(\xi/d_\tau\le1/16\). The identity
\[
\frac{\beta H}{V\Upsilon}
=1-\frac{(1-V)\xi}{WV}\ge\frac{15}{16}
\]
now yields \(b^2/(64\epsilon^2)\ge12/13>49/64\), proving
\(7\epsilon\le b\le8\epsilon\).
The normalization of \(X\) gives \(\max\{X,\|\Lambda\|\}=1/3\).

To bound \(a\), set
\(a_{\rm p}:=\nu(W-1)/(\eta\alpha)\) and \(a_{\rm d}:=\xi/\eta\).
The definitions imply
\begin{equation}
\label{eq:rot-algebra-identities}
\eta a=\frac{\nu(W-1)}\alpha+\xi,
\qquad \eta\tau b^2=\nu W+V\xi,
\qquad U-D=\frac{(1-\alpha)\nu(W-1)}\alpha.
\end{equation}
Hence \(a=a_{\rm p}+a_{\rm d}\),
\(a_{\rm p}\le\tau b^2/\alpha\), and \(a_{\rm d}\le sb^2\).
Since \(c_1\le1/224\) and \(\epsilon<1/64\),
\[
0<\kappa\le(\rho+s)b^2+\frac{\tau b^2}\alpha
\le64c_1\epsilon+\frac{\tau}{16\max\{1,4\tau,4\gamma\tau\}}
\le\frac1{32}.
\]
For the amplitude bounds, put \(E_\beta:=\beta^2+(1-\beta)\nu\).
The formula for \(\Lambda\) gives
\(\|\Lambda\|^2/X^2=\tau WE_\beta/(\eta\Upsilon)\).
Moreover,
\[
\Upsilon\le
\operatorname{Re}\bigl((\zeta-1+V)\overline{(\zeta-1+\beta)}\bigr)
\le\sqrt{HE_\beta},
\]
where the difference in the first inequality is \((V-\beta)\nu/2\).
Consequently,
\[
\frac{a_{\rm p}X}{b\|\Lambda\|}
=\frac{\sqrt\nu(W-1)\Upsilon}{\alpha W\sqrt{HE_\beta}}
\le\frac{\sqrt\nu}\alpha\le\frac18,
\]
\[
\frac{a_{\rm p}X^2}{b^2\|\Lambda\|^2}
=\frac{\eta(W-1)\Upsilon^2}{\alpha W^2HE_\beta}
\le\frac\eta D\le2\gamma.
\]
Adding the bound on \(a_{\rm d}\) proves the asserted bound on \(aX\)
and also gives
\begin{equation}
\label{eq:rot-gap-curvature}
\kappa X^2\le(\rho+s+2\gamma)b^2\max\{X^2,\|\Lambda\|^2\}.
\end{equation}
Finally, \(b\le8\epsilon\) and~\eqref{eq:rot-thresholds} imply
\[
bs\Bigl(\frac1\chi+\frac52\Bigr)
\le8c_1\Bigl(\frac1\chi+\frac52\Bigr)\le\frac18,
\]
\[
b\Bigl(2s\Bigl(\frac1\chi+\frac54\Bigr)+\frac54\rho\Bigr)
\le16c_1\Bigl(\frac1\chi+\frac54\Bigr)
\le\frac{1/\chi+5/4}{4(1/\chi+5/2)}<\frac14.
\]
\end{proof}

\begin{proof}[Proof of Lemma~\ref{lem:rot-estimates}~(i)]
We use the bounds on \(\mathbb T^k\br\) proved in
Appendix~\ref{app:lb-rotation-transient-proof}; that proof uses only
Lemma~\ref{lem:rot-parameter-bounds} and the spectral identities, and does
not use membership in the problem class. At \(k=0\), those bounds and
\eqref{eq:rot-shift} give
\[
\|x_c\|\le\frac54X=R_*,\qquad
\|\lambda_c\|\le\|\Lambda\|+s\Bigl(\frac1\chi+\frac54\Bigr)bX
\le\frac13\left(1+\frac{1/\chi+5/4}{8(1/\chi+5/2)}\right)<\frac12.
\]
The function \(\psi\) and its first derivative agree at \(R_*\) and
\(3R_*\). On each smooth radial piece, the Hessian of
\(\psi(\|x-x_c\|)\) has radial eigenvalue \(\psi''(r)\) and tangential
eigenvalue \(\psi'(r)/r\), both of absolute value at most \(\kappa\).
Near \(x_c\), this function is quadratic. Integrating its continuous
gradient along line segments therefore shows that
\(\operatorname{Lip}(\nabla f)\le\kappa\le1/32\).
Also \(|\psi'(r)|\le\kappa R_*\) for every \(r\ge0\), so
\[
\|\nabla f(x)\|\le\kappa R_*+b\|\lambda_c\|
\le\frac5{384}+\frac12<1.
\]
Here we used \(R_*\le5/12\) and \(b\le8\epsilon<1\).
The affine mapping \(\calA\) has constant Jacobian of norm \(b\le1\),
and \(h=\|\cdot\|_2\) is convex and 1-Lipschitz.

It remains to bound the initial gap. Since \(\psi(r)\ge-\kappa R_*^2\)
and \(\|\lambda_c\|<1/2\),
\[
F(x)=\psi(\|x-x_c\|)-b\langle\lambda_c,x-x_c\rangle+b\|x-x_c\|
\ge-\kappa R_*^2.
\]
At the origin, \(\|x_c\|\le R_*\) and \(\psi(\|x_c\|)\le0\), whence
\[
F(0)-\inf F\le\frac32bR_*+\kappa R_*^2
\le5\epsilon+\frac{100}{9}(\rho+s+2\gamma)\epsilon^2.
\]
By \(\epsilon(\rho+s)\le c_1\le1/224\), this is at most
\(6\epsilon+32\gamma\epsilon^2\le\Delta\), where the last inequality
uses \(\epsilon<c_2\). Thus the instance belongs to
\(\mathcal I(\mathbf C_0,\Delta)\).
\end{proof}

\subsection{Proof of Lemma~\ref{lem:rot-spectral}}
\label{app:lb-rotation-spectral-proof}
Define
\[
\begin{aligned}
r_x&:=-\frac{(W-1)\frac{\zeta-1}{\zeta-1+\alpha}+\frac{\xi\zeta}{\zeta-1+V}}
  {W+\xi}\,X,\\
r_\lambda&:=\frac{1}{\eta b}\Bigl[\frac{\eta}{\gamma}
  \frac{(1-\alpha)(\zeta-1)}{\zeta-1+\alpha}X+(\eta a+\alpha W)r_x
  -\xi\Bigl((1-\alpha)r_x+\frac{(V-\alpha)\zeta}{\zeta-1+V}X\Bigr)\Bigr].
\end{aligned}
\]
Set
\begin{equation}
x_c:=-(X+r_x),\qquad \lambda_c:=-(\Lambda+r_\lambda),\qquad
\br:=-\bw_c-\bv .
\label{eq:rot-shift}
\end{equation}
\begin{proof}
For this calculation, write \(Z:=\bv_z\), \(M:=\bv_\mu\), and
\(\vartheta:=(1-V)\xi\). The characteristic polynomial factors as
\begin{equation}
\label{eq:rot-characteristic-factorization}
\begin{aligned}
\det((1+t)I-\mathbb T)
&=t^2(t+U)(t+V)-\eta a\,t(t+\alpha)(t+V)\\
&\quad+\eta\tau b^2(t+1)(t+\alpha)(t+\beta)\\
&=(t^2+\nu t+\nu)
  \bigl((t+D)(t+V)-\vartheta(t+\alpha)\bigr).
\end{aligned}
\end{equation}
The first equality follows by eliminating the two center coordinates in
\eqref{eq:rot-matrix}; expansion using~\eqref{eq:rot-algebra-identities}
gives the second equality. Since
\((\zeta-1)^2+\nu(\zeta-1)+\nu=0\), substitution of
\eqref{eq:rot-eigvec} in the four rows of \(\mathbb T\) gives
\[
\mathbb T\bv=\zeta\bv.
\]
For example, the center rows follow from
\((\zeta-1+\alpha)Z=\alpha\zeta X\) and
\((\zeta-1+\beta)M=\beta\zeta\Lambda\); the dual row reduces to
\((\zeta-1)(\zeta-1+V)\Lambda
=\tau b\zeta(\zeta-1+\beta)X\).
The primal row then follows from the first factor in
\eqref{eq:rot-characteristic-factorization}. The bounds on \(\nu\)
proved above give \(\zeta\notin\RR\) and \(|\zeta|=1\).

We next prove
\begin{equation}
\label{eq:rot-annihilator}
Q_{\rm f}(\mathbb T)\br=0,
\qquad Q_{\rm f}(t):=(t-1+D)(t-1+V)-\vartheta(t-1+\alpha).
\end{equation}
First consider distinct roots \(t_1,t_2\) of \(Q_{\rm f}\), with
\(\alpha<1\), \(\beta<1\), \(d_\tau<1\), \(V\ne\alpha\), and all
following denominators nonzero. An eigenvector associated with a root \(t\) is
\[
\left(1,\ \frac{\tau bt(t-1+\beta)}{(t-1)(t-1+V)},\
\frac{\alpha t}{t-1+\alpha},\
\frac{\tau b\beta t^2}{(t-1)(t-1+V)}\right)^\top.
\]
Let its coefficients at \(t_1,t_2\) be \(c_1',c_2'\).
Requiring that their sum with \(\bv\) have equal first and third
coordinates, and equal second and fourth coordinates, gives
\[
\sum_{j=1}^2c_j'\frac{t_j-1}{t_j-1+\alpha}
=-\frac{\zeta-1}{\zeta-1+\alpha}X,
\qquad
\sum_{j=1}^2c_j'\frac{t_j}{t_j-1+V}
=-\frac\zeta{\zeta-1+V}X.
\]
To solve these two equations, set
\(S_0:=c_1'+c_2'\), \(S_1:=\sum_j(t_j-1)c_j'\),
\(C_{\rm f}:=D+V-\vartheta\), and \(D_{\rm f}:=DV-\alpha\vartheta\).
At a root \(v=t_j-1\) of \(v^2+C_{\rm f}v+D_{\rm f}\),
\[
\frac v{v+\alpha}
=\frac{D_{\rm f}+\alpha v}{\alpha^2-C_{\rm f}\alpha+D_{\rm f}},
\qquad
\frac{v+1}{v+V}
=\frac{D_{\rm f}-C_{\rm f}+V-(1-V)v}{V^2-C_{\rm f}V+D_{\rm f}}.
\]
The two denominators equal \(\alpha(W-1)(V-\alpha)\) and
\(\vartheta(V-\alpha)\), respectively. The resulting system has
coefficient determinant \(\alpha(1-V)(\alpha-V)(W+\xi)\ne0\), and its solution is
\[
S_0=r_x,\qquad
S_1=-Dr_x+\xi\left((1-\alpha)r_x+
\frac{(V-\alpha)\zeta}{\zeta-1+V}X\right).
\]
The primal row of \(\mathbb T\) now gives
\(\eta br_\lambda=\chi(X-Z)+\eta ar_x-S_1\), which is the stated
formula for \(r_\lambda\). The center synchronization gives
\(r_z=X-Z+r_x\) and \(r_\mu=\Lambda-M+r_\lambda\).
Thus this linear combination is precisely \(\br\), proving
\eqref{eq:rot-annihilator} in the stated nondegenerate case.

All entries of \(Q_{\rm f}(\mathbb T)\br\) are continuous functions
of the admissible parameters: \(\alpha,b,V>0\), \(W+\xi>0\), and
\(|\zeta-1+\alpha|,|\zeta-1+V|>0\).
To justify the density used here, first take
\(\alpha,\beta,d_\tau\in(0,1)\) and \(V\ne\alpha\).
Then \(W-1>0\) and \(\vartheta>0\), so the two displayed denominators
and the coefficient determinant are nonzero. Moreover,
\[
Q_{\rm f}(1)=DV-\alpha\vartheta
=\frac{\alpha\beta\eta\tau b^2}{\nu}>0,
\]
so neither root equals one. The remaining exclusion is a repeated root.
For fixed \(\alpha,\beta,d_\tau\), its discriminant
\[
(D+V-\vartheta)^2-4(DV-\alpha\vartheta)
\]
is a rational function of \(\epsilon^2\) and tends to \((U-V)^2\) as
\(\epsilon\downarrow0\). If \(U\ne V\), this rational function is not
identically zero, so its zeros are isolated. Since
\(V=\beta+d_\tau(1-\beta)\) varies strictly with \(\beta\) in the interior,
the restrictions \(V\ne\alpha\) and \(V\ne U\) also define a dense subset.
These perturbations can be made within the admissible parameter set:
first decrease \(\epsilon\) slightly if
\(\epsilon(\rho+s)=c_1\), and then perturb the other parameters.
Taking limits therefore proves~\eqref{eq:rot-annihilator} also when roots
coincide, \(V=\alpha\), or \(\alpha=1\), \(\beta=1\), or \(d_\tau=1\).

To locate the roots of \(Q_{\rm f}\), note that
\(0<D,V\le1\) and, by~\eqref{eq:rot-small-feedback},
\[
\frac{2\vartheta}{DV}\le\frac{4\xi}{\chi d_\tau}
\le2\xi C_{\rm tr}\le\frac18.
\]
If \(|t|\ge1\), then
\[
|(t-1+D)(t-1+V)|\ge DV|t|^2
>2\vartheta|t|\ge|\vartheta(t-1+\alpha)|.
\]
Thus every root of \(Q_{\rm f}\) lies inside the unit circle.
The subspace \(\ker Q_{\rm f}(\mathbb T)\) is invariant under \(\mathbb T\),
and every eigenvalue of the restriction is a root of \(Q_{\rm f}\).
It contains \(\br\) by~\eqref{eq:rot-annihilator}, proving (ii), including
repeated roots. Finally, \(\bv+\br=-\bw_c\) follows from
\eqref{eq:rot-shift}, proving (iii).
\end{proof}

\subsection{Proof of Lemma~\ref{lem:rot-estimates}~(ii)}
\label{app:lb-rotation-transient-proof}
\begin{proof}
Write \((\mathbb T^k\br)_x\), \((\mathbb T^k\br)_\lambda\), and
\((\mathbb T^k\br)_\mu\) as \(p_k,\ell_k,m_k\), respectively.
These are components of the linear recursion, not new algorithm iterates.
By the preceding spectral proof, all these components tend to zero.
Put \(r_{\rm p}:=1-D\), \(r_{\rm d}:=1-V\), and
\(\vartheta:=(1-V)\xi\). Equation~\eqref{eq:rot-annihilator} gives
\[
p_{k+2}-(r_{\rm p}+r_{\rm d})p_{k+1}+r_{\rm p}r_{\rm d}p_k
=\vartheta\bigl(p_{k+1}-(1-\alpha)p_k\bigr).
\]
The initial values are \(p_0=r_x\) and
\(p_1=r_{\rm p}r_x+\omega_x\), where
\[
\omega_x:=\xi\left((1-\alpha)r_x+
\frac{(V-\alpha)\zeta}{\zeta-1+V}X\right).
\]
Since \(|\zeta-1|=\sqrt\nu\), the formula for \(r_x\) implies
\[
|r_x|\le\left(\frac{\sqrt\nu}{\alpha}+\frac\xi{d_\tau}\right)X
\le\frac{3X}{16}.
\]
As \(|V-\alpha|\le1\), it follows that
\(|\omega_x|\le\xi(1+d_\tau^{-1})X\). Hence
\begin{equation}
\label{eq:rot-initial-transient}
|r_x|+\frac{|\omega_x|}D
\le\left(\frac18+\xi C_{\rm tr}\right)X\le\frac{3X}{16}.
\end{equation}

For \(j\ge0\), define the convolution kernel
\(G_j:=\sum_{i=0}^jr_{\rm p}^{j-i}r_{\rm d}^i\).
Since \(0\le r_{\rm p},r_{\rm d}<1\),
\[
0\le G_j\le\sum_{i=0}^jr_{\rm p}^{j-i}\le\frac1D,
\qquad \sum_{j\ge0}G_j=\frac1{DV}.
\]
Solving the scalar recurrence gives, for \(k\ge1\),
\[
p_k=r_{\rm p}^kr_x+G_{k-1}\omega_x+
\vartheta\sum_{j=0}^{k-2}G_{k-2-j}\bigl(p_{j+1}-(1-\alpha)p_j\bigr),
\]
with an empty sum when \(k=1\). This formula follows by substitution
in the recurrence and its two initial values.
For each finite \(N\), let \(M_N:=\max_{0\le k\le N}|p_k|\).
Using~\eqref{eq:rot-initial-transient} and \(2\vartheta/(DV)\le1/4\),
\[
M_N\le\frac{3X}{16}+\frac{2\vartheta}{DV}M_N,
\qquad M_N\le\frac X4.
\]
Since \(N\) is arbitrary, this proves the primal bound.
Summing absolute values on a finite prefix instead gives
\[
\sum_{k=0}^N|p_k|
\le\frac{|r_x|}D+\frac{|\omega_x|}{DV}
+\frac{2\vartheta}{DV}\sum_{k=0}^N|p_k|.
\]
Because \(D,V\le1\),
\[
\frac{|r_x|}D+\frac{|\omega_x|}{DV}
\le\frac{|r_x|+|\omega_x|/D}{DV}.
\]
Letting \(N\) increase yields
\begin{equation}
\label{eq:rot-transient-sum}
\sum_{k\ge0}|p_k|\le\frac X{4DV}\le\frac X{\chi d_\tau}.
\end{equation}

For the dual bound, introduce the difference and weighted average
\[
e_k:=\ell_k-m_k,
\qquad q_k:=\frac{\beta(1-d_\tau)\ell_k+d_\tau m_k}{V}.
\]
The dual and center rows of \(\mathbb T\) give
\[
e_{k+1}=(1-V)e_k+\tau b(1-\beta)p_{k+1},
\qquad q_{k+1}=q_k+\frac{\tau b\beta}{V}p_{k+1}.
\]
The synchronization at \(k=0\) implies
\[
e_0=-\frac{\tau b(1-\beta)\zeta}{\zeta-1+V}X,
\qquad |e_0|\le sbX.
\]
Iterating the first recurrence and using \(|p_k|\le X/4\) gives
\[
|e_k|\le sbX+\frac{\tau b(1-\beta)}V\frac X4\le\frac54sbX.
\]
Since \(q_k\to0\), the second recurrence and the absolutely summable
bound~\eqref{eq:rot-transient-sum} give
\[
q_k=-\frac{\tau b\beta}V\sum_{j=k+1}^{\infty}p_j,
\qquad |q_k|\le\frac{s}{\chi}bX.
\]
Finally, \(\ell_k=q_k+(d_\tau/V)e_k\), with \(d_\tau/V\le1\).
Thus
\[
\|(\mathbb T^k\br)_\lambda\|=|\ell_k|
\le s\left(\frac1\chi+\frac54\right)bX
=s\left(\frac\gamma\eta+\frac54\right)bX,
\]
as required.
\end{proof}

% !TEX root = ../main.tex
\section{Proof of Theorem~\ref{thm:upper-beta1}}
\label{app:upper-beta1}

\subsection{Curvature and smoothness of the potential function}
\label{subsec:curvature-smoothness}
Let us define the constants
\begin{equation}
    L_{\calD,\bx}
    :=
    L_f + L_{\calA}L_h + \rho G_{\calA}^{2},
    \label{eq:x-weak-convexity-modulus}
\end{equation}
and 
\[
    \mu_\gamma
    :=
    \frac{1}{\gamma}-L_{\calD,\bx},
    \qquad
    L_x
    :=
    L_{\calD,\bx}+\frac{1}{\gamma},
    \qquad
    L_\lambda(s)
    :=
    \frac{1}{\rho}+\frac{1}{s},
\] where $L_f$, $L_{\calA}$, $G_{\calA}$, and $L_h$ are the constants of
Assumptions~\ref{ass:smoothness} and~\ref{ass:h}. 

The following proposition summarizes the curvature and smoothness properties
of the smoothed perturbed potential function $\mathcal P_s$.
The proof is deferred to Appendix~\ref{app:curvature-smoothness}.

\begin{proposition}
\label{prop:curvature-smoothness}
Suppose that Assumptions~\ref{ass:smoothness} and~\ref{ass:h} hold, and let the smoothing parameter  
$ 0 < \gamma < \frac{1}{2L_{\calD,\bx}}$. Then, for all $s>0$, $\mathcal P_s$ is
$\mu_\gamma $-strongly convex in $\bx$ and
$1/s$-strongly concave in $\blmbd$. Moreover, for all
$\bx,\bx'\in\RR^n$,
$\blmbd,\blmbd'\in\RR^m$,
$\bz\in\RR^n$, and
$\bmu\in\RR^m$, it holds that 
\begin{align}
    &\bigl\|
    \nabla_{\bx}\mathcal{P}_s(\bx,\blmbd,\bz,\bmu)
    -
    \nabla_{\bx}\mathcal{P}_s(\bx',\blmbd',\bz,\bmu)
    \bigr\|
    \le
    L_x\|\bx-\bx'\|
    +
    G_{\calA}\|\blmbd-\blmbd'\|,
    \label{eq:payoff-x-smoothness}
    \\
    &\bigl\|
    \nabla_{\blmbd}\mathcal{P}_s(\bx,\blmbd,\bz,\bmu)
    -
    \nabla_{\blmbd}\mathcal{P}_s(\bx',\blmbd',\bz,\bmu)
    \bigr\|
    \le
    G_{\calA}\|\bx-\bx'\|
    +
    L_\lambda(s)\|\blmbd-\blmbd'\|
    \label{eq:payoff-lambda-smoothness}.
\end{align}
In particular, $\mathcal P_s$ is $L_x$-smooth in $\bx$ and $L_\lambda(s)$-smooth in $\blmbd$.
\end{proposition}

\subsection{Proof of the iteration complexity bound}

Throughout this appendix, let
\[
L:=\sqrt2\,\max\Bigl\{L_{\calD,\bx},\ G_{\calA},\ \frac1\rho\Bigr\},
\qquad
\bar\gamma:=\frac1{4L},
\qquad
\bar\tau:=\frac1{36L}\Bigl(\frac\eta\gamma\Bigr)^2,
\]
and let \(\mathcal Y:=\{\blmbd\in\RR^m:\|\blmbd\|\le L_h+1\}\). We write
\([\mathrm{ZXSL}]\) for~\cite{zhang2020smoothed}; equation and lemma numbers
refer to its arXiv version~2010.15768v3.

\paragraph{Step 1: reduction to Smoothed GDA}
Since \(\beta=1\) and \(\bmu^0=\blmbd^0=\zero\), induction gives
\(\bmu^k=\blmbd^k\) for all \(k\). Hence the gradient of the dual regularizer
vanishes at every dual step, and the updates~\eqref{eq:splalm-updates} read
\[
\bx^{k+1}=\bx^k-\eta\Bigl[\nabla_{\bx}\calD_\rho(\bx^k,\blmbd^k)
  +\frac1\gamma(\bx^k-\bz^k)\Bigr],
\qquad
\blmbd^{k+1}=\blmbd^k+\tau\nabla_{\blmbd}\calD_\rho(\bx^{k+1},\blmbd^k),
\]
together with \(\bz^{k+1}=\bz^k+\alpha(\bx^{k+1}-\bz^k)\). None of these
updates involves \(s\). By~\eqref{eq:prox-multiplier},
\(\nabla_{\blmbd}\calD_\rho(\bx,\blmbd)=\calA(\bx)-\by(\bx,\blmbd)
=(\bq(\bx,\blmbd)-\blmbd)/\rho\), so
\[
\blmbd^{k+1}=\Bigl(1-\frac\tau\rho\Bigr)\blmbd^k+\frac\tau\rho\,\bq^k,
\qquad \|\bq^k\|\le L_h,
\]
where the bound holds because \(\bq^k\in\partial h(\by^{k+1})\) and \(h\) is
\(L_h\)-Lipschitz. Since \(L\ge1/\rho\), we have
\(\tau\le\bar\tau\le1/(144L)\le\rho\), and induction gives
\(\|\blmbd^k\|\le L_h\) for all \(k\). Hence every \(\blmbd^k\) lies in the
interior of \(\mathcal Y\), and the projection onto \(\mathcal Y\) does not
change the iterates. Consequently, \(\{(\bx^k,\blmbd^k,\bz^k)\}\) is the
sequence generated by Smoothed GDA~\cite[Algorithm~2]{zhang2020smoothed}
applied to
\(\min_{\bx\in\RR^n}\max_{\blmbd\in\mathcal Y}\calD_\rho(\bx,\blmbd)\) from
zero, with smoothing parameter \(p=1/\gamma\), primal stepsize \(\eta\), dual
stepsize \(\tau\), and averaging parameter \(\alpha\).

\paragraph{Step 2: assumptions of~\cite[Theorem~3.4]{zhang2020smoothed}}
The proof of Proposition~\ref{prop:curvature-smoothness} gives
\[
\begin{aligned}
\|\nabla_{\bx}\calD_\rho(\bx,\blmbd)-\nabla_{\bx}\calD_\rho(\bx',\blmbd')\|
&\le L_{\calD,\bx}\|\bx-\bx'\|+G_{\calA}\|\blmbd-\blmbd'\|,\\
\|\nabla_{\blmbd}\calD_\rho(\bx,\blmbd)-\nabla_{\blmbd}\calD_\rho(\bx',\blmbd')\|
&\le G_{\calA}\|\bx-\bx'\|+\frac1\rho\|\blmbd-\blmbd'\|.
\end{aligned}
\]
By the Cauchy--Schwarz inequality, both gradients are \(L\)-Lipschitz on
\(\RR^n\times\RR^m\). The set \(\mathcal Y\) is convex and compact,
\(\RR^n\) is closed and convex, and \(\calD_\rho(\bx,\cdot)\) is concave. By
Lemma~\ref{lem:effective-penalty},
\[
\calD_\rho(\bx,\blmbd)
=\max_{\bq}\Bigl\{\calL(\bx,\bq)-\frac1{2\rho}\|\bq-\blmbd\|^2\Bigr\}
\le\max_{\bq}\calL(\bx,\bq)=F(\bx),
\]
with equality for \(\blmbd\in\partial h(\calA(\bx))\). Every such
\(\blmbd\) satisfies \(\|\blmbd\|\le L_h\) and hence lies in \(\mathcal Y\).
Therefore \(\max_{\blmbd\in\mathcal Y}\calD_\rho(\bx,\blmbd)=F(\bx)\), which
is bounded below by \(\inf F>-\infty\). Thus Assumption~3.3 of
\([\mathrm{ZXSL}]\) holds with gradient Lipschitz constant \(L\).

\paragraph{Step 3: parameter conditions}
We verify condition~(3.3) of \([\mathrm{ZXSL}]\). First,
\(\gamma\le\bar\gamma\) gives \(p=1/\gamma\ge4L>3L\). Second, since
\(\gamma L\le1/4\),
\(\gamma/(1+\gamma L)\ge4\gamma/5>\gamma/2\ge\eta\), hence
\(\eta<1/(p+L)\). Third, put \(a:=\eta(p-L)\). Then
\(3\eta/(4\gamma)\le a<\eta p\le1/2\), so
\[
\frac{\eta^2(p-L)^2}{4L(1+\eta(p-L))^2}
>\frac{a^2}{16L}
\ge\frac9{256L}\Bigl(\frac\eta\gamma\Bigr)^2
>\bar\tau\ge\tau,
\]
and \(\bar\tau\le1/(144L)<1/(11L)\). Finally, the averaging parameter
\(\alpha=c\epsilon^2\le c\) satisfies
\(\alpha\le\min\{1/36,(p-L)^2/(384p(p+L)^2)\}\) by the choice of \(c\) in
Step~5.

\paragraph{Step 4: initial gap}
Let \(K(\bx,\bz;\blmbd):=\calD_\rho(\bx,\blmbd)+\frac1{2\gamma}\|\bx-\bz\|^2\),
\(d(\blmbd,\bz):=\min_{\bx}K(\bx,\bz;\blmbd)\), and
\(P(\bz):=\min_{\bx}\max_{\blmbd\in\mathcal Y}K(\bx,\bz;\blmbd)\). The
potential function of~\cite[(4.1)]{zhang2020smoothed} at the zero
initialization is \(\varphi^0=K(\zero,\zero;\zero)-2d(\zero,\zero)+2P(\zero)\).
By Proposition~\ref{prop:curvature-smoothness}, \(K(\cdot,\zero;\zero)\) is
\(\mu_\gamma\)-strongly convex, and
\(\|\nabla_{\bx}\calD_\rho(\zero,\zero)\|
=\|\nabla f(\zero)+J_{\calA}(\zero)^\top\bq(\zero,\zero)\|\le G_f+G_{\calA}L_h\).
Hence
\[
d(\zero,\zero)\ge\calD_\rho(\zero,\zero)-\frac{(G_f+G_{\calA}L_h)^2}{2\mu_\gamma}.
\]
By Step~2, \(P(\zero)=\min_{\bx}\{F(\bx)+\frac1{2\gamma}\|\bx\|^2\}\le F(\zero)\).
Moreover,
\(\calD_\rho(\zero,\zero)=f(\zero)+\min_{\by}\{h(\by)+\frac\rho2\|\calA(\zero)-\by\|^2\}\),
and the Lipschitz continuity of \(h\) gives
\(0\le F(\zero)-\calD_\rho(\zero,\zero)\le L_h^2/(2\rho)\). Combining these
bounds,
\[
\varphi^0-\inf F\le\hat\Delta
:=\Delta+\frac{L_h^2}{2\rho}+\frac{(G_f+G_{\calA}L_h)^2}{\mu_\gamma},
\]
where \(\mu_\gamma\ge1/\gamma-L\ge3/(4\gamma)>0\).

\paragraph{Step 5: rate}
For \(t\ge0\), set
\[
S_{t+1}:=\max\bigl\{\|\nabla_{\bx}\calD_\rho(\bx^{t+1},\blmbd^{t+1})\|,\
\|\nabla_{\blmbd}\calD_\rho(\bx^{t+1},\blmbd^{t+1})\|\bigr\}.
\]
Since \(\blmbd^{t+1}\) lies in the interior of \(\mathcal Y\), the normal
cone of \(\mathcal Y\) at \(\blmbd^{t+1}\) is \(\{\zero\}\), and
\((\bx^{t+1},\blmbd^{t+1})\) is an \(S_{t+1}\)-solution in the sense
of~\cite[Definition~3.1]{zhang2020smoothed}. Fix \(T\ge1\). The proof
of~\cite[Theorem~3.4]{zhang2020smoothed} distinguishes the cases (B.52) and
(B.53) and yields some \(t\in\{0,\dots,T-1\}\) with either
\[
\text{(i)}\quad S_{t+1}\le\lambda_1\max\{\alpha,\sqrt\alpha\}
\qquad\text{or}\qquad
\text{(ii)}\quad S_{t+1}\le\lambda_2\sqrt{\frac{\varphi^0-\inf F}{T\alpha}},
\]
by (B.62)--(B.63) with Lemma~B.11 in case (i) and by Lemma~B.12 in case (ii).
Here \(\lambda_1\) and \(\lambda_2\) depend only on \(L\), \(p\), \(\eta\),
\(\tau\), and the diameter \(2(L_h+1)\) of \(\mathcal Y\). Choose
\[
c:=\min\Bigl\{\frac1{36},\ \frac{(p-L)^2}{384p(p+L)^2},\
\Bigl(\frac{10}{11\lambda_1}\Bigr)^2\Bigr\},
\qquad
T:=\Bigl\lceil\Bigl(\frac{11}{10}\Bigr)^2\frac{\lambda_2^2\hat\Delta}{c\,\epsilon^4}\Bigr\rceil.
\]
Since \(\alpha=c\epsilon^2\le1\), case (i) gives
\(S_{t+1}\le\lambda_1\sqrt c\,\epsilon\le\frac{10}{11}\epsilon\), and case
(ii) gives \(S_{t+1}\le\frac{10}{11}\epsilon\) by Step~4 and the choice of
\(T\).

\paragraph{Step 6: from \(S_{t+1}\) to the KKT residual}
By~\eqref{eq:kkt-residual} and~\eqref{eq:prox-multiplier},
\(\mathcal R_{\mathrm{KKT}}(\bx^{t+1},\by^{t+1},\bq^t)
=\max\{\|\nabla_{\bx}\calD_\rho(\bx^{t+1},\blmbd^t)\|,\|\mathbf g\|\}\), where
\(\mathbf g:=\nabla_{\blmbd}\calD_\rho(\bx^{t+1},\blmbd^t)\) and
\(\blmbd^{t+1}-\blmbd^t=\tau\mathbf g\). By Step~2 and \(\tau/\rho\le1/144\),
\(\|\nabla_{\blmbd}\calD_\rho(\bx^{t+1},\blmbd^{t+1})-\mathbf g\|\le\|\mathbf g\|/144\),
hence \(\|\mathbf g\|\le\frac{144}{143}S_{t+1}\). Similarly,
\(G_{\calA}\tau\le1/144\) gives
\(\|\nabla_{\bx}\calD_\rho(\bx^{t+1},\blmbd^t)\|
\le S_{t+1}+G_{\calA}\tau\|\mathbf g\|\le\frac{144}{143}S_{t+1}\). Therefore
\(\mathcal R_{\mathrm{KKT}}(\bx^{t+1},\by^{t+1},\bq^t)\le\frac{11}{10}S_{t+1}\le\epsilon\),
and
\[
T^{\mathrm{KKT}}_\epsilon(\theta,\alpha,1;I)\le t+1\le T\le C\epsilon^{-4},
\qquad
C:=\Bigl(\frac{11}{10}\Bigr)^2\frac{\lambda_2^2\hat\Delta}{c}+1,
\]
where the last inequality uses \(\epsilon\le1\). The constant \(c\) depends
only on \(L\), \(\gamma\), \(\eta\), \(\tau\), and \(L_h\), hence only on
\(\theta\) and \(\mathbf C\); the constant \(C\) depends in addition on
\(\Delta\) through \(\hat\Delta\). Neither depends on the instance or on the
dimensions.
\hfill\(\square\)

% !TEX root = ../main.tex
\section{Proof of Theorem~\ref{thm:upper-beta0}}
\label{app:upper-beta0}

All references to~\cite{li2026smoothing} in this appendix use the numbering
of arXiv:2602.14185v2. Let
\[
\ell:=\max\Bigl\{L_{\calD,\bx},G_{\calA},\frac1\rho\Bigr\},
\qquad r:=\frac1\gamma,\qquad \delta:=\frac1s,
\qquad \mathcal Y:=\{\blmbd\in\RR^m:\|\blmbd\|\le2L_h\},
\qquad D_{\mathcal Y}:=4L_h,
\]
where \(L_{\calD,\bx}\) is defined in~\eqref{eq:x-weak-convexity-modulus}.
We take \(L_h>0\). If the supplied class constant is zero, we use the
admissible Lipschitz bound \(L_h=1/16\) throughout this appendix; the
choice \(c_s=\max\{1,16L_h\}\) below remains \(1\), and the algorithm is
unchanged. All auxiliary constants are chosen from the fixed parameters
and class bounds, independently of the instance and its dimensions.

\paragraph{Step 1: reduction to Perturbed Smoothed GDA}
Since \(\beta=0\) and \(\bmu^0=\zero\), we have \(\bmu^k=\zero\) for all
\(k\). The potential function becomes
\[
\mathcal P_s(\bx,\blmbd;\bz,\zero)
=\calD_\rho(\bx,\blmbd)+\frac r2\|\bx-\bz\|^2
-\frac\delta2\|\blmbd\|^2.
\]
The correspondence with~\cite[(3.1) and Algorithm~3]{li2026smoothing} is
\[
f_{\mathrm{Li}}=\calD_\rho,\quad
(r_x,r_y^t,c^t,\alpha^t_{\mathrm{Li}},\beta^t_{\mathrm{Li}})
=(r,\delta,\eta,\tau,\alpha),\quad
\mathcal X=\RR^n.
\]
In particular, both algorithms use the old multiplier in the primal step,
the new primal iterate in the dual step, and the new primal iterate in the
center step. By~\eqref{eq:prox-multiplier}, the dual update reads
\[
\blmbd^{k+1}
=\Bigl(1-\frac\tau\rho-\frac\tau s\Bigr)\blmbd^k
+\frac\tau\rho\bq^k,
\qquad \|\bq^k\|\le L_h.
\]
The step-size conditions in Step~3 give
\(\tau\le1/(11(\ell+1/s))<1/(1/\rho+1/s)\).
Thus both coefficients are nonnegative and sum to \(1-\tau/s\le1\).
Induction yields \(\|\blmbd^k\|\le L_h\). Every multiplier lies in the
interior of \(\mathcal Y\), so the projection onto \(\mathcal Y\) does not
change the iterates and its normal cone at \(\blmbd^k\) is
\(\{\zero\}\). The primal projection is the identity. Consequently,
SPLALM generates exactly the sequence of Perturbed Smoothed GDA applied to
\(\calD_\rho\) on \(\RR^n\times\mathcal Y\), from zero.

\paragraph{Step 2: assumptions and the unperturbed objective}
The gradient estimates in Step~2 of Appendix~\ref{app:upper-beta1} imply
\[
\begin{aligned}
\|\nabla_{\bx}\calD_\rho(\bx,\blmbd)
-\nabla_{\bx}\calD_\rho(\bx',\blmbd')\|
&\le\ell\bigl(\|\bx-\bx'\|+\|\blmbd-\blmbd'\|\bigr),\\
\|\nabla_{\blmbd}\calD_\rho(\bx,\blmbd)
-\nabla_{\blmbd}\calD_\rho(\bx',\blmbd')\|
&\le\ell\bigl(\|\bx-\bx'\|+\|\blmbd-\blmbd'\|\bigr).
\end{aligned}
\]
Hence Assumption~2.1 of~\cite{li2026smoothing} holds. Moreover,
\(\calD_\rho(\bx,\cdot)\) is concave and \(\mathcal Y\) is a compact
convex set containing zero, as required by its Assumption~2.2.
Taking \(\by=\calA(\bx)\) in~\eqref{eq:reduced-al} gives
\(\calD_\rho(\bx,\blmbd)\le F(\bx)\). For any
\(\blmbd\in\partial h(\calA(\bx))\), proximal optimality gives
\(\by(\bx,\blmbd)=\calA(\bx)\) and equality. Such a multiplier belongs
to \(\mathcal Y\), since its norm is at most \(L_h\). Therefore
\[
\max_{\blmbd\in\mathcal Y}\calD_\rho(\bx,\blmbd)=F(\bx),
\qquad \inf F>-\infty.
\]
The descent argument of~\cite[Section~4.2]{li2026smoothing} uses a lower
bound on the potential function. Accordingly, its minimum objective value
can be replaced by the infimum; attainment of \(\inf F\) is not needed.
The primal subproblems with quadratic regularization attain their minima
by strong convexity and coercivity.

\paragraph{Step 3: step-size conditions}
Set \(a_s:=r-\ell-\delta\) and \(b_s:=\ell+\delta\), and define
\[
\omega_s:=\frac1{\sqrt\delta}\,
\frac{a_s+\tau b_s(3r-2b_s)}{\tau a_s^{3/2}}.
\]
Under the correspondence in Step~1, Condition~3.2(ii)
of~\cite{li2026smoothing} reads
\begin{equation}
\begin{aligned}
&r=\Theta(\ell),\quad r>3\ell,\quad a_s>0,\quad \delta=1/s>0,\\
&\eta=\Theta(1/\ell),\quad 0<\eta<\frac1{r+\ell+\delta},\\
&\tau=\Theta(1/\ell),\quad
0<\tau\le\min\Bigl\{\frac1{11b_s},
\frac{\eta^2a_s^2}{4b_s(1+\eta a_s)^2}\Bigr\},\\
&\alpha=\Theta(\delta/\ell),\quad
0<\alpha\le\min\Bigl\{\frac1{36},
\frac{a_s^2}{384r(r+\ell+\delta)^2},
\frac1{384r\tau\omega_s^2}\Bigr\}.
\end{aligned}
\label{eq:beta0-stepsizes}
\end{equation}
The constants in these order relations are independent of \(s\).
For fixed \(r,\eta,\tau,\ell\) satisfying the first three conditions for
all sufficiently large \(s\), we have \(\omega_s^2=\Theta(s)\).
We can therefore choose a fixed, sufficiently small \(\kappa>0\), as a
function of these fixed quantities, so that
\begin{equation}
\alpha:=\frac\kappa{\ell s}
\label{eq:beta0-alpha}
\end{equation}
satisfies the last condition for all sufficiently large \(s\).
The conditions also give \(\eta<\gamma\) and \(\tau<s\), as required
by Definition~\ref{def:splalm}.

For example, all conditions hold for
\[
\gamma=\frac1{4\ell},\qquad \eta=\frac1{12\ell},\qquad
\tau=\frac1{800\ell},\qquad
\alpha=\frac1{2{,}000{,}000\,\ell s},\qquad
s\ge\max\{1,1/\ell\}.
\]
Indeed, \(2\ell\le a_s\le3\ell\) and \(b_s\le2\ell\), and
\[
\frac{\eta^2a_s^2}{4b_s(1+\eta a_s)^2}
\ge\frac1{392\ell}>\tau,\qquad
\frac{a_s^2}{384r(r+\ell+\delta)^2}\ge\frac1{13824\ell},
\qquad \omega_s^2\le800000\,\ell s.
\]
Thus \(1/(384r\tau\omega_s^2)\ge1/(1536000\,\ell s)>\alpha\),
and the remaining inequalities follow directly.

\paragraph{Step 4: initial gap for the perturbed objective}
Define
\[
\Phi_s(\bx):=\max_{\blmbd\in\mathcal Y}
\Bigl\{\calD_\rho(\bx,\blmbd)-\frac1{2s}\|\blmbd\|^2\Bigr\}.
\]
The maximization in~\eqref{eq:effective-penalty}, with \(\bmu=\zero\),
has an optimal multiplier \(s\bq/(\rho+s)\), where
\(\|\bq\|\le L_h\). This multiplier belongs to \(\mathcal Y\), so
restricting the maximization to \(\mathcal Y\) does not change its value.
It follows that
\begin{equation}
\Phi_s(\bx)=\calD_{\rho+s}(\bx,\zero),\qquad
F(\bx)-\frac{L_h^2}{2(\rho+s)}\le\Phi_s(\bx)\le F(\bx).
\label{eq:beta0-perturbed-objective}
\end{equation}
Let
\[
\begin{aligned}
v_0&:=\calD_\rho(\zero,\zero),\qquad
 d_0:=\min_{\bx}\Bigl\{\calD_\rho(\bx,\zero)
                         +\frac1{2\gamma}\|\bx\|^2\Bigr\},\\
p_0&:=\min_{\bx}\Bigl\{\Phi_s(\bx)+\frac1{2\gamma}\|\bx\|^2\Bigr\}.
\end{aligned}
\]
Since \(\mu_\gamma=1/\gamma-L_{\calD,\bx}>0\), the argument in
Step~4 of Appendix~\ref{app:upper-beta1} gives
\[
d_0\ge v_0-\frac{(G_f+G_{\calA}L_h)^2}{2\mu_\gamma},\qquad
p_0\le F(\zero),\qquad v_0\ge F(\zero)-\frac{L_h^2}{2\rho}.
\]
At zero initialization, the initial gap in
\cite[Definition~3.1(ii)]{li2026smoothing}, with the infimum convention
explained in Step~2, is
\(\Delta_{\Psi_2}=v_0-2d_0+2p_0-\inf\Phi_s\). Hence
\begin{equation}
\begin{aligned}
\Delta_{\Psi_2}
&\le\Delta+\frac{L_h^2}{2\rho}
+\frac{(G_f+G_{\calA}L_h)^2}{\mu_\gamma}
+\frac{L_h^2}{2(\rho+s)}\\
&\le\bar\Delta:=\Delta+\frac{L_h^2}{\rho}
+\frac{(G_f+G_{\calA}L_h)^2}{\mu_\gamma}.
\end{aligned}
\label{eq:beta0-initial-gap}
\end{equation}
The bound \(\bar\Delta\) is uniform over the problem class and independent
of \(s\), the target accuracy, and the dimensions.

\paragraph{Step 5: iteration complexity for game stationarity}
For \(k\ge0\), set
\[
S_{k+1}:=\max\Bigl\{
\|\nabla_{\bx}\calD_\rho(\bx^{k+1},\blmbd^{k+1})\|,
\|\nabla_{\blmbd}\calD_\rho(\bx^{k+1},\blmbd^{k+1})\|\Bigr\}.
\]
Both normal-cone terms vanish by Step~1, so this is the game-stationarity
measure for the unperturbed function \(\calD_\rho\) in
\cite[Theorem~3.3(i)]{li2026smoothing}.
The descent estimate in its Section~4.2 and estimates (4.19)--(4.23),
with~\eqref{eq:beta0-alpha}, give a constant
\(C_*=C_*(\eta\ell,\tau\ell,r/\ell,\kappa)>0\) such that, among the first
\(K\ge1\) updates, some \(0\le k<K\) satisfies
\begin{equation}
S_{k+1}\le C_*\ell\sqrt{\frac{s\bar\Delta}{K}}
+\frac{D_{\mathcal Y}}s.
\label{eq:beta0-gs-rate}
\end{equation}
Here we take \(s\) sufficiently large that \(\ell s\ge1\), so the
term of order \(\sqrt{\ell\bar\Delta/K}\) is absorbed into the first
term. We average the descent estimate over \(k=0,\ldots,K-1\);
a fixed factor in this averaging is included in \(C_*\).
Thus the index in~\eqref{eq:beta0-gs-rate} refers to an output among the
first \(K\) updates of Definition~\ref{def:complexity-measure}.

Choose
\begin{equation}
c_s:=\max\{1,4D_{\mathcal Y}\}=\max\{1,16L_h\},\qquad
s:=\frac{c_s}{\epsilon},\qquad
K:=\max\Bigl\{1,\Bigl\lceil
16C_*^2\ell^2c_s\bar\Delta\,\epsilon^{-3}
\Bigr\rceil\Bigr\}.
\label{eq:beta0-budget}
\end{equation}
Then \(D_{\mathcal Y}/s\le\epsilon/4\), and the first term
in~\eqref{eq:beta0-gs-rate} is at most \(\epsilon/4\).
Consequently, some \(0\le k<K\) has \(S_{k+1}\le\epsilon/2\).

\paragraph{Step 6: from \(S_{k+1}\) to the KKT residual}
Write \(\mathbf d_k:=\blmbd^{k+1}-\blmbd^k\). The dual update and
Step~2 yield
\[
\|\mathbf d_k\|\le\tau\Bigl(S_{k+1}
+\frac{\|\mathbf d_k\|}\rho+\frac{L_h}s\Bigr),\qquad
\|\mathbf d_k\|\le\frac\tau{1-\tau/\rho}
\Bigl(S_{k+1}+\frac{L_h}s\Bigr).
\]
By~\eqref{eq:kkt-residual} and~\eqref{eq:prox-multiplier}, the output
residual is
\[
R_k:=\mathcal R_{\mathrm{KKT}}(\bx^{k+1},\by^{k+1},\bq^k)
=\max\Bigl\{\|\nabla_{\bx}\calD_\rho(\bx^{k+1},\blmbd^k)\|,
\|\nabla_{\blmbd}\calD_\rho(\bx^{k+1},\blmbd^k)\|\Bigr\}.
\]
Both gradients change by at most \(\ell\|\mathbf d_k\|\) when the
multiplier changes from \(\blmbd^{k+1}\) to \(\blmbd^k\). Since
\(\tau\ell\le1/11\) and \(\tau/\rho\le1/11\), we obtain
\begin{equation}
R_k\le\Bigl(1+\frac{\tau\ell}{1-\tau/\rho}\Bigr)S_{k+1}
+\frac{\tau\ell}{1-\tau/\rho}\frac{L_h}s
\le\frac{11}{10}S_{k+1}+\frac{L_h}{10s}.
\label{eq:beta0-residual-transfer}
\end{equation}
For the index selected in Step~5, \(c_s\ge16L_h\) gives
\[
R_k\le\frac{11}{20}\epsilon+\frac1{160}\epsilon
=\frac{89}{160}\epsilon<\epsilon.
\]
For all sufficiently small \(\epsilon\in(0,1]\), it follows that
\[
T^{\mathrm{KKT}}_\epsilon(\theta(\epsilon),\alpha,0;I)
\le k+1\le K\le C'\epsilon^{-3},\qquad
C':=1+16C_*^2\ell^2c_s\bar\Delta.
\]
Finally, \(\alpha=\kappa\epsilon/(\ell c_s)\), so
\(c_\alpha=\kappa/(\ell c_s)>0\). Our fixed choice of \(\kappa\)
ensures that \(c_\alpha\) depends only on
\(\rho,\gamma,\eta,\tau,\mathbf C\), and \(C'\) depends in addition
on \(\Delta\). Neither constant depends on the instance, the dimensions,
or \(\epsilon\).
\hfill\(\square\)

% Verification status (2026-10-14): the existence and parameter dependence
% of C_* in eq:beta0-gs-rate are taken from GPT Pro's conversion of Li26 v2,
% Section 4.2, (4.19)--(4.23), recorded in
% notes/gpt-proofs/beta0-upper-check-20261013.md, item 7.
% That conversion, including its fixed averaging factor, has not been
% independently verified line by line. Reading the cited statements and
% estimates does not constitute such a verification.

% !TEX root = ../main.tex
\section{An accuracy-independent schedule with increasing \texorpdfstring{$s_k$}{s\_k}}
\label{sec:upper-bound}

This appendix gives a parameter schedule that does not depend on
\(\epsilon\). With \(\rho\) fixed, \(s_k=\Theta((k+1)^{1/3})\), and diminishing
center stepsizes, SPLALM finds an \(\epsilon\)-KKT point in
\(\mathcal O(\epsilon^{-3})\) iterations (Theorem~\ref{thm:main}). The
parameters vary with \(k\), so the schedule lies outside the family
\(\calC_\theta\) of Definition~\ref{def:algorithm-family}.

Throughout this section, we additionally impose the following assumption.
It is not part of the problem class \(\mathcal I(\mathbf C,\Delta)\) of
Definition~\ref{def:problem-class}.

\begin{assumption}[Level-bounded objective]
\label{ass:level-boundedness}
For every \(c\in\RR\), the sublevel set
\(\{\bx\in\RR^n:F(\bx)\le c\}\) is empty or compact.
\end{assumption}

\subsection{Algorithm with increasing \texorpdfstring{$s_k$}{s\_k}}
\label{sec:method}

This section uses the updates of Definition~\ref{def:splalm} with
iteration-dependent parameters \(s_k\), \(\eta_k\), \(\tau_k\), \(\alpha_k\),
and \(\beta_k\), and with an arbitrary initialization; see
Algorithm~\ref{alg:ddlalm}. By~\eqref{eq:reduced-al}--\eqref{eq:potential},
the first two updates read
\begin{equation}\label{eq:explicit-updates}
    \begin{aligned}
        \bx^{k+1}
        &=
        \bx^k-\eta_k
        \Bigl[
            \nabla f(\bx^k)
            +J_{\calA}(\bx^k)^\top\bq(\bx^k,\blmbd^k)
            +\frac{1}{\gamma}(\bx^k-\bz^k)
        \Bigr], \\
        \blmbd^{k+1}
        &=
        \blmbd^k+\tau_k
        \Bigl[
            \calA(\bx^{k+1})
            -\by(\bx^{k+1},\blmbd^k)
            -\frac{1}{s_k}(\blmbd^k-\bmu^k)
        \Bigr],
    \end{aligned}
\end{equation}
and the outputs \((\bx^{k+1},\by^{k+1},\bq^k)\) are those of
Definition~\ref{def:splalm}.

\begin{algorithm}[t]
    \caption{Smoothed-Perturbed Linearized Augmented Lagrangian Method (SPLALM)}
    \label{alg:ddlalm}
    \small
    \begin{algorithmic}[1]
        \STATE \textbf{Input:} Initial
        \(\bx^0,\bz^0\in\RR^n\) and
        \(\blmbd^0,\bmu^0\in\RR^m\); 
        Parameters: \(\rho>0\) and \(\gamma>0\);
        sequences: \(\{\eta_k\}\), \(\{\tau_k\}\), \(\{\alpha_k\}\), and
        \(\{\beta_k\}\);
        and iteration limit \(K\).
        \FOR{\(k=0,1,\ldots,K-1\)}
            \STATE  \(\bx^{k+1}=\bx^k- \eta_k \nabla_{\bx} \mathcal P_{s_k}(\bx^{k},\blmbd^{k};\bz^{k},\bmu^{k})\);
            \STATE  \(\blmbd^{k+1}=\blmbd^k+\tau_k
            \nabla_{\blmbd}\mathcal P_{s_k}
            (\bx^{k+1},\blmbd^{k};\bz^{k},\bmu^{k})\).
            \STATE 
            \(\bz^{k+1}=\bz^k+ \alpha_k(\bx^{k+1}-\bz^k)\).
            \STATE  \(\bmu^{k+1}=\bmu^k+ \beta_k(\blmbd^{k+1}-\bmu^k)\).
        \ENDFOR
    \end{algorithmic}
\end{algorithm}

We use fixed primal and dual stepsizes together with diminishing center
stepsizes. Choose \(\rho>0\) and \(\gamma>0\) satisfying
Proposition~\ref{prop:curvature-smoothness}. With \(L_x\) and
\(\mu_\gamma\) defined in Section~\ref{subsec:curvature-smoothness} and
\(L_{d,0}\) defined in Section~\ref{sec:complexity-analysis}, set \(\eta_k\equiv\eta\) and
\(\tau_k\equiv\tau\). An admissible choice is
\begin{equation}\label{eq:ddlalm-parameter-schedule}
\begin{aligned}
    s_k &= c_s(k+1)^{1/3},
    \qquad
    \eta = \frac{1}{2L_x},
    \qquad
    \tau =
    \left(
        2L_{d,0}
        +\frac{72L_xG_{\cal A}^2}{\mu_\gamma^2}
    \right)^{-1},
    \\
    \alpha_k &=
    \frac{\min\{s_0,\gamma\mu_\gamma\tau\}}{36s_k},
    \qquad
    \beta_k=\frac{\theta_\mu}{k+k_0},
\end{aligned}
\end{equation}
where \(c_s>0\), \(\theta_\mu\in(0,1]\), and \(k_0\ge1\) are constants
independent of \(K\). Consequently,
\[
    s_k=\Theta((k+1)^{1/3}),
    \qquad
    \alpha_k=\Theta((k+1)^{-1/3}),
    \qquad
    \beta_k=\Theta((k+1)^{-1}).
\]

\begin{theorem}
\label{thm:main}
Suppose that Assumptions~\ref{ass:level-boundedness},
\ref{ass:smoothness}, and~\ref{ass:h} hold.
Let the sequence generated by \Cref{alg:ddlalm} use
\eqref{eq:ddlalm-parameter-schedule}, with $\rho$ and $\gamma$ as specified
above. For $K\ge2$, choose a tail-best index
\[
    \kappa\in\argmin_{k\in\{\lceil K/2\rceil,\ldots,K-1\}}
    \left\{
        \|\Delta_{\bx}^{k+1}\|^2
        +\|\Delta_{\blmbd}^{k+1}\|^2
        +\alpha_k\|\bx^{k+1}-\bz^k\|^2
    \right\}.
\]
Then it holds that
\[
\max\left\{
\left[
    \left\|
        \nabla f({\bx}^{\kappa+1})
        +
        J_{\cal A}({\bx}^{\kappa+1})^\top{\bq}^{\kappa}
    \right\|^2
\right],
\;
\left[
    \left\|
        {\cal A}({\bx}^{\kappa+1})-{\by}^{\kappa+1}
    \right\|^2
\right]
\right\}
=
\mathcal O(K^{-2/3}),
\]
with \({\bq}^{\kappa}\in\partial h({\by}^{\kappa+1})\).
Consequently, given \(\epsilon>0\), the iteration complexity of
\Cref{alg:ddlalm} to obtain an \(\epsilon\)-KKT point
is of order \(\mathcal O(\epsilon^{-3})\).
\end{theorem}



\subsection{Complexity Analysis}
\label{sec:complexity-analysis}
\label{sec:basic-descent}
We use the constants \(L_{\calD,\bx}\), \(\mu_\gamma\), \(L_x\), and
\(L_\lambda(s)\) and Proposition~\ref{prop:curvature-smoothness} from
Section~\ref{subsec:curvature-smoothness}.
 
\subsubsection{Value Functions and Lyapunov Function}
\label{subsec:value-functions}

Fix $s>0$ and $(\bz,\bmu)$. By Proposition~\ref{prop:curvature-smoothness},
$\mathcal P_s(\cdot,\cdot;\bz,\bmu)$ is strongly convex in $\bx$ and strongly
concave in $\blmbd$. Hence the following value functions are well defined,
and the minimizer and maximizer below are unique:
\begin{align}
    d_s(\blmbd;\bz,\bmu)
    &:= \min_{\bx\in\RR^n}\mathcal P_s(\bx,\blmbd;\bz,\bmu),
    &
    \bx_s(\blmbd;\bz,\bmu)
    &:= \argmin_{\bx\in\RR^n}\mathcal P_s(\bx,\blmbd;\bz,\bmu),
    \label{eq:dual-value}\\
    p_s(\bz,\bmu)
    &:= \max_{\blmbd\in\RR^m} d_s(\blmbd;\bz,\bmu),
    &
    \blmbd_s^\star(\bz,\bmu)
    &:= \argmax_{\blmbd\in\RR^m} d_s(\blmbd;\bz,\bmu).
    \label{eq:saddle-value}
\end{align}
Strong convexity--strong concavity also yields strong duality,
\begin{equation}
    p_s(\bz,\bmu)
    = \min_{\bx\in\RR^n}\max_{\blmbd\in\RR^m}
      \mathcal P_s(\bx,\blmbd;\bz,\bmu),
    \label{eq:minimax-equality}
\end{equation}
and the unique saddle point is
$(\bx_s^\star(\bz,\bmu),\blmbd_s^\star(\bz,\bmu))$, where
\[
    \bx_s^\star(\bz,\bmu)
    :=\bx_s\bigl(\blmbd_s^\star(\bz,\bmu);\bz,\bmu\bigr).
\]

\blue{Following the three-value Lyapunov construction for smoothed nonconvex--concave minimax methods~\cite{zhang2020smoothed}, and its perturbed-smoothed extension~\cite{li2026smoothing}, we adapt the construction to the reduced augmented-Lagrangian setting and define}
\begin{equation}
\begin{aligned}
    \Psi_s(\bx,\blmbd;\bz,\bmu)
    &:={}
    \mathcal P_s(\bx,\blmbd;\bz,\bmu)
    -2d_s(\blmbd;\bz,\bmu)
    +2p_s(\bz,\bmu) \\
    &= \underbrace{\mathcal P_s(\bx,\blmbd;\bz,\bmu)-d_s(\blmbd;\bz,\bmu)}_{\text{primal gap at fixed }\blmbd}
     + \underbrace{p_s(\bz,\bmu)-d_s(\blmbd;\bz,\bmu)}_{\text{dual gap}}
     + p_s(\bz,\bmu).
\end{aligned}
\label{eq:lyapunov-decomposition}
\end{equation}
Both gaps are nonnegative, so $\Psi_s\ge p_s$.

\paragraph{Iteration-$k$ shorthand.}
At iteration $k$ we evaluate the above objects at $(s,\bmu)=(s_k,\bmu^k)$,
replace the subscript $s$ by $k$, and suppress $\bmu^k$. For example,
$\mathcal P_k(\bx,\blmbd;\bz):=\mathcal P_{s_k}(\bx,\blmbd;\bz,\bmu^k)$ and
$d_k(\blmbd;\bz):=d_{s_k}(\blmbd;\bz,\bmu^k)$; the quantities
$p_k$, $\bx_k$, $\blmbd_k^\star$, $\bx_k^\star$, and $\Psi_k$ are defined
analogously. We also write
$\Delta_{\bx}^{k+1}:=\bx^{k+1}-\bx^k$, and define
$\Delta_{\blmbd}^{k+1}$, $\Delta_{\bz}^{k+1}$, and $\Delta_{\bmu}^{k+1}$
in the same way.

\subsubsection{Basic One-Step Estimate}
\label{subsec:basic-one-step}

Throughout this subsection, we assume that 
\begin{equation}
    s_{k+1}\ge s_k>0,\qquad
    0 < \tau_k\le s_k,\qquad
    0<\beta_k\le1,\qquad
    0\le\alpha_k\le2.
    \label{eq:basic-parameter-condition}
\end{equation} 
The following lemma gives a basic one-step descent estimate of the potential function $\mathcal P$ from iteration $k$ to $k+1$.

\begin{lemma}
\label{lem:basic-one-step}
Suppose Assumptions~\ref{ass:smoothness} and~\ref{ass:h} hold.
Let the iterates be generated by Algorithm~\ref{alg:ddlalm} with the
parameters satisfying \eqref{eq:basic-parameter-condition}.  Then, for
every $k\ge0$, the following holds:
\begin{equation}
\label{eq:basic-one-step}
\begin{aligned}
    &\mathcal P_{k+1}(\bx^{k+1},\blmbd^{k+1};\bz^{k+1})
    -\mathcal P_k(\bx^k,\blmbd^k;\bz^k) \\
    \le\;&
    -\left(\frac{1}{\eta_k}-\frac{L_x}{2}\right)
      \|\Delta_{\bx}^{k+1}\|^2
    +\left(\frac{1}{\tau_k}-\frac{1}{2s_k}\right)
      \|\Delta_{\blmbd}^{k+1}\|^2 \\
    &-\frac{\alpha_k(2-\alpha_k)}{2\gamma}
      \|\bx^{k+1}-\bz^k\|^2
    +\frac{\beta_k(2-\beta_k)}{2s_k}
      \|\blmbd^{k+1}-\bmu^k\|^2 +\frac{1}{2}
      \left(\frac{1}{s_k}-\frac{1}{s_{k+1}}\right)
      \|\blmbd^{k+1}-\bmu^{k+1}\|^2 .
\end{aligned}
\end{equation}
\end{lemma}


As shown in Lemma~\ref{lem:basic-one-step}, the positive ascent term
generated by the $\blmbd$-update prevents the potential
$\mathcal P_k$ from yielding a direct descent estimate. We therefore relate
this term to the one-step increase of the dual value function, as developed
in Lemma~\ref{lem:dual-value-gain}.  This relation allows us to construct a Lyapunov function with a descent property. Before proceeding, we establish several regularity properties of the value functions \(d_k\) and \(p_k\), with their gradient formulas following from Danskin's theorem.
\begin{lemma}
\label{lem:v2-value-regularity}
Under the conditions of Proposition~\ref{prop:curvature-smoothness}, for all $k\ge0$, the following properties hold:
\begin{enumerate}[label=(\roman*)]
\item
For every \(\bz \in \RR^n \), the minimizer
\(\bx_k(\blmbd;\bz)\) is unique and
\(\frac{G_{\cal A}}{\mu_\gamma}\)-Lipschitz continuous in
\(\blmbd\).

\item for every \(\bz\in\RR^n\),
\(d_k(\cdot;\bz)\) is \(1/s_k\)-strongly concave and
\(L_{d,k}\)-smooth, with gradient
\[
    \nabla_{\blmbd}d_k(\blmbd;\bz)
    =
    \nabla_{\blmbd}\mathcal P_k
    \bigl(\bx_k(\blmbd;\bz),\blmbd;\bz\bigr),
\]
where \(L_{d,k}:=\frac1\rho+\frac1{s_k}+\frac{G_{\cal A}^2}{\mu_\gamma}\).

\item
The function \(p_k(z)\) is continuously differentiable, with gradient
\[
    \nabla p_k(\bz)
    =
    \frac1\gamma
    \bigl(\bz-\bx_k^\star(\bz)\bigr),
\]
and has a \(1/\gamma\)-Lipschitz continuous gradient.
\end{enumerate}
\end{lemma}

The following lemma controls the dual ascent term.
\begin{lemma}
\label{lem:dual-value-gain}
Suppose that the conditions of
Lemma~\ref{lem:v2-value-regularity} hold, and let the iterates
be generated by Algorithm~\ref{alg:ddlalm} with parameters satisfying
\eqref{eq:basic-parameter-condition}.
For every $k\ge 0$, assume that
\begin{equation}
\label{eq:dual-gain-condition}
    \frac{1}{\tau_k}+\frac{1}{2s_k} \ge 2L_{d,k}.
\end{equation}
Then the following holds
\begin{equation*}
\begin{aligned}
& \left(\frac{1}{\tau_k}-\frac{1}{2s_k}\right)
  \|\Delta_{\boldsymbol{\lambda}}^{k+1}\|^2 \\
& \le 2\!\bigl[
    d_k(\boldsymbol{\lambda}^{k+1};\mathbf{z}^k)
    -d_k(\boldsymbol{\lambda}^k;\mathbf{z}^k)
  \bigr]
  -\frac{L_{d,k}}{2}
    \|\Delta_{\boldsymbol{\lambda}}^{k+1}\|^2
  +\frac{4G_{\mathcal A}^2}{L_{d,k}\mu_\gamma}
  \left[
    \mathcal P_k(
      \mathbf{x}^{k+1},\boldsymbol{\lambda}^k;\mathbf{z}^k)
    -d_k(\boldsymbol{\lambda}^k;\mathbf{z}^k)
  \right].
\end{aligned}
\end{equation*}
Moreover, it holds that
\begin{equation}
\begin{aligned}
&\mathcal P_{k+1}(\bx^{k+1},\blmbd^{k+1};\bz^{k+1})
-\mathcal P_k(\bx^k,\blmbd^k;\bz^k)
\\
&\le
-\left[
\frac1{\eta_k}-\frac{L_x}{2}
-
\frac{2G_{\mathcal A}^2(1/\eta_k+L_x)^2}
{L_{d,k}\mu_\gamma^2}
\right]
\|\Delta_{\bx}^{k+1}\|^2
-\frac{L_{d,k}}2\|\Delta_{\blmbd}^{k+1}\|^2 \\
&\qquad
+2\bigl[d_k(\blmbd^{k+1};\bz^k)-d_k(\blmbd^k;\bz^k)\bigr]
-\frac{\alpha_k(2-\alpha_k)}{2\gamma}
\|\bx^{k+1}-\bz^k\|^2 \\
&\qquad
+\frac{\beta_k(2-\beta_k)}{2s_k}
\|\blmbd^{k+1}-\bmu^k\|^2
+\frac12\left(\frac1{s_k}-\frac1{s_{k+1}}\right)
\|\blmbd^{k+1}-\bmu^{k+1}\|^2.
\end{aligned}
\label{eq:dual-value-full-step}
\end{equation}
\end{lemma}

The potential $\mathcal P_k$ increases along the dual update, so its variation alone does not capture progress in
both variables. We therefore seek a merit function that accounts for both
primal suboptimality and progress in dual maximization. To account for
these two effects, consider the partial potential
\begin{equation}
\label{eq:partial-potential}
    \mathcal P_k-cd_k
    =(\mathcal P_k-d_k)-(c-1)d_k,
    \qquad c>1.
\end{equation}
The first term measures the primal optimality gap, while the second term
decreases as the dual value increases. In the following analysis, we set
$c=2$, consistent with the coefficient of the dual-value increment in
\eqref{eq:dual-value-full-step}.

\begin{lemma}
\label{lem:partial-lyapunov-z-step}
Under the conditions of Lemma~\ref{lem:dual-value-gain}, for every $k\ge0$, the following holds
\begin{equation}
\begin{aligned}
&\Big[
\mathcal P_{k+1}(\bx^{k+1},\blmbd^{k+1};\bz^{k+1})
-2d_k(\blmbd^{k+1};\bz^{k+1})
\Big] -
\Big[
\mathcal P_k(\bx^k,\blmbd^k;\bz^k)
-2d_k(\blmbd^k;\bz^k)
\Big]\\
\le\,&
-\left[
\frac1{\eta_k}-\frac{L_x}{2}
-\frac{2G_{\mathcal A}^2(1/\eta_k+L_x)^2}
{L_{d,k}\mu_\gamma^2}
\right]
\|\Delta_{\bx}^{k+1}\|^2 -\frac{L_{d,k}}{2}\|\Delta_{\blmbd}^{k+1}\|^2 \\
&\qquad
+\frac{\alpha_k(2-\alpha_k)}{2\gamma}
\|\bx^{k+1}-\bz^k\|^2
-\frac{2\alpha_k}{\gamma}
\left\langle
\bx^{k+1}
-\bx_k(\blmbd^{k+1};\bz^{k+1}),
\bx^{k+1}-\bz^k
\right\rangle  \\
&\qquad
+\frac{\beta_k(2-\beta_k)}{2s_k}
\|\blmbd^{k+1}-\bmu^k\|^2
+\frac12\left(\frac1{s_k}-\frac1{s_{k+1}}\right)
\|\blmbd^{k+1}-\bmu^{k+1}\|^2.
\end{aligned}
\label{eq:partial-lyapunov-z-step}
\end{equation}
\end{lemma}
Equation~\eqref{eq:partial-lyapunov-z-step} gives the one-step estimate
for the partial potential in \eqref{eq:partial-potential} with $c=2$.
The $\bx$- and
$\blmbd$-updates contribute the first two terms on the right-hand side.
The $\bz$-update, however, produces a positive quadratic term and an
indefinite inner-product term. Hence, the partial potential does not yet
provide a sufficient descent estimate for the full iteration.

At an exact saddle point of
\(\min_{\bx}\max_{\blmbd}\mathcal P_k(\bx,\blmbd;\bz)\),
strong duality gives
\[
    \mathcal P_k\bigl(\bx_k^\star(\bz),
    \blmbd_k^\star(\bz);\bz\bigr)
    =d_k\bigl(\blmbd_k^\star(\bz);\bz\bigr)
    =p_k(\bz),
\]
and therefore
\[
    \mathcal P_k\bigl(\bx_k^\star(\bz),
    \blmbd_k^\star(\bz);\bz\bigr)
    -2d_k\bigl(\blmbd_k^\star(\bz);\bz\bigr)
    =-p_k(\bz).
\]
Consequently, any $\bz$-update that decreases the saddle value $p_k$
increases the partial potential $\mathcal P_k-2d_k$. This motivates us
to examine whether an additional value-function term can restore a descent
estimate.

We therefore add the value-function increment to the partial potential in
\eqref{eq:partial-potential}:
$2[p_k(\bz^{k+1})-p_k(\bz^k)]$ and combine the two estimates directly.
The resulting one-step bound is stated next.

\begin{lemma}
\label{lem:z-step-p}
Under the conditions of Lemma~\ref{lem:v2-value-regularity} and
$0\le\alpha_k\le2$, the $\bz$-update satisfies
\begin{equation}
\begin{aligned}
p_k(\bz^{k+1})-p_k(\bz^k)
\le\,
-\frac{\alpha_k(2-\alpha_k)}{2\gamma}
\|\bx^{k+1}-\bz^k\|^2 +
\frac{\alpha_k}{\gamma}
\left\langle
\bx^{k+1}-\bx_k^\star(\bz^k),
\bx^{k+1}-\bz^k
\right\rangle.
\end{aligned}
\label{eq:z-step-p}
% \notag
\end{equation}
\end{lemma}

\begin{lemma}
\label{lem:augmented-one-step}
Under the conditions of Lemma~\ref{lem:partial-lyapunov-z-step}, for every
$k\ge0$,
\begin{equation}
\begin{aligned}
&\Big[
\mathcal P_{k+1}(\bx^{k+1},\blmbd^{k+1};\bz^{k+1})
-2d_k(\blmbd^{k+1};\bz^{k+1})
+2p_k(\bz^{k+1})
\Big] \\
&\quad-
\Big[
\mathcal P_k(\bx^k,\blmbd^k;\bz^k)
-2d_k(\blmbd^k;\bz^k)
+2p_k(\bz^k)
\Big] \\
&\le
-\left[
\frac1{\eta_k}-\frac{L_x}{2}
-\frac{2G_{\mathcal A}^2(1/\eta_k+L_x)^2}
{L_{d,k}\mu_\gamma^2}
\right]
\|\Delta_{\bx}^{k+1}\|^2
-\frac{L_{d,k}}{2}\|\Delta_{\blmbd}^{k+1}\|^2 \\
&\qquad
-\frac{\alpha_k(2-\alpha_k)}{2\gamma}
\|\bx^{k+1}-\bz^k\|^2
+\frac{2\alpha_k}{\gamma}
\left\langle
\bx_k(\blmbd^{k+1};\bz^{k+1})
-\bx_k^\star(\bz^k),
\bx^{k+1}-\bz^k
\right\rangle \\
&\qquad
+\frac{\beta_k(2-\beta_k)}{2s_k}
\|\blmbd^{k+1}-\bmu^k\|^2
+\frac12\left(\frac1{s_k}-\frac1{s_{k+1}}\right)
\|\blmbd^{k+1}-\bmu^{k+1}\|^2.
\end{aligned}
\label{eq:augmented-one-step}
\end{equation}
\end{lemma}

Two gaps remain before $\Phi_k$ admits a genuine descent inequality: the augmented estimate still carries an indefinite cross term and positive multiplier/center-increment terms, and its left-hand side compares $d_k, p_k$ at both endpoints rather than the true $k \to k+1$ potential difference, which requires $d_{k+1}, p_{k+1}$. 

For that, we establish boundedness of the multipliers directly, rather than assuming compactness a priori: both multipliers are controlled by the AL structure and the inner residual, not by projection onto a compact set.
\begin{lemma}
\label{lem:current-multiplier-boundedness}
Suppose Assumptions~\ref{ass:smoothness} and~\ref{ass:h} hold.
Let the iterates be generated by Algorithm~\ref{alg:ddlalm} with
\eqref{eq:basic-parameter-condition} and
\begin{equation}
\label{eq:multiplier-boundedness-condition}
    \tau_k\left(\frac1\rho+\frac1{s_k}\right)\le1.
\end{equation}
Define the constant independent of $k$
\[
    M:=\max\{\|\blmbd^0\|,\|\bmu^0\|,L_h\}.
\]
Then, for every $k\ge0$, the following holds
\[
    \|\blmbd^k\|\le M,
    \qquad
    \|\bmu^k\|\le M.
\]
Moreover, for every $k\ge0$ and every $\bz\in\RR^n$, it holds that
\[
    \|\blmbd_k^\star(\bz)\|\le M.
\]
\end{lemma}

% The augmented one-step estimate also contains the indefinite term
% \[
% \left\langle
% \bx_k(\blmbd^{k+1};\bz^{k+1})-\bx_k^\star(\bz^k),
% \bx^{k+1}-\bz^k
% \right\rangle.
% \]
% Its first factor combines two effects: the dual iterate may differ from
% the dual maximizer, and the saddle-point primal component may move when
% $\bz$ changes. The next lemma quantifies both effects. The first estimate
% relates the dual suboptimality gap $p_k-d_k$ to the corresponding primal
% distance, while the second estimate controls the sensitivity of the
% primal minimizers to the center variable.

\begin{lemma}
\label{lem:saddle-center-stability}
Under the conditions of Proposition~\ref{prop:curvature-smoothness}, for every $k\ge0$, $\blmbd\in\RR^m$, and $\bz\in\RR^n$, it holds that
\begin{equation}
\label{eq:current-saddle-stability}
\|\bx_k(\blmbd;\bz)-\bx_k^\star(\bz)\|^2
\le
\frac{2}{\mu_\gamma}
\bigl[p_k(\bz)-d_k(\blmbd;\bz)\bigr].
\end{equation}
For any $\bz,\bz'\in\RR^n$,
\begin{equation}
\label{eq:current-center-stability}
\max\bigl\{
\|\bx_k(\blmbd;\bz')-\bx_k(\blmbd;\bz)\|,
\|\bx_k^\star(\bz')-\bx_k^\star(\bz)\|
\bigr\}
\le\frac{1}{\gamma\mu_\gamma}\|\bz'-\bz\|.
\end{equation}
\end{lemma}

We now establish an approximate descent inequality for
\Cref{alg:ddlalm} with respect to the Lyapunov function defined in \eqref{eq:lyapunov-decomposition}. We assume 
the parameters are chosen as
\begin{equation}
\label{eq:parameter-choice}
\begin{gathered}
    0<\gamma<\frac{1}{2L_{\calD,\bx}},\qquad
    0<s_k\le s_{k+1},\qquad
    0<\beta_k\le1,\qquad
    0<\eta_k\le\frac{1}{2L_x},\\
    0<\tau_k\le\left(\max\left\{2L_{d,k},\,
        \frac{36G_{\mathcal A}^2}{\eta_k\mu_\gamma^2}\right\}\right)^{-1},\qquad
    0<\alpha_k\le\frac1{36}\min\left\{1,\,
        \frac{\gamma\mu_\gamma\tau_k}{s_k}\right\},
\end{gathered}
\end{equation}
where $L_{d,k}:=\rho^{-1}+s_k^{-1}+G_{\mathcal A}^2/\mu_\gamma$.
Since $L_{d,k}\ge\rho^{-1}+s_k^{-1}$, \eqref{eq:parameter-choice} implies the conditions of Proposition~\ref{prop:curvature-smoothness},
Lemma~\ref{lem:current-multiplier-boundedness},
\eqref{eq:basic-parameter-condition},
\eqref{eq:dual-gain-condition}, and
\eqref{eq:multiplier-boundedness-condition}.

\begin{lemma}
\label{lem:moving-three-value-descent}
\blue{Suppose that Assumptions~\ref{ass:level-boundedness},
\ref{ass:smoothness}, and~\ref{ass:h} hold, and let}
$\{(\bx^k,\blmbd^k,\bz^k,\bmu^k)\}$ be generated by \Cref{alg:ddlalm}
with parameters satisfying~\eqref{eq:parameter-choice}. Then for every
$k\ge0$:
\begin{enumerate}[label=(\alph*)]
    \item for all $\bx\in\RR^n$, $\blmbd\in\RR^m$, and $\bz\in\RR^n$,
    \begin{equation}
    \label{eq:current-lyapunov-lower-bound}
        \Psi_k(\bx,\blmbd;\bz)\ge p_k(\bz)
        \ge\inf_{\bx\in\RR^n}F(\bx)-\frac{(L_h+M)^2}{2s_0};
    \end{equation}
    \item the iterates satisfy
    \begin{equation}
    \label{eq:full-lyapunov-descent}
    \begin{aligned}
        &\Psi_{k+1}(\bx^{k+1},\blmbd^{k+1};\bz^{k+1})
         -\Psi_k(\bx^k,\blmbd^k;\bz^k)\\
        &\qquad\le
         -\frac1{4\eta_k}\|\Delta_{\bx}^{k+1}\|^2
         -\frac1{8\tau_k}\|\Delta_{\blmbd}^{k+1}\|^2
         -\frac{\alpha_k}{2\gamma}\|\bx^{k+1}-\bz^k\|^2
         +\delta_k,
    \end{aligned}
    \end{equation}
    where
    $\delta_k:=10M^2\bigl(s_k^{-1}-s_{k+1}^{-1}\bigr)+20M^2\beta_k/s_k$.
\end{enumerate}
\end{lemma}


% For $s_k=c_s(k+1)^{1/3}$ and
% $\beta_k=\theta_\mu/(k+k_0)$ with $k_0\ge1$, the total error satisfies
% \begin{equation}
% \label{eq:explicit-transport-budget}
% \sum_{k=0}^\infty
% \left[
% 10M^2\left|\frac1{s_{k+1}}-\frac1{s_k}\right|
% +20M^2\frac{\beta_k}{s_k}
% \right]
% \le\frac{(10+80\theta_\mu)M^2}{c_s}.
% \end{equation}

% The augmented estimate adds $2p_k$ to the partial potential in
% \eqref{eq:partial-potential}. It uses the $1/\gamma$-smoothness of $p_k$
% once. The remaining
% inner product is the only new $\bz$-direction error that must be controlled
% by the subsequent stability estimates.

% \begin{lemma}[Saddle-point stability]
% \label{lem:saddle-point-stability}
% Under the conditions of
% Proposition~\ref{prop:curvature-smoothness}, for each $k\ge 0$ and
% each fixed $\bz$, let
% \[
%     \blmbd_k^\star(\bz)
%     :=\argmax_{\blmbd\in\RR^m}d_k(\blmbd;\bz).
% \]
% Then
% \[
%     \bx_k^\star(\bz)
%     =
%     \bx_k\bigl(\blmbd_k^\star(\bz),\bz\bigr),
% \]
% and
% \begin{equation}
% \label{eq:saddle-primal-stability}
% \left\|
%     \bx_k(\blmbd;\bz)-\bx_k^\star(\bz)
% \right\|
% \le
% \frac{G_{\mathcal A}}{\mu_\gamma}
% \left\|
%     \blmbd-\blmbd_k^\star(\bz)
% \right\|.
% \end{equation}
% Moreover,
% \begin{equation}
% \label{eq:dual-gap-to-multiplier-distance}
% \left\|
%     \blmbd-\blmbd_k^\star(\bz)
% \right\|^2
% \le
% 2s_k
% \bigl[
%     p_k(\bz)-d_k(\blmbd;\bz)
% \bigr].
% \end{equation}
% Consequently,
% \begin{equation}
% \label{eq:saddle-point-stability-final}
% \left\|
%     \bx_k(\blmbd;\bz)-\bx_k^\star(\bz)
% \right\|^2
% \le
% \frac{2G_{\mathcal A}^2s_k}{\mu_\gamma^2}
% \bigl[
%     p_k(\bz)-d_k(\blmbd;\bz)
% \bigr].
% \end{equation}
% \end{lemma}

The proof is given in Appendix~\ref{app:curvature-smoothness}.


% !TEX root = ../main.tex
\section{Proofs for Appendix~\ref{sec:upper-bound}}
\label{app:curvature-smoothness}

\begin{proof}[Proof of the curvature part of Proposition~\ref{prop:curvature-smoothness}]
For fixed $\blmbd$, the reduced payoff satisfies
\[
    \nabla_{\bx}\mathcal D_\rho(\bx,\blmbd)
    =
    \nabla f(\bx)+J_{\calA}(\bx)^\top\bq(\bx,\blmbd),
\]
where
\[
    \bq(\bx,\blmbd)
    =
    \operatorname{prox}_{\rho h^*}
    \bigl(\blmbd+\rho\calA(\bx)\bigr).
\]

The nonexpansiveness of the proximal mapping, the bounded Jacobian, and
the Lipschitz Jacobian imply
\[
    \|\bq(\bx,\blmbd)-\bq(\bx',\blmbd)\|
    \le \rho G_{\calA}\|\bx-\bx'\|.
\]
Assumption~\ref{ass:h} gives
\(\|\bq(\bx,\blmbd)\|\le L_h\).  Therefore
\[
    \|\nabla_{\bx}\mathcal D_\rho(\bx,\blmbd)
    -\nabla_{\bx}\mathcal D_\rho(\bx',\blmbd)\|
    \le L_{\calD,\bx}\|\bx-\bx'\|.
\]

Thus $\mathcal D_\rho(\cdot,\blmbd)$ is
$L_{\calD,\bx}$-weakly convex.  Adding
$\frac{1}{2\gamma}\|\bx-\bz\|^2$ gives
$\mu_\gamma$-strong convexity in $\bx$.

For fixed $(\bx,\bz,\bmu)$, $\mathcal D_\rho(\bx,\cdot)$ is concave in
$\blmbd$ because it is the pointwise infimum of affine functions of
$\blmbd$.  The term
$-\frac{1}{2s}\|\blmbd-\bmu\|^2$ is $1/s$-strongly concave.  Hence
$\mathcal P_s$ is $1/s$-strongly concave in $\blmbd$.
\end{proof}

\begin{proof}[Proof of the smoothness part of Proposition~\ref{prop:curvature-smoothness}]
The $\bx$-block bound follows from the Lipschitz estimate for
$\nabla_{\bx}\mathcal D_\rho$ in the proof of
Proposition~\ref{prop:curvature-smoothness} and the quadratic term
$\frac{1}{2\gamma}\|\bx-\bz\|^2$.

For the $\blmbd$-block, the identity
\[
    \nabla_{\blmbd}\mathcal P_s(\bx,\blmbd;\bz,\bmu)
    =
    \frac{1}{\rho}
    \bigl(\bq(\bx,\blmbd)-\blmbd\bigr)
    -\frac{1}{s}(\blmbd-\bmu)
\]
and the nonexpansiveness of
$\operatorname{prox}_{\rho h^*}$ give
$L_\lambda(s)=1/\rho+1/s$.
The same estimate, together with
\[
    \|\calA(\bx)-\calA(\bx')\|
    \le G_{\calA}\|\bx-\bx'\|,
\]
gives both cross Lipschitz bounds with constant $G_{\calA}$.
\end{proof}

\begin{proof}[Proof of Lemma~\ref{lem:dual-value-gain}]
Set
\[
    g^k
    :=
    \nabla_{\blmbd}\mathcal P_k
    (\bx^{k+1},\blmbd^k,\bz^k),
    \qquad
    \Delta_{\blmbd}^{k+1}=\tau_k g^k,
\]
where the second identity follows from the dual update.
By the $L_{d,k}$-smoothness of $d_k(\cdot;\bz^k)$,
\[
\begin{aligned}
&d_k(\blmbd^{k+1};\bz^k)-d_k(\blmbd^k;\bz^k) \\
&\quad\ge
\left\langle
\nabla_{\blmbd}d_k(\blmbd^k;\bz^k),
\Delta_{\blmbd}^{k+1}
\right\rangle
-\frac{L_{d,k}}{2}
\|\Delta_{\blmbd}^{k+1}\|^2.
\end{aligned}
\]
Danskin's theorem and Lemma~\ref{lem:v2-value-regularity} give
\[
\nabla_{\blmbd}d_k(\blmbd^k;\bz^k)
=
\nabla_{\blmbd}\mathcal P_k
\bigl(\bx_k(\blmbd^k;\bz^k),\blmbd^k,\bz^k\bigr).
\]
The cross-Lipschitz estimate in
Proposition~\ref{prop:curvature-smoothness} therefore yields
\[
\begin{aligned}
&\left\langle
\nabla_{\blmbd}d_k(\blmbd^k;\bz^k),
\Delta_{\blmbd}^{k+1}
\right\rangle \\
&\quad=
\tau_k\left\langle
\nabla_{\blmbd}\mathcal P_k
\bigl(\bx_k(\blmbd^k;\bz^k),\blmbd^k,\bz^k\bigr),g^k
\right\rangle \\
&\quad\ge
\tau_k\|g^k\|^2
-\tau_kG_{\cal A}
\|\bx^{k+1}-\bx_k(\blmbd^k;\bz^k)\|\|g^k\|.
\end{aligned}
\]
Using
\[
    \tau_k\|g^k\|^2
    =\frac{1}{\tau_k}\|\Delta_{\blmbd}^{k+1}\|^2,
    \qquad
    \tau_k\|g^k\|
    =\|\Delta_{\blmbd}^{k+1}\|,
\]
we obtain
\begin{equation}
\begin{aligned}
&d_k(\blmbd^{k+1};\bz^k)-d_k(\blmbd^k;\bz^k) \\
&\quad\ge
\left(\frac{1}{\tau_k}-\frac{L_{d,k}}{2}\right)
\|\Delta_{\blmbd}^{k+1}\|^2
-G_{\cal A}\|\bx^{k+1}-\bx_k(\blmbd^k;\bz^k)\|
\|\Delta_{\blmbd}^{k+1}\|.
\end{aligned}
\label{eq:dual-value-gain-appendix}
\end{equation}

Let
\[
    b_k:=
    \frac{1}{\tau_k}+\frac{1}{2s_k}-L_{d,k}.
\]
Multiplying \eqref{eq:dual-value-gain-appendix} by $2$ and rearranging gives
\[
\begin{aligned}
&\left(\frac{1}{\tau_k}-\frac{1}{2s_k}\right)
\|\Delta_{\blmbd}^{k+1}\|^2 \\
&\quad\le
2\bigl[d_k(\blmbd^{k+1};\bz^k)-d_k(\blmbd^k;\bz^k)\bigr]
-b_k\|\Delta_{\blmbd}^{k+1}\|^2 \\
&\qquad+
2G_{\cal A}\|\bx^{k+1}-\bx_k(\blmbd^k;\bz^k)\|
\|\Delta_{\blmbd}^{k+1}\|.
\end{aligned}
\]
The assumed parameter condition implies $b_k\ge L_{d,k}>0$. Young's inequality
with parameter $L_{d,k}$ gives
\[
2G_{\cal A}\|\bx^{k+1}-\bx_k(\blmbd^k;\bz^k)\|
\|\Delta_{\blmbd}^{k+1}\|
\le
\frac{L_{d,k}}{2}\|\Delta_{\blmbd}^{k+1}\|^2
+\frac{2G_{\cal A}^2}{L_{d,k}}
\|\bx^{k+1}-\bx_k(\blmbd^k;\bz^k)\|^2.
\]
Since $b_k\ge L_{d,k}$, we have
\[
    -b_k\|\Delta_{\blmbd}^{k+1}\|^2
    +\frac{L_{d,k}}{2}\|\Delta_{\blmbd}^{k+1}\|^2
    \le
    -\frac{L_{d,k}}{2}\|\Delta_{\blmbd}^{k+1}\|^2.
\]
Finally, the $\mu_\gamma$-strong convexity of
$\mathcal P_k(\cdot,\blmbd^k;\bz^k)$ gives
\[
\|\bx^{k+1}-\bx_k(\blmbd^k;\bz^k)\|^2
\le
\frac{2}{\mu_\gamma}
\left[
\mathcal P_k(\bx^{k+1},\blmbd^k;\bz^k)
-d_k(\blmbd^k;\bz^k)
\right].
\]
Substituting this estimate proves the first inequality.
\end{proof}

\begin{proof}[Proof of the second estimate in Lemma~\ref{lem:dual-value-gain}]
Fix $k\ge0$. By Proposition~\ref{prop:curvature-smoothness},
$\mathcal P_k(\cdot,\blmbd^k;\bz^k)$ is
$\mu_\gamma$-strongly convex. Applying the strong-convexity inequality
at $\bx^{k+1}$ to the minimizer $\bx_k(\blmbd^k;\bz^k)$ gives
\[
\begin{aligned}
&\mathcal P_k(\bx^{k+1},\blmbd^k;\bz^k)-d_k(\blmbd^k;\bz^k) \\
&\quad\le
\left\langle
\nabla_{\bx}\mathcal P_k(\bx^{k+1},\blmbd^k;\bz^k),
\bx^{k+1}-\bx_k(\blmbd^k;\bz^k)
\right\rangle \\
&\qquad
-\frac{\mu_\gamma}{2}
\|\bx^{k+1}-\bx_k(\blmbd^k;\bz^k)\|^2 \\
&\quad\le
\frac{1}{2\mu_\gamma}
\|\nabla_{\bx}\mathcal P_k(\bx^{k+1},\blmbd^k;\bz^k)\|^2,
\end{aligned}
\]
where the last inequality follows from Young's inequality.
The $\bx$-update and the $L_x$-Lipschitz continuity of
$\nabla_{\bx}\mathcal P_k(\cdot,\blmbd^k;\bz^k)$ imply
\[
\begin{aligned}
&\|\nabla_{\bx}\mathcal P_k(\bx^{k+1},\blmbd^k;\bz^k)\| \\
&\quad\le
\|\nabla_{\bx}\mathcal P_k(\bx^k,\blmbd^k;\bz^k)\|
+L_x\|\Delta_{\bx}^{k+1}\| \\
&\quad=
\left(\frac{1}{\eta_k}+L_x\right)
\|\Delta_{\bx}^{k+1}\|.
\end{aligned}
\]
Combining these estimates yields
\begin{equation}
\label{eq:primal-gap-control}
\mathcal P_k(\bx^{k+1},\blmbd^k;\bz^k)
-d_k(\blmbd^k;\bz^k)
\le
\frac{(1/\eta_k+L_x)^2}{2\mu_\gamma}
\|\Delta_{\bx}^{k+1}\|^2.
\end{equation}
Substituting \eqref{eq:primal-gap-control} into the first inequality of
Lemma~\ref{lem:dual-value-gain} gives a bound on the dual ascent term,
where the coefficient of
$\|\Delta_{\bx}^{k+1}\|^2$ becomes
\[
\frac{4G_{\mathcal A}^2}{L_{d,k}\mu_\gamma}
\frac{(1/\eta_k+L_x)^2}{2\mu_\gamma}
=
\frac{2G_{\mathcal A}^2(1/\eta_k+L_x)^2}
{L_{d,k}\mu_\gamma^2}.
\]
Substitute the resulting bound on
$(1/\tau_k-1/(2s_k))\|\Delta_{\blmbd}^{k+1}\|^2$
into \eqref{eq:basic-one-step} and combine the coefficients of
$\|\Delta_{\bx}^{k+1}\|^2$. The $\bz$-step term and both parameter-drift
terms remain unchanged. This proves \eqref{eq:dual-value-full-step}.
\end{proof}

\begin{proof}[Proof of Lemma~\ref{lem:partial-lyapunov-z-step}]
Subtract
$2[d_k(\blmbd^{k+1};\bz^k)-d_k(\blmbd^k;\bz^k)]$
from both sides of \eqref{eq:dual-value-full-step}. To replace
$d_k(\blmbd^{k+1};\bz^k)$ on the left by
$d_k(\blmbd^{k+1};\bz^{k+1})$, add
$-2[d_k(\blmbd^{k+1};\bz^{k+1})-d_k(\blmbd^{k+1};\bz^k)]$
to both sides.
By the minimum-value definition,
\[
d_k(\blmbd^{k+1};\bz^k)
\le
\mathcal P_k\bigl(\bx_k(\blmbd^{k+1};\bz^{k+1}),
\blmbd^{k+1};\bz^k\bigr).
\]
Subtracting this inequality from the corresponding equality at
$\bz^{k+1}$ and expanding the quadratic terms yields
\[
\begin{aligned}
&d_k(\blmbd^{k+1};\bz^{k+1})-d_k(\blmbd^{k+1};\bz^k) \\
&\quad\ge
\frac1\gamma\left\langle
\bz^k-\bx_k(\blmbd^{k+1};\bz^{k+1}),
\Delta_{\bz}^{k+1}\right\rangle
+\frac1{2\gamma}\|\Delta_{\bz}^{k+1}\|^2 \\
&\quad=
-\frac{\alpha_k(2-\alpha_k)}{2\gamma}
\|\bx^{k+1}-\bz^k\|^2 \\
&\qquad+
\frac{\alpha_k}{\gamma}\left\langle
\bx^{k+1}-\bx_k(\blmbd^{k+1};\bz^{k+1}),
\bx^{k+1}-\bz^k\right\rangle.
\end{aligned}
\]
Multiply this inequality by $-2$ and substitute it into the preceding
estimate. Its squared-step term combines with the $\bz$-step term in
\eqref{eq:dual-value-full-step} to give
$\alpha_k(2-\alpha_k)\|\bx^{k+1}-\bz^k\|^2/(2\gamma)$.
Both parameter-drift terms remain unchanged, proving
\eqref{eq:partial-lyapunov-z-step}.
\end{proof}

\begin{proof}[Proof of Lemma~\ref{lem:z-step-p}]
By Lemma~\ref{lem:v2-value-regularity}, $p_k$ has a
$1/\gamma$-Lipschitz continuous gradient and
\[
    \nabla p_k(\bz^k)
    =
    \frac1\gamma\bigl(\bz^k-\bx_k^\star(\bz^k)\bigr).
\]
The $\bz$-update gives
\[
    \bz^{k+1}-\bz^k
    =\alpha_k(\bx^{k+1}-\bz^k).
\]
Applying the smoothness inequality for $p_k$ and expanding the inner
product yields \eqref{eq:z-step-p}.
\end{proof}

\begin{proof}[Proof of Lemma~\ref{lem:augmented-one-step}]
By Lemma~\ref{lem:z-step-p},
\begin{equation}
\begin{aligned}
2\bigl[p_k(\bz^{k+1})-p_k(\bz^k)\bigr]
\le{}&
-\frac{\alpha_k(2-\alpha_k)}{\gamma}
\|\bx^{k+1}-\bz^k\|^2 \\
&+
\frac{2\alpha_k}{\gamma}
\left\langle
\bx^{k+1}-\bx_k^\star(\bz^k),
\bx^{k+1}-\bz^k
\right\rangle.
\end{aligned}
\end{equation}
Adding this inequality to
\eqref{eq:partial-lyapunov-z-step}, the quadratic terms combine as
\[
\frac{\alpha_k(2-\alpha_k)}{2\gamma}
-\frac{\alpha_k(2-\alpha_k)}{\gamma}
=
-\frac{\alpha_k(2-\alpha_k)}{2\gamma}.
\]
The two inner products combine into
\[
\frac{2\alpha_k}{\gamma}
\left\langle
\bx_k(\blmbd^{k+1};\bz^{k+1})
-\bx_k^\star(\bz^k),
\bx^{k+1}-\bz^k
\right\rangle.
\]
Substitution into \eqref{eq:partial-lyapunov-z-step} proves
\eqref{eq:augmented-one-step}.
\end{proof}

\begin{proof}[Proof of Lemma~\ref{lem:current-multiplier-boundedness}]
The proximal optimality condition gives
$\bq(\bx,\blmbd)\in\partial h(\by(\bx,\blmbd))$, and hence
$\|\bq(\bx,\blmbd)\|\le L_h$. Substituting
$\calA(\bx)-\by(\bx,\blmbd)
=(\bq(\bx,\blmbd)-\blmbd)/\rho$ into the dual update yields
\[
\blmbd^{k+1}
=\left[1-\tau_k\left(\frac1\rho+\frac1{s_k}\right)\right]\blmbd^k
+\frac{\tau_k}{\rho}\bq(\bx^{k+1},\blmbd^k)
+\frac{\tau_k}{s_k}\bmu^k.
\]
The coefficients are nonnegative and sum to one. Together with
$\bmu^{k+1}=(1-\beta_k)\bmu^k+\beta_k\blmbd^{k+1}$, this proves
$\|\blmbd^k\|,\|\bmu^k\|\le M$ by induction from the definition of $M$.
The center update then gives
$\|\bmu^{k+1}-\bmu^k\|\le2M\beta_k$.

At the saddle point, the dual first-order condition gives
\[
\blmbd_k^\star(\bz)
=\frac{s_k}{\rho+s_k}
  \bq\bigl(\bx_k^\star(\bz),\blmbd_k^\star(\bz)\bigr)
+\frac{\rho}{\rho+s_k}\bmu^k.
\]
Since both vectors in this convex combination have norm at most $M$,
$\|\blmbd_k^\star(\bz)\|\le M$.
\end{proof}

\begin{proof}[Proof of Lemma~\ref{lem:saddle-center-stability}]
Fix $k$ and $\bz$. Strong convexity gives the unique minimizer
$\bx_k(\blmbd;\bz)$ for each $\blmbd$. The finite, smooth,
$1/s_k$-strongly concave function $d_k(\cdot;\bz)$ tends to $-\infty$
as $\|\blmbd\|\to\infty$, so it has a unique maximizer.
The minimax equality and strong convexity in $\bx$ identify the unique
saddle pair and give
$\bx_k^\star(\bz)=\bx_k(\blmbd_k^\star(\bz);\bz)$.

The maximizing property at $\bx_k^\star(\bz)$ and strong convexity at
fixed $\blmbd$ imply
\[
\begin{aligned}
p_k(\bz)-d_k(\blmbd;\bz)
&\ge
\mathcal P_k(\bx_k^\star(\bz),\blmbd;\bz)
-\mathcal P_k(\bx_k(\blmbd;\bz),\blmbd;\bz) \\
&\ge\frac{\mu_\gamma}{2}
\|\bx_k(\blmbd;\bz)-\bx_k^\star(\bz)\|^2.
\end{aligned}
\]
This proves \eqref{eq:current-saddle-stability}.
For the center sensitivity of $\bx_k(\blmbd;\bz)$, add the
strong-convexity inequalities at the minimizers for $\bz$ and $\bz'$.
The only center-dependent term is $\|\bx-\bz\|^2/(2\gamma)$, so
\[
\mu_\gamma\|\bx_k(\blmbd;\bz')-\bx_k(\blmbd;\bz)\|^2
\le\frac1\gamma\left\langle
\bz'-\bz,\bx_k(\blmbd;\bz')-\bx_k(\blmbd;\bz)\right\rangle.
\]
For the saddle primal points, use the same argument for
$\max_{\blmbd}\mathcal P_k(\cdot,\blmbd;\bz)$, whose strong-convexity
modulus is also $\mu_\gamma$. Cauchy--Schwarz gives
\eqref{eq:current-center-stability} in both cases.
The latter argument also gives
\[
\left\langle\bz'-\bz,\bx_k^\star(\bz')-\bx_k^\star(\bz)\right\rangle
\ge\gamma\mu_\gamma
\|\bx_k^\star(\bz')-\bx_k^\star(\bz)\|^2.
\]
Together with the Danskin formula for $\nabla p_k$ and
$2\gamma\mu_\gamma>1$, this verifies the smoothness bound used above:
\[
\begin{aligned}
\gamma^2\|\nabla p_k(\bz')-\nabla p_k(\bz)\|^2
&\le\|\bz'-\bz\|^2
-(2\gamma\mu_\gamma-1)
\|\bx_k^\star(\bz')-\bx_k^\star(\bz)\|^2 \\
&\le\|\bz'-\bz\|^2.
\end{aligned}
\]
\end{proof}

\begin{proof}[Proof of Lemma~\ref{lem:moving-three-value-descent}]
Let $u_{k+1}:=(\bx^{k+1},\blmbd^{k+1},\bz^{k+1})$ and define
\[
\widehat\Psi_{k+1|k}:=
\mathcal P_{k+1}(\bx^{k+1},\blmbd^{k+1};\bz^{k+1})
-2d_k(\blmbd^{k+1};\bz^{k+1})+2p_k(\bz^{k+1}).
\]
Then
\[
\Psi_{k+1}(u_{k+1})-\Psi_k(u_k)
=\widehat\Psi_{k+1|k}-\Psi_k(u_k)+R_k,
\]
where
\[
R_k:=2\bigl[d_k(\blmbd^{k+1};\bz^{k+1})
-d_{k+1}(\blmbd^{k+1};\bz^{k+1})\bigr]
+2\bigl[p_{k+1}(\bz^{k+1})-p_k(\bz^{k+1})\bigr].
\]

Apply Lemma~\ref{lem:augmented-one-step} to the first difference.
Decompose the indefinite term at $\bx_k^\star(\bz^{k+1})$ and apply
Lemma~\ref{lem:saddle-center-stability}, followed by Young's inequality.
The four bounds in \eqref{eq:parameter-choice} absorb the resulting
cross terms and leave
\[
\begin{aligned}
\widehat\Psi_{k+1|k}-\Psi_k(u_k)
\le{}&
-\frac1{4\eta_k}\|\Delta_{\bx}^{k+1}\|^2
-\frac1{8\tau_k}\|\Delta_{\blmbd}^{k+1}\|^2
-\frac{\alpha_k}{2\gamma}\|\bx^{k+1}-\bz^k\|^2 \\
&+\frac{\beta_k(2-\beta_k)}{2s_k}
\|\blmbd^{k+1}-\bmu^k\|^2
\\
&+\frac12\left(\frac1{s_k}-\frac1{s_{k+1}}\right)
\|\blmbd^{k+1}-\bmu^{k+1}\|^2.
\end{aligned}
\]

It remains to bound $R_k$. Lemma~\ref{lem:current-multiplier-boundedness}
gives $\|\blmbd^j\|,\|\bmu^j\|,\|\blmbd_j^\star(\bz)\|\le M$ and
$\|\bmu^{k+1}-\bmu^k\|\le2M\beta_k$. Direct comparison of the two
quadratic damping terms gives, uniformly in $\bx$ and $\bz$,
\[
|d_{k+1}(\blmbd;\bz)-d_k(\blmbd;\bz)|
\le\Delta_k,
\qquad
|p_{k+1}(\bz)-p_k(\bz)|\le\Delta_k,
\]
where
\[
\Delta_k:=2M^2\left(\frac1{s_k}-\frac1{s_{k+1}}\right)
+\frac{2M}{s_k}\|\bmu^{k+1}-\bmu^k\|.
\]
Consequently,
\[
|R_k|\le4\Delta_k
\le8M^2\left(\frac1{s_k}-\frac1{s_{k+1}}\right)
+16M^2\frac{\beta_k}{s_k}.
\]
The two explicit positive terms above satisfy
\[
\frac{\beta_k(2-\beta_k)}{2s_k}
\|\blmbd^{k+1}-\bmu^k\|^2
\le4M^2\frac{\beta_k}{s_k},
\]
and
\[
\frac12\left(\frac1{s_k}-\frac1{s_{k+1}}\right)
\|\blmbd^{k+1}-\bmu^{k+1}\|^2
\le2M^2\left(\frac1{s_k}-\frac1{s_{k+1}}\right).
\]
Adding the estimates proves \eqref{eq:full-lyapunov-descent}.

For the stated schedule, the two error series satisfy
\[
\sum_{k=0}^\infty\left(\frac1{s_k}-\frac1{s_{k+1}}\right)=\frac1{c_s},
\qquad
\sum_{k=0}^\infty\frac{\beta_k}{s_k}
\le\frac{4\theta_\mu}{c_s},
\]
which makes the accumulated transport error finite. The lower bound
\eqref{eq:current-lyapunov-lower-bound} follows as follows. For any
$\blmbd\in\partial h(\calA(\bx))$, proximal optimality gives
$\by(\bx,\blmbd)=\calA(\bx)$ and hence
$\mathcal D_\rho(\bx,\blmbd)=F(\bx)$. Since
$\|\blmbd\|\le L_h$ and $\|\bmu^k\|\le M$,
\[
\max_{\blmbd}\mathcal P_k(\bx,\blmbd;\bz)
\ge F(\bx)+\frac1{2\gamma}\|\bx-\bz\|^2
-\frac{(L_h+M)^2}{2s_k}.
\]
Minimizing in $\bx$ and using $s_k\ge s_0$ proves
\eqref{eq:current-lyapunov-lower-bound}, because
$\Psi_k=(\mathcal P_k-d_k)+(p_k-d_k)+p_k\ge p_k$.
Summing \eqref{eq:full-lyapunov-descent} and using the error budget
then gives
\[
\begin{aligned}
&\sum_{k=0}^{N-1}\left[
\frac1{4\eta_k}\|\Delta_{\bx}^{k+1}\|^2
+\frac1{8\tau_k}\|\Delta_{\blmbd}^{k+1}\|^2
+\frac{\alpha_k}{2\gamma}\|\bx^{k+1}-\bz^k\|^2\right] \\
&\quad\le
\Psi_0(\bx^0,\blmbd^0;\bz^0)-\inf_{\bx\in\RR^n}F(\bx)
+\frac{(L_h+M)^2}{2s_0}
+\frac{(10+80\theta_\mu)M^2}{c_s},
\qquad N\ge1.
\end{aligned}
\]
\end{proof}



\begin{proof}[Proof of Theorem~\ref{thm:main}]
The schedule \eqref{eq:ddlalm-parameter-schedule} satisfies
\eqref{eq:parameter-choice}: $\eta_kL_x=1/2$,
$\tau_kL_{d,k}\le\tau_kL_{d,0}\le1/2$,
$\tau_kG_{\cal A}^2\le\mu_\gamma^2/(72L_x)
=\eta_k\mu_\gamma^2/36$, and the two bounds on $\alpha_k$ follow
from $s_k\ge s_0$ and its definition. Moreover,
$s_{k+1}\ge s_k>0$, $\eta_k,\tau_k>0$, and
$0<\beta_k\le1$. Since $L_{d,k}\ge1/\rho+1/s_k$,
the multiplier-boundedness condition also holds.
By assumption, $\inf_{\bx}F(\bx)>-\infty$. The master budget proved after
Lemma~\ref{lem:moving-three-value-descent} is therefore finite and
independent of $K$.

The tail has $\lfloor K/2\rfloor\ge K/3$ elements. Since the primal
and dual stepsizes are fixed and
$\alpha_k\ge\alpha_{K-1}=\Theta(K^{-1/3})$ on the tail, the master
budget and the definition of $\kappa$ give
\[
\|\Delta_{\bx}^{\kappa+1}\|^2=O(K^{-1}),\qquad
\|\Delta_{\blmbd}^{\kappa+1}\|^2=O(K^{-1}),\qquad
\|\bx^{\kappa+1}-\bz^\kappa\|^2=O(K^{-2/3}).
\]
The output certificate follows directly from
\eqref{eq:prox-multiplier}. For stationarity, the $\bx$-update
and $L_{\calD,\bx}$-smoothness of $\mathcal D_\rho(\cdot,\blmbd^k)$ give
\[
\|\nabla f(\bx^{k+1})+J_{\cal A}(\bx^{k+1})^\top\bq^k\|
\le\left(\frac1{\eta_k}+L_x\right)\|\Delta_{\bx}^{k+1}\|
+\frac1\gamma\|\bx^{k+1}-\bz^k\|.
\]
For feasibility, the dual update gives
\[
{\cal A}(\bx^{k+1})-\by^{k+1}
=\frac1{\tau_k}\Delta_{\blmbd}^{k+1}
+\frac1{s_k}(\blmbd^k-\bmu^k),
\]
and hence
\[
\|{\cal A}(\bx^{k+1})-\by^{k+1}\|
\le\frac1{\tau_k}\|\Delta_{\blmbd}^{k+1}\|+\frac{2M}{s_k}.
\]
Square the last two bounds using $(a+b)^2\le2a^2+2b^2$ and use
$s_\kappa^{-2}=O(K^{-2/3})$. Both squared residuals are
$O(K^{-2/3})$, so their norms are $O(K^{-1/3})$. All constants depend
only on the problem data and fixed initialization and parameters, so
$K=O(\epsilon^{-3})$ suffices.
\end{proof}


\end{document}
